% La science — Philosophie Tle Tech — Cours signé OnBuch+
% Programme officiel MINESEC. Compiler avec : tectonic N10-science.tex
\newcommand{\DOCMATIERE}{Philosophie}
\newcommand{\DOCNIVEAU}{Tle Tech}
\newcommand{\DOCMODULE}{Philosophie – classe de Terminale technique}
\newcommand{\DOCLECON}{N10}
\newcommand{\DOCTITRE}{La science}
% OnBuch+ — préambule « cours signés » ENRICHI (compilé avec tectonic / XeLaTeX)
% Système visuel « Pop » d'OnBuch+ : fond crème (jamais blanc), bordures encre
% épaisses, ombres « dures » décalées, titres Archivo Black en capitales.
% Version MATHÉMATIQUES : graphes (pgfplots/tikz), boîtes pédagogiques
% (définition, propriété, méthode, attention, le savais-tu, exercice résolu,
% exercice, TP, bilan), page de garde et signature OnBuch+ ancrée en bas.
%
% Macros attendues dans main.tex AVANT \input{preamble} :
%   \DOCMATIERE  \DOCNIVEAU  \DOCMODULE  \DOCLECON (numéro)  \DOCTITRE
\documentclass[11pt]{article}

\usepackage{fontspec}
\usepackage[french]{babel}
\usepackage[a4paper,margin=2cm,top=3.4cm,bottom=2.8cm,headheight=1.4cm,footskip=1.3cm]{geometry}
\usepackage{amsmath,amssymb,amsfonts}
\usepackage{mathtools} % \xmapsto, \coloneqq... souvent utilisés par les modèles
\usepackage{cancel} % simplifications barrées (bilans de forces)
% \abs{z} (module, valeur absolue) et \norm{u} (norme), utilisés par les modèles
\providecommand{\abs}{}\let\abs\relax\DeclarePairedDelimiter{\abs}{\lvert}{\rvert}
\providecommand{\norm}{}\let\norm\relax\DeclarePairedDelimiter{\norm}{\lVert}{\rVert}
\usepackage[table]{xcolor} % [table] : fournit aussi \rowcolors (alternance de lignes)
\usepackage{pagecolor}
\usepackage{tikz}
\usetikzlibrary{arrows.meta,positioning,calc,shapes.geometric,shapes.misc,shapes.arrows,decorations.pathmorphing,decorations.pathreplacing,decorations.markings,patterns,shadows,mindmap,trees,fit,backgrounds,matrix,intersections,angles,quotes}
\usepackage{pgfplots}
\usepackage[version=4]{mhchem} % équations nucléaires \ce{^{235}_{92}U}
\usepackage[european,siunitx]{circuitikz} % schémas électriques
\pgfplotsset{compat=1.18}
\usepgfplotslibrary{fillbetween,groupplots} % aires entre courbes (calcul intégral)
\usepackage{enumitem}
\usepackage{fancyhdr}
% [most] charge toutes les bibliothèques tcolorbox (listings, minted, raster,
% documentation...) et gonfle inutilement la mémoire TeX sur un document long
% (~300 boîtes) : on ne charge que ce qui est réellement utilisé.
\usepackage[skins,breakable]{tcolorbox}
\usepackage{graphicx}
\usepackage{colortbl}
\usepackage{booktabs}
\usepackage{tabularx}
\usepackage{multirow}
\usepackage{diagbox} % case d'en-tête diagonale des tableaux à double entrée
% Colonnes extensibles centrée / gauche / droite (C, L, R), souvent utilisées par les modèles
\newcolumntype{C}{>{\centering\arraybackslash}X}
\newcolumntype{L}{>{\raggedright\arraybackslash}X}
\newcolumntype{R}{>{\raggedleft\arraybackslash}X}
\usepackage{array}
\usepackage{multicol}
\usepackage{siunitx}
% parse-numbers=false : accepte aussi \SI{2,5 \times 10^{-3}}{...}
\sisetup{locale=FR,output-decimal-marker={,},group-minimum-digits=4,parse-numbers=false}
\DeclareSIUnit\ohm{\ensuremath{\symup{\Omega}}} % U+2126 absent de la police
\DeclareSIUnit\division{div} % réglages d'oscilloscope (ms/div, V/div)
\DeclareSIUnit\tr{tr} % tours (vitesse de rotation, ex. \SI{1500}{\tr\per\minute})
\DeclareSIUnit\molar{M} % concentration molaire (ex. \SI{3}{\molar}, \SI{1}{\micro\molar})
\DeclareSIUnit\an{an} % année (le modèle confond parfois avec \year, un compteur TeX) — ex. \SI{5}{\kilo\meter\per\an}
\DeclareSIUnit\FCFA{FCFA} % franc CFA (ex. \SI{500}{\FCFA\per\kilo\gram})
\DeclareSIUnit\degreeBrix{\textdegree Brix} % teneur en sucre d'un sirop/jus (réfractométrie)
\DeclareSIUnit\month{mois} % le modèle utilise parfois \month, un compteur TeX, comme unité de durée
\DeclareSIUnit\microliter{\micro\liter} % le modèle écrit parfois \microliter en un seul token
\DeclareSIUnit\million{millions} % dénombrement (ex. \SI{15}{\million\per\mL} spermatozoïdes)
\usepackage{qrcode}
% \degree : commande fréquemment utilisée par les modèles pour un angle ou une
% température en dehors de \SI{}{} (non fournie par les paquets chargés).
\providecommand{\degree}{^{\circ}}
% \qed (fin de démonstration), utilisé par les modèles sans amsthm
\providecommand{\qed}{\hfill$\square$}
% \barre : double barre des valeurs interdites dans un tableau de variations (tkz-tab)
\providecommand{\barre}{\ensuremath{\|}}
% environnement proof (amsthm non chargé) utilisé par les modèles
\ifdefined\proof\else\newenvironment{proof}[1][Démonstration]{\par\noindent\textit{#1.}\ }{\qed\par}\fi
\usepackage{lastpage}
\usepackage{hyperref}
\hypersetup{hidelinks}

% ============================================================
% Palette « Pop » OnBuch+ — mêmes hex que la classe P (plus_ui.dart)
% ============================================================
\definecolor{poporange}{HTML}{F59321}
\definecolor{popdark}{HTML}{E07A0C}
\definecolor{popcream}{HTML}{FFF6E8}
\definecolor{popink}{HTML}{1C1714}
\definecolor{poppurple}{HTML}{7A5AE0}
\definecolor{popgreen}{HTML}{2E9E6A}
\definecolor{poppink}{HTML}{E0457B}
\definecolor{popblue}{HTML}{2F7DE1}
\definecolor{popgold}{HTML}{C98A12}
\definecolor{popmuted}{HTML}{8A7F76}
\definecolor{popline}{HTML}{EADFD2}
\definecolor{popgray}{HTML}{8A7F76}
\colorlet{popgrey}{popgray}
% Alias tolérants (le modèle nomme parfois les couleurs différemment)
\colorlet{popwhite}{white}
\colorlet{popblack}{popink}
\colorlet{popred}{poppink}
\colorlet{popyellow}{popgold}
% Teintes claires (fonds de tableaux/encadrés) — le modèle nomme parfois
% les couleurs avec un suffixe L (ex. popblueL) sans les avoir définies.
\colorlet{poporangeL}{poporange!20}
\colorlet{popdarkL}{popdark!20}
\colorlet{poppurpleL}{poppurple!20}
\colorlet{popgreenL}{popgreen!20}
\colorlet{poppinkL}{poppink!20}
\colorlet{popblueL}{popblue!20}
\colorlet{popgoldL}{popgold!20}
\colorlet{popredL}{popred!20}
\colorlet{popgrayL}{popgray!20}
% Teintes claires pour fonds de tableaux / zones
\colorlet{popblueL}{popblue!10!white}
\colorlet{poporangeL}{poporange!14!white}
\colorlet{popgreenL}{popgreen!12!white}
\colorlet{poppurpleL}{poppurple!12!white}
\colorlet{poppinkL}{poppink!12!white}
\colorlet{popgoldL}{popgold!14!white}

\pagecolor{popcream}

% ============================================================
% Typographie
% ============================================================
\defaultfontfeatures{Ligatures=TeX}
\setmainfont{PlusJakartaSans-Medium.ttf}[Path=../../../fonts/,BoldFont=PlusJakartaSans-Bold.ttf]
% Maths en sans-serif (Fira Math) avec chiffres et lettres droites de Plus
% Jakarta, pour que les formules se fondent dans le texte.
\usepackage[math-style=ISO,bold-style=ISO]{unicode-math}
\setmathfont{FiraMath-Regular.otf}[Path=../../../fonts/]
\setmathfont{PlusJakartaSans-Medium.ttf}[Path=../../../fonts/,range={up/num,up/{Latin,latin},\mathup}]
\usepackage{icomma} % virgule décimale sans espace en mode math
% Caractères mathématiques Unicode absents des polices de texte (Plus Jakarta,
% Archivo Black) : rendus par leur équivalent mathématique, en texte comme en
% formule — sinon ils s'affichent en carré vide (ex. « Arithmétique dans ℤ »).
\usepackage{newunicodechar}
\newunicodechar{ℝ}{\ensuremath{\mathbb{R}}}
\newunicodechar{ℤ}{\ensuremath{\mathbb{Z}}}
\newunicodechar{ℂ}{\ensuremath{\mathbb{C}}}
\newunicodechar{ℕ}{\ensuremath{\mathbb{N}}}
\newunicodechar{ℚ}{\ensuremath{\mathbb{Q}}}
\newunicodechar{𝔻}{\ensuremath{\mathbb{D}}}
\newunicodechar{α}{\ensuremath{\alpha}}
\newunicodechar{β}{\ensuremath{\beta}}
\newunicodechar{γ}{\ensuremath{\gamma}}
\newunicodechar{δ}{\ensuremath{\delta}}
\newunicodechar{ε}{\ensuremath{\varepsilon}}
\newunicodechar{θ}{\ensuremath{\theta}}
\newunicodechar{λ}{\ensuremath{\lambda}}
\newunicodechar{μ}{\ensuremath{\mu}}
\newunicodechar{σ}{\ensuremath{\sigma}}
\newunicodechar{τ}{\ensuremath{\tau}}
\newunicodechar{φ}{\ensuremath{\varphi}}
\newunicodechar{ψ}{\ensuremath{\psi}}
\newunicodechar{ω}{\ensuremath{\omega}}
\newunicodechar{Σ}{\ensuremath{\Sigma}}
% majuscules produites par \MakeUppercase dans les titres (α-aminés -> Α)
\newunicodechar{Α}{\ensuremath{\alpha}}
\newunicodechar{Β}{\ensuremath{\beta}}
\newunicodechar{Γ}{\ensuremath{\Gamma}}
\newunicodechar{Δ}{\ensuremath{\Delta}}
\newunicodechar{Ε}{\ensuremath{\varepsilon}}
\newunicodechar{Θ}{\ensuremath{\Theta}}
\newunicodechar{Λ}{\ensuremath{\Lambda}}
\newunicodechar{Μ}{\ensuremath{\mu}}
\newunicodechar{Τ}{\ensuremath{\tau}}
\newunicodechar{Φ}{\ensuremath{\Phi}}
\newunicodechar{Ψ}{\ensuremath{\Psi}}
\newunicodechar{Ω}{\ensuremath{\Omega}}
\newunicodechar{↦}{\ensuremath{\mapsto}}
\newunicodechar{∈}{\ensuremath{\in}}
\newunicodechar{∉}{\ensuremath{\notin}}
\newunicodechar{∧}{\ensuremath{\wedge}}
\newunicodechar{⇔}{\ensuremath{\Leftrightarrow}}
\newunicodechar{⟺}{\ensuremath{\Longleftrightarrow}}
\newunicodechar{⇒}{\ensuremath{\Rightarrow}}
\newunicodechar{⟹}{\ensuremath{\Longrightarrow}}
\newunicodechar{∘}{\ensuremath{\circ}}
\newunicodechar{∪}{\ensuremath{\cup}}
\newunicodechar{∩}{\ensuremath{\cap}}
\newunicodechar{≡}{\ensuremath{\equiv}}
\newunicodechar{⊂}{\ensuremath{\subset}}
\newunicodechar{⊥}{\ensuremath{\perp}}
\newunicodechar{∃}{\ensuremath{\exists}}
\newunicodechar{∀}{\ensuremath{\forall}}
\newunicodechar{∛}{\ensuremath{\sqrt[3]{\,}}}
\newunicodechar{∜}{\ensuremath{\sqrt[4]{\,}}}
\newunicodechar{⃗}{\ensuremath{{}^{\to}}}
\newunicodechar{ⁿ}{\ifmmode^{n}\else\textsuperscript{\ensuremath{n}}\fi}
\newunicodechar{ˣ}{\ifmmode^{x}\else\textsuperscript{\ensuremath{x}}\fi}
\newunicodechar{ᵇ}{\ifmmode^{b}\else\textsuperscript{\ensuremath{b}}\fi}
\newunicodechar{ʳ}{\ifmmode^{r}\else\textsuperscript{\ensuremath{r}}\fi}
\newunicodechar{ᵘ}{\ifmmode^{u}\else\textsuperscript{\ensuremath{u}}\fi}
\newunicodechar{ʸ}{\ifmmode^{y}\else\textsuperscript{\ensuremath{y}}\fi}
\newunicodechar{ᵃ}{\ifmmode^{a}\else\textsuperscript{\ensuremath{a}}\fi}
\newunicodechar{ᵉ}{\ifmmode^{e}\else\textsuperscript{\ensuremath{e}}\fi}
\newunicodechar{ᵏ}{\ifmmode^{k}\else\textsuperscript{\ensuremath{k}}\fi}
\newunicodechar{ᵖ}{\ifmmode^{p}\else\textsuperscript{\ensuremath{p}}\fi}
\newunicodechar{ᴺ}{\ifmmode^{N}\else\textsuperscript{\ensuremath{N}}\fi}
\newunicodechar{ᶜ}{\ifmmode^{c}\else\textsuperscript{\ensuremath{c}}\fi}
\newunicodechar{ʰ}{\ifmmode^{h}\else\textsuperscript{\ensuremath{h}}\fi}
\newunicodechar{ᵠ}{\ifmmode^{q}\else\textsuperscript{\ensuremath{q}}\fi}
\newunicodechar{ₙ}{\ifmmode_{n}\else\textsubscript{\ensuremath{n}}\fi}
\newunicodechar{ₐ}{\ifmmode_{a}\else\textsubscript{\ensuremath{a}}\fi}
\newunicodechar{ᵢ}{\ifmmode_{i}\else\textsubscript{\ensuremath{i}}\fi}
\newunicodechar{ᵣ}{\ifmmode_{r}\else\textsubscript{\ensuremath{r}}\fi}
\newunicodechar{ₖ}{\ifmmode_{k}\else\textsubscript{\ensuremath{k}}\fi}
\newunicodechar{ᵦ}{\ifmmode_{\beta}\else\textsubscript{\ensuremath{\beta}}\fi}
\newunicodechar{ₓ}{\ifmmode_{x}\else\textsubscript{\ensuremath{x}}\fi}
\newfontfamily\popdisplay{ArchivoBlack-Regular.ttf}[Path=../../../fonts/]
\newfontfamily\poplogo{SpaceGrotesk-Bold.ttf}[Path=../../../fonts/]
\newfontfamily\popsemi{PlusJakartaSans-SemiBold.ttf}[Path=../../../fonts/]
\newfontfamily\popxbold{PlusJakartaSans-ExtraBold.ttf}[Path=../../../fonts/]
\newfontfamily\popxboldls{PlusJakartaSans-ExtraBold.ttf}[Path=../../../fonts/,LetterSpace=3.0]

\setlist[enumerate]{leftmargin=1.7em,itemsep=5pt,topsep=4pt}
\setlist[itemize]{leftmargin=1.7em,itemsep=4pt,topsep=4pt,label={\color{poporange}\textbullet}}
\setlength{\parindent}{0pt}
\setlength{\parskip}{5pt}
\renewcommand{\arraystretch}{1.25}

% pgfplots : style Pop par défaut
\pgfplotsset{
  every axis/.append style={axis line style={popink,line width=1pt},tick style={popink},
    grid style={popline,line width=0.5pt},label style={font=\small},tick label style={font=\footnotesize}},
  popcurve/.style={poporange,line width=1.8pt,no markers,smooth},
}
% Même style disponible dans un \draw TikZ simple (hors pgfplots)
\tikzset{popcurve/.style={poporange,line width=1.8pt}}
\tikzset{popgrid/.style={popline,very thin}}
\colorlet{popcurve}{poporange} % le nom du style est parfois utilisé comme couleur

% ============================================================
% Pill / squiggle
% ============================================================
\newcommand{\poppill}[3]{%
  \tikz[baseline=-0.55ex]\node[draw=popink,line width=1.2pt,rounded corners=2.6pt,%
    fill=#2,inner xsep=6.5pt,inner ysep=3pt]{%
    \popxboldls\fontsize{7}{7}\selectfont\color{#3}\MakeUppercase{#1}};%
}
\newcommand{\popsquiggle}[1]{%
  \tikz[baseline=-0.4ex]{
    \draw[#1,line width=2.6pt,line cap=round]
      plot [smooth] coordinates {(0,0.09) (0.34,0.01) (0.68,0.21) (1.02,0.01) (1.36,0.10)};
  }%
}
% Mot-clé surligné (marqueur orange sous le texte)
\newcommand{\cle}[1]{{\popxbold\color{popdark}#1}}

% ============================================================
% En-tête / pied
% ============================================================
\pagestyle{fancy}
\fancyhf{}
\renewcommand{\headrulewidth}{0pt}
\renewcommand{\footrulewidth}{0pt}
\lhead{\raisebox{-0.15ex}{\poplogo\fontsize{13}{13}\selectfont\color{popink}OnBuch\color{poporange}+}}
\rhead{\poppill{\DOCMATIERE\ \textbullet\ \DOCNIVEAU\ \textbullet\ Leçon \DOCLECON}{white}{popink}}
\lfoot{\popxbold\fontsize{7.6}{7.6}\selectfont\color{popmuted}\MakeUppercase{Cours signé OnBuch+}\ \textbullet\ {\popsemi\DOCTITRE}}
\rfoot{\tikz[baseline=-0.6ex]\node[rounded corners=8.5pt,fill=popink,minimum height=17pt,minimum width=17pt,inner xsep=5pt,inner sep=0]{\popxbold\fontsize{8}{8}\selectfont\color{popcream}\thepage\,/\,\pageref*{LastPage}};}

% ============================================================
% Titres
% ============================================================
\newcommand{\chaptitre}[2]{%
  \begin{tcolorbox}[enhanced,colback=poporange,colframe=popink,boxrule=2.6pt,arc=11pt,
      drop shadow=popink,left=15pt,right=15pt,top=12pt,bottom=11pt,before skip=3pt,after skip=18pt]
    \poppill{#2}{popink}{popcream}\par\vspace{9pt}
    {\raggedright\hyphenpenalty=10000\exhyphenpenalty=10000\popdisplay\fontsize{22}{24}\selectfont\color{popcream}\MakeUppercase{#1}\par}
  \end{tcolorbox}
}
% Section numérotée : pastille orange avec le numéro + titre Archivo
\newcounter{popsec}
\newcommand{\coursec}[1]{%
  \stepcounter{popsec}\setcounter{popsub}{0}%
  \par\vspace{16pt}\noindent
  \begin{minipage}{\linewidth}
  \tikz[baseline=-0.7ex]\node[draw=popink,line width=1.6pt,fill=poporange,rounded corners=4pt,minimum size=22pt,inner sep=0]{\popdisplay\fontsize{12}{12}\selectfont\color{popcream}\thepopsec};%
  \hspace{8pt}{\popdisplay\fontsize{15}{17}\selectfont\color{popink}\MakeUppercase{#1}}\par
  \vspace{4pt}\noindent{\color{poporange}\rule{36pt}{3.6pt}}
  \end{minipage}\par\nopagebreak\vspace{8pt}
}
\newcounter{popsub}
\newcommand{\courssub}[1]{%
  \stepcounter{popsub}%
  \par\vspace{9pt}\noindent{\popxbold\fontsize{11.5}{13}\selectfont\color{popink}{\color{poporange}\thepopsec.\thepopsub}\hspace{6pt}#1}\par\nopagebreak\vspace{4pt}
}

% ============================================================
% Boîtes pédagogiques — même recette : fond blanc, bordure épaisse colorée,
% ombre dure encre, badge plein. Titre optionnel [#1].
% ============================================================
\tcbset{popbase/.style={breakable,enhanced,colback=white,boxrule=2.3pt,arc=9pt,
  shadow={3pt}{-3pt}{0pt}{popink},left=14pt,right=14pt,top=11pt,bottom=11pt,before skip=11pt,after skip=15pt,
  fontupper=\popsemi\fontsize{10.6}{14.5}\selectfont\color{popink},
  fontlower=\popsemi\fontsize{10.6}{14.5}\selectfont\color{popink}}}
% Coche dessinée (le glyphe ✓ n'existe pas dans les polices du document)
\newcommand{\popcheck}[1]{\tikz[baseline=-0.5ex]\draw[#1,line width=1.6pt,line cap=round,line join=round](-0.13,0.0)--(-0.04,-0.09)--(0.13,0.1);}
\providecommand{\coche}{\popcheck{popgreen}}
\providecommand{\resultat}[1]{\par\smallskip\noindent\fcolorbox{popgreen}{popgreenL}{\parbox{\dimexpr\linewidth-2\fboxsep-2\fboxrule\relax}{\textbf{Résultat.} #1}}\par\smallskip}
\newcommand{\popboxhead}[3]{\poppill{#1}{#2}{white}\ifx\relax#3\relax\else\hspace{6pt}{\popxbold\fontsize{10.6}{12}\selectfont\color{popink}#3}\fi\par\vspace{7pt}}

\newtcolorbox{aretenir}[1][]{popbase,colframe=popgreen,before upper={\popboxhead{À retenir}{popgreen}{#1}}}
\newtcolorbox{exemplebox}[1][]{popbase,colframe=poporange,before upper={\popboxhead{Exemple}{poporange}{#1}}}
\newtcolorbox{definition}[1][]{popbase,colframe=popblue,before upper={\popboxhead{Définition}{popblue}{#1}}}
\newtcolorbox{propriete}[1][]{popbase,colframe=poppurple,before upper={\popboxhead{Propriété}{poppurple}{#1}}}
\newtcolorbox{methode}[1][]{popbase,colframe=popgold,before upper={\popboxhead{Méthode}{popgold}{#1}}}
\newtcolorbox{attention}[1][]{popbase,colframe=poppink,before upper={\popboxhead{Attention}{poppink}{#1}}}
\newtcolorbox{experience}[1][]{popbase,colframe=popblue,colback=popblueL,before upper={\popboxhead{Expérience}{popblue}{#1}}}
% « Le savais-tu ? » : carte violette pleine (contexte, culture, Cameroun)
\newtcolorbox{savaistu}[1][]{popbase,colback=poppurple,colframe=popink,
  fontupper=\popsemi\fontsize{10.4}{14.2}\selectfont\color{white},
  before upper={\poppill{Le savais-tu\,?}{white}{poppurple}\ifx\relax#1\relax\else\hspace{6pt}{\popxbold\color{white}#1}\fi\par\vspace{7pt}}}
% Exercice résolu : énoncé en haut, \tcblower puis la solution sur fond crème
\newcounter{popexo}\newcounter{popexr}
\newtcolorbox{exoresolu}[1][]{popbase,colframe=poporange,colbacklower=popcream,
  segmentation style={popink,line width=1.2pt,dashed},
  before upper={\stepcounter{popexr}\popboxhead{Exercice résolu \thepopexr}{poporange}{#1}},
  before lower={\poppill{Solution}{popink}{popcream}\par\vspace{6pt}}}
% Exercice (non résolu en place) — la correction est dans la partie Corrigés
\newtcolorbox{exercice}[1][]{popbase,colframe=popink,
  before upper={\stepcounter{popexo}\popboxhead{Exercice \thepopexo}{popink}{#1}}}
% Corrigé
\newtcolorbox{corrige}[1][]{popbase,colframe=popgreen,colback=popgreenL,
  before upper={\popboxhead{Corrigé}{popgreen}{#1}}}
% Objectifs / prérequis / bilan
\newtcolorbox{objectifs}{popbase,colframe=popink,colback=poporangeL,
  before upper={\popboxhead{Objectifs de la leçon}{popink}{}}}
\newtcolorbox{prerequis}{popbase,colframe=popink,colback=popblueL,
  before upper={\popboxhead{Prérequis}{popblue}{}}}
\newtcolorbox{bilan}{popbase,colframe=popink,colback=white,boxrule=2.8pt,
  before upper={\popboxhead{Fiche bilan}{poporange}{L'essentiel à connaître pour l'examen}}}
% Figure encadrée avec légende
\newtcolorbox{popfigure}[1][]{enhanced,breakable=false,colback=white,colframe=popline,boxrule=1.4pt,arc=8pt,
  left=10pt,right=10pt,top=10pt,bottom=8pt,before skip=10pt,after skip=12pt,halign=center,
  lower separated=false,halign lower=center,
  fontlower=\popsemi\fontsize{9}{11.5}\selectfont\color{popmuted},
  #1}
\newcommand{\legende}[1]{\par\vspace{6pt}{\popsemi\fontsize{9}{11.5}\selectfont\color{popmuted}#1}}

% ============================================================
% Signature OnBuch+
% ============================================================
\newcommand{\poplogoinline}[1]{{\poplogo\fontsize{#1}{#1}\selectfont\color{popink}OnBuch\color{poporange}+}}
\newcommand{\signatureonbuch}{%
  \begin{tcolorbox}[enhanced,colback=popink,colframe=popink,boxrule=2.2pt,arc=10pt,
      drop shadow=poporange,left=14pt,right=14pt,top=10pt,bottom=10pt,before skip=0pt,after skip=0pt]
    \tikz[baseline=-0.7ex]\node[circle,fill=popgreen,draw=popcream,line width=1.3pt,minimum size=22pt,inner sep=0]{\popcheck{white}};%
    \hspace{9pt}\begin{minipage}[c]{0.62\linewidth}
      {\popxbold\fontsize{12}{13}\selectfont\color{popcream}Signé\ \textbullet\ {\poplogo OnBuch\color{poporange}+}}\par\vspace{2pt}
      {\raggedright\popsemi\fontsize{8}{10.5}\selectfont\color{popline}\MakeUppercase{Équipe pédagogique OnBuch+}\\\MakeUppercase{Conforme au programme officiel MINESEC}\par}
    \end{minipage}\hfill
    {\poplogo\fontsize{22}{22}\selectfont\color{popcream}OnBuch\color{poporange}+}
  \end{tcolorbox}%
}

% ============================================================
% Page de garde
% #1 titre · #2 matière · #3 niveau · #4 module · #5 n° leçon · #6 durée ·
% #7 description / objectifs (texte ou itemize)
% ============================================================
\newcommand{\pagegarde}[7]{%
  \thispagestyle{empty}
  \newgeometry{a4paper,margin=1.7cm,top=1.6cm,bottom=1.4cm}%
  \begin{tikzpicture}[remember picture,overlay]
    \fill[poporange!18] ([xshift=-3.2cm,yshift=-3.6cm]current page.north east) circle (4.6cm);
    \fill[poppurple!14] ([xshift=2.4cm,yshift=7.6cm]current page.south west) circle (3.8cm);
  \end{tikzpicture}%
  {\poplogo\fontsize{30}{30}\selectfont\color{popink}OnBuch\color{poporange}+}\hfill
  \raisebox{0.6ex}{\poppill{Cours signé \textbullet\ Premium}{popink}{popcream}}\par
  \vspace{1pt}\popsquiggle{poppurple}\par
  \vspace{8pt}
  \begin{tcolorbox}[enhanced,colback=poporange,colframe=popink,boxrule=2.8pt,arc=13pt,
      drop shadow=popink,left=18pt,right=18pt,top=12pt,bottom=13pt,after skip=10pt]
    \poppill{#2 \textbullet\ #3}{popink}{popcream}\hspace{5pt}\poppill{#4}{white}{popink}\par\vspace{8pt}
    {\popdisplay\fontsize{14}{14}\selectfont\color{popink}LEÇON #5}\par\vspace{4pt}
    {\raggedright\hyphenpenalty=10000\exhyphenpenalty=10000\popdisplay\fontsize{25}{27}\selectfont\color{popcream}\MakeUppercase{#1}\par}
    \vspace{8pt}
    {\popxbold\fontsize{9.5}{9.5}\selectfont\color{popink}\MakeUppercase{Durée conseillée : #6}}
  \end{tcolorbox}
  \begin{tcolorbox}[enhanced,colback=white,colframe=popink,boxrule=2.2pt,arc=10pt,
      drop shadow=popink,left=16pt,right=16pt,top=10pt,bottom=10pt,after skip=8pt]
    \poppill{Dans cette leçon}{poppurple}{white}\par\vspace{5pt}
    \popsemi\fontsize{9.3}{12.6}\selectfont\color{popink}#7
  \end{tcolorbox}
  \noindent\begin{minipage}[t]{0.487\linewidth}
    \begin{tcolorbox}[enhanced,colback=popgreen,colframe=popink,boxrule=2.2pt,arc=10pt,
        drop shadow=popink,left=11pt,right=11pt,top=10pt,bottom=10pt,height=3.1cm,valign=center]
      \tcbox[colback=white,colframe=popink,boxrule=1.3pt,arc=5pt,boxsep=0pt,nobeforeafter,
        left=3pt,right=3pt,top=3pt,bottom=3pt,valign=center]{\qrcode[height=1.9cm]{https://play.google.com/store/apps/details?id=com.luvvix.onbuch&hl=fr}}%
      \hspace{8pt}\begin{minipage}[c]{0.5\linewidth}
        {\popdisplay\fontsize{11}{12.5}\selectfont\color{white}\MakeUppercase{Télécharge OnBuch}}\par\vspace{4pt}
        {\raggedright\popsemi\fontsize{8.4}{11}\selectfont\color{white}Gratuit \textbullet\ Google Play\\Tuteur IA, annales, concours, résultats\par}
      \end{minipage}
    \end{tcolorbox}
  \end{minipage}\hfill
  \begin{minipage}[t]{0.487\linewidth}
    \begin{tcolorbox}[enhanced,colback=poppurple,colframe=popink,boxrule=2.2pt,arc=10pt,
        drop shadow=popink,left=11pt,right=11pt,top=10pt,bottom=10pt,height=3.1cm,valign=center]
      \tikz[baseline=-0.7ex]\node[circle,fill=white,draw=popink,line width=1.3pt,minimum size=1.6cm,inner sep=0]{\popdisplay\fontsize{24}{24}\selectfont\color{poppurple}+};%
      \hspace{8pt}\begin{minipage}[c]{0.6\linewidth}
        {\popdisplay\fontsize{11}{12.5}\selectfont\color{white}\MakeUppercase{Passe à OnBuch+}}\par\vspace{4pt}
        {\raggedright\popsemi\fontsize{8.4}{11}\selectfont\color{white}Tous les cours signés, Tuteur IA illimité, fiches PDF — onglet OnBuch+ de l'app\par}
      \end{minipage}
    \end{tcolorbox}
  \end{minipage}
  \vspace{8pt}
  \popcontenu
  \vfill
  \signatureonbuch
  \restoregeometry
  \newpage
}

% Rangée « Ce cours contient » (page de garde)
\newcommand{\poptile}[3]{% #1 couleur #2 gros texte #3 légende
  \begin{tcolorbox}[enhanced,colback=#1,colframe=popink,boxrule=2pt,arc=9pt,shadow={2.5pt}{-2.5pt}{0pt}{popink},
      left=6pt,right=6pt,top=9pt,bottom=9pt,halign=center,valign=center,height=1.75cm,width=\linewidth,nobeforeafter]
    {\popdisplay\fontsize{12.5}{13}\selectfont\color{white}\MakeUppercase{#2}}\par\vspace{3pt}
    {\centering\popsemi\fontsize{7.8}{9.5}\selectfont\color{white}#3\par}
  \end{tcolorbox}}
\newcommand{\popcontenu}{%
  \par\noindent{\popxboldls\fontsize{7.5}{7.5}\selectfont\color{popmuted}CE COURS SIGNÉ CONTIENT}\par\vspace{7pt}
  \noindent\begin{minipage}[t]{0.235\linewidth}\poptile{popblue}{Cours}{complet, illustré, pas à pas}\end{minipage}\hfill
  \begin{minipage}[t]{0.235\linewidth}\poptile{poporange}{Méthodes}{et exercices résolus}\end{minipage}\hfill
  \begin{minipage}[t]{0.235\linewidth}\poptile{poppink}{Exercices}{type examen + corrigés}\end{minipage}\hfill
  \begin{minipage}[t]{0.235\linewidth}\poptile{popgreen}{Fiche bilan}{l'essentiel à retenir}\end{minipage}\par}

% Sommaire
\newtcolorbox{sommaire}{enhanced,colback=white,colframe=popink,boxrule=2.2pt,arc=10pt,
  drop shadow=popink,left=18pt,right=18pt,top=13pt,bottom=13pt,before skip=6pt,after skip=20pt,
  before upper={\poppill{Sommaire}{poppurple}{white}\par\vspace{9pt}\popsemi\fontsize{10.6}{14.5}\selectfont\color{popink}}}

% Signature de fin de cours — poussée tout en bas de la dernière page
\newcommand{\findecours}{%
  \par\vfill\nopagebreak
  \begin{center}\popsquiggle{poporange}\end{center}\vspace{4pt}
  \signatureonbuch
}
\AtBeginDocument{\def\llbracket{\mathopen{[\mkern-3.5mu[}}\def\rrbracket{\mathclose{]\mkern-3.5mu]}}} % intervalles d'entiers (glyphe ⟦ absent de la police)
% Glyphes absents de Fira Math : redéfinis à partir de glyphes disponibles
\AtBeginDocument{%
  \def\setminus{\mathbin{\backslash}}\let\smallsetminus\setminus
  \def\vdots{\mathord{\vbox{\baselineskip3.2pt\lineskiplimit0pt\kern4pt\hbox{.}\hbox{.}\hbox{.}}}}%
  \def\ddots{\mathinner{\mkern1mu\raise6.5pt\hbox{.}\mkern2mu\raise3.6pt\hbox{.}\mkern2mu\raise0.7pt\hbox{.}\mkern1mu}}%
  \def\triangle{\mathord{\scalebox{0.9}{$\symup{\Delta}$}}}%
  \def\nabla{\mathord{\raisebox{\depth}{\scalebox{1}[-1]{$\symup{\Delta}$}}}}%
  \def\checkmark{\popcheck{popdark}}%
  \def\star{\mathbin{\ast}}\def\bigstar{\mathord{\ast}}%
  \def\diamond{\mathbin{\scalebox{0.6}{$\Diamond$}}}%
  \def\bigcup{\mathop{\mathchoice{\raisebox{-0.3em}{\scalebox{1.7}{$\cup$}}}{\scalebox{1.25}{$\cup$}}{\cup}{\cup}}\displaylimits}%
  \def\bigcap{\mathop{\mathchoice{\raisebox{-0.3em}{\scalebox{1.7}{$\cap$}}}{\scalebox{1.25}{$\cap$}}{\cap}{\cap}}\displaylimits}%
  \def\lesseqqgtr{\mathrel{\lessgtr}}\def\bigoplus{\mathop{\oplus}\displaylimits}%
}
\providecommand{\ssi}{si et seulement si} % abréviation fréquente
\AtBeginDocument{\providecommand{\wideparen}{\overparen}} % arc de cercle

%% FIGFIX-BEGIN v5 — correctifs globaux des figures (docs/onbuch-generation/tools/figfix)
\usepackage{adjustbox}\usepackage{etoolbox}\tcbuselibrary{raster}
% toute figure TikZ/pgfplots plus large que la ligne est réduite à la largeur disponible
\BeforeBeginEnvironment{tikzpicture}{\begin{adjustbox}{max width=\linewidth}}
\AfterEndEnvironment{tikzpicture}{\end{adjustbox}}
% légendes pgfplots lisibles par-dessus les courbes
\pgfplotsset{every axis/.append style={legend style={fill=white,fill opacity=0.92,text opacity=1,draw=black!15,inner sep=3pt}}}
% étiquettes d'arêtes (syntaxe edge["..."]) avec fond blanc
\tikzset{every edge quotes/.append style={fill=white,inner sep=1.2pt}}
%% FIGFIX-END
\begin{document}

\pagegarde{La science}{Philosophie}{Tle Tech}{Philosophie – classe de Terminale technique}{N10}{4 séances (fiche de progression harmonisée)}{Cette leçon interroge la nature de la science, ses différents types et les conditions d'émergence de l'esprit scientifique face aux obstacles épistémologiques. Elle évalue ensuite la valeur de la science, de la méthode expérimentale aux problèmes bioéthiques et écologiques qui se posent aujourd'hui au Cameroun et dans le monde.
\vspace{4pt}
\begin{multicols}{2}\small
\begin{itemize}
\item À la fin de cette leçon, je sais définir la science et distinguer ses caractères essentiels (objectivité, rationalité, vérifiabilité).
\item À la fin de cette leçon, je sais distinguer et caractériser les différents types de sciences : mathématiques, sciences de la matière et sciences humaines.
\item À la fin de cette leçon, je sais expliquer les conditions historiques et intellectuelles d'émergence de l'esprit scientifique.
\item À la fin de cette leçon, je sais identifier et illustrer les principaux obstacles épistémologiques selon Bachelard.
\item À la fin de cette leçon, je sais décrire les étapes de la méthode expérimentale et l'appliquer à un cas concret.
\item À la fin de cette leçon, je sais discuter la valeur de la science (utilité, limites, responsabilité).
\end{itemize}\end{multicols}}

\begin{sommaire}
\begin{multicols}{2}
\begin{enumerate}
\item Qu'est-ce que la science ? Nature et caractères
\item Les différents types de sciences
\item L'émergence de l'esprit scientifique et les obstacles épistémologiques
\item La méthode expérimentale et ses applications
\item La valeur de la science : utilité, limites et problèmes
\item Science, bioéthique et question écologique
\item Activité d'intégration
\item Méthodes et exercices résolus
\item Exercices
\item Corrigés des exercices
\item Fiche bilan
\end{enumerate}
\end{multicols}
\end{sommaire}

\begin{objectifs}
\begin{itemize}
\item À la fin de cette leçon, je sais définir la science et distinguer ses caractères essentiels (objectivité, rationalité, vérifiabilité).
\item À la fin de cette leçon, je sais distinguer et caractériser les différents types de sciences : mathématiques, sciences de la matière et sciences humaines.
\item À la fin de cette leçon, je sais expliquer les conditions historiques et intellectuelles d'émergence de l'esprit scientifique.
\item À la fin de cette leçon, je sais identifier et illustrer les principaux obstacles épistémologiques selon Bachelard.
\item À la fin de cette leçon, je sais décrire les étapes de la méthode expérimentale et l'appliquer à un cas concret.
\item À la fin de cette leçon, je sais discuter la valeur de la science (utilité, limites, responsabilité).
\item À la fin de cette leçon, je sais analyser des problèmes de bioéthique et la question écologique à partir de situations camerounaises.
\item À la fin de cette leçon, je sais rédiger et justifier une démarche argumentée, une réponse ou une conclusion en dissertation philosophique.
\end{itemize}
\end{objectifs}

\begin{prerequis}
\begin{itemize}
\item Notion de vérité et distinction opinion/connaissance (leçon sur la vérité, début d'année).
\item Notion de raison et de rationalité (leçon sur la raison et l'expérience).
\item Notions de technique et de progrès (leçon sur la technique).
\item Repères d'histoire des sciences vus en classe de Première (Galilée, Descartes).
\item Capacité à lire et interpréter un document (texte, tableau de données, photographie).
\end{itemize}
\end{prerequis}

\newpage

\coursec{Qu'est-ce que la science ? Nature et caractères}

\courssub{Accroche et problématique}

À Yaoundé, un jeune diplômé hésite : sa grand-mère lui conseille un traitement traditionnel contre le paludisme, à base d'écorces d'arbres et de prières aux ancêtres, tandis que le médecin du Centre Hospitalier Universitaire lui prescrit un médicament testé en laboratoire, dont la posologie est calibrée au milligramme près. Lequel des deux a raison, et \emph{pourquoi} ? Par ailleurs, les téléphones mobiles conçus grâce à la science inondent le marché camerounais, facilitant la communication et le commerce, mais leurs déchets électroniques polluent les quartiers de Douala et les cours d'eau. La science est-elle toujours un bien ? Qu'est-ce qui distingue fondamentalement un savoir scientifique d'une simple opinion, d'une croyance ou d'un savoir-faire technique ? C'est à ces questions cruciales que cette première section va répondre, en posant les fondements épistémologiques de la notion de science.

\courssub{Étymologie et définition rigoureuse}

Le mot « science » vient du latin \emph{scientia}, qui signifie « savoir », « connaissance », dérivé du verbe \emph{scire} (savoir). Mais dans l'acception philosophique et épistémologique moderne, la science ne désigne pas n'importe quel savoir. Elle désigne un \textbf{savoir rationnel, méthodique, objectif et vérifiable}, constitué de propositions logiquement liées les unes aux autres et portant sur un domaine déterminé de la réalité.

\begin{definition}[La science]
La science est un ensemble de connaissances \textbf{exactes}, \textbf{méthodiques}, \textbf{objectives} et \textbf{vérifiables}, portant sur un domaine déterminé de la réalité, obtenues par une démarche rationnelle rigoureuse et soumises au contrôle de l'expérience et de la raison.
\end{definition}

Contrairement à l'opinion (\emph{doxa}) qui est subjective, instable et non fondée sur la preuve, la science vise l'universalité et la nécessité. Elle ne se contente pas de décrire les faits (le « comment »), elle cherche à les expliquer par des lois et des théories (le « pourquoi »).

\courssub{L'opposition science / opinion (doxa) : le critère de la preuve}

L'opposition fondamentale, déjà tracée par Platon dans \emph{La République}, distingue l'\emph{épistémè} (science, savoir vrai) de la \emph{doxa} (opinion, croyance). L'opinion repose sur le sensible, l'apparence, le préjugé ou l'autorité de la tradition. Elle est variable selon les individus, les époques et les cultures. La science, elle, repose sur la \textbf{preuve} (logique ou expérimentale) et prétend à l'universalité.

Prenons un exemple concret camerounais : les causes du paludisme.
\begin{itemize}
    \item \textbf{Opinion / Croyance locale :} « Le paludisme vient de la "mauvaise air" (\emph{mal aria}) des marécages, ou d'un mauvais sort jeté par un voisin jaloux, ou encore de la colère des ancêtres. » Ce savoir est transmis oralement, non testé, non mesuré. Il varie d'un village à l'autre.
    \item \textbf{Science (médecine) :} Le paludisme est une maladie parasitaire causée par un protozoaire du genre \emph{Plasmodium} (falciparum, vivax, etc.), transmis à l'homme par la piqûre d'un moustique femelle \emph{Anopheles} infecté. Ce savoir repose sur l'observation microscopique (Laveran, 1880), l'identification du vecteur (Ross, 1897), la reproductibilité du diagnostic (goutte épaisse, test rapide) et l'efficacité thérapeutique vérifiée statistiquement (artémisinine, ACT).
\end{itemize}

\begin{popfigure}
\begin{tikzpicture}[
    box/.style={draw=popblue, thick, rounded corners, fill=popblueL, minimum height=1.2cm, text width=3.8cm, align=center, font=\small},
    header/.style={draw=popdark, thick, rounded corners, fill=popdark, text=white, minimum height=0.9cm, text width=3.8cm, align=center, font=\small\bfseries},
    crit/.style={draw=poporange, thick, rounded corners, fill=poporangeL, minimum height=1.0cm, text width=2.5cm, align=center, font=\small\bfseries},
    cell/.style={draw=popline, minimum height=1.0cm, text width=3.8cm, align=center, font=\small},
    cellc/.style={draw=popline, minimum height=1.0cm, text width=2.5cm, align=center, font=\small\itshape},
]
% En-têtes
\node[header] (op) at (0,0) {Opinion / Croyance (\emph{Doxa})};
\node[header] (sc) at (8,0) {Science (\emph{Epistémè})};
\node[crit] (cr) at (4,0) {Critères de distinction};

% Ligne 1 : Fondement
\node[crit] (c1) at (4,-1.1) {Fondement};
\node[cell] (o1) at (0,-1.1) {Tradition, autorité, intuition, sentiment, "on dit", texte sacré};
\node[cell] (s1) at (8,-1.1) {Observation rigoureuse, expérience contrôlée, raisonnement logique, calcul};

% Ligne 2 : Preuve
\node[crit] (c2) at (4,-2.2) {Preuve / Validation};
\node[cell] (o2) at (0,-2.2) {Aucune preuve exigée ; adhésion par foi, habitude ou conformisme social};
\node[cell] (s2) at (8,-2.2) {Preuve intersubjective : reproductibilité, falsifiabilité (Popper), consensus de la communauté scientifique};

% Ligne 3 : Universalité
\node[crit] (c3) at (4,-3.3) {Portée / Universalité};
\node[cell] (o3) at (0,-3.3) {Relative, subjective, changeante selon les cultures, les époques, les individus};
\node[cell] (s3) at (8,-3.3) {Universelle : vraie pour tous, en tout lieu, en tout temps (dans les conditions définies)};

% Ligne 4 : Erreur
\node[crit] (c4) at (4,-4.4) {Rapport à l'erreur};
\node[cell] (o4) at (0,-4.4) {L'erreur est niée ou sacralisée ; dogmatisme};
\node[cell] (s4) at (8,-4.4) {L'erreur est moteur de progrès ; révision constante, cumulativité (Bachelard)};

% Ligne 5 : Exemple
\node[crit] (c5) at (4,-5.5) {Exemple : Paludisme};
\node[cell] (o5) at (0,-5.5) {"Mauvais air", sorcellerie, déséquilibre humoral, punition divine};
\node[cell] (s5) at (8,-5.5) {Plasmodium + vecteur Anophèle $\to$ diagnostic biologique $\to$ traitement ciblé (ACT)};
\end{tikzpicture}
\legende{Tableau comparatif : Opinion / Croyance vs Science. La science se distingue par son fondement rationnel, sa validation par la preuve, son universalité et sa capacité d'autocorrection.}
\end{popfigure}

\begin{attention}[Piège à éviter : Confondre conviction et certitude scientifique]
Ne confondez pas la \textbf{certitude psychologique} (je suis \emph{sûr} de moi, j'y crois fermement) et la \textbf{certitude épistémologique} (cette proposition est démontrée et vérifiable par tous). Un traditionaliste peut être intimement convaincu que l'écorce d'arbre guérit le paludisme ; sa conviction est sincère mais non scientifique tant que l'efficacité n'a pas été isolée, dosée, testée en double aveugle contre placebo et reproduite. La science exige la \textbf{validation intersubjective} : n'importe quel chercheur compétent, n'importe où dans le monde, doit pouvoir refaire l'expérience et obtenir le même résultat.
\end{attention}

\courssub{Les caractères fondamentaux de la science}

La science ne se définit pas seulement par opposition à l'opinion ; elle se caractérise positivement par un ensemble de traits indissociables. On retient classiquement six caractères majeurs :

\begin{enumerate}
    \item \textbf{Objectivité :} Le savoir scientifique est indépendant du sujet qui le produit. Le résultat ne dépend pas de la personnalité, des désirs, de la nationalité ou de la religion du savant. L'eau bout à 100°C (au niveau de la mer) que l'expérience soit menée à Yaoundé, Paris ou Tokyo, par un chrétien, un musulman ou un athée. La \emph{méthode} (protocole, mesure, calcul) garantit cette objectivité en éliminant le subjectif.
    \item \textbf{Universalité :} Corollaire de l'objectivité. Une loi scientifique vaut pour tous les objets de même nature, en tout lieu et en tout temps. La loi de la gravitation universelle (Newton) ou la structure de l'ADN (Watson, Crick, Franklin) sont valides au Cameroun comme en Antarctique.
    \item \textbf{Rationalité :} La science procède par concepts clairs, définis, et par raisonnements logiques (déduction, induction, abduction). Elle refuse le contradictoire (principe de non-contradiction) et exige la cohérence interne des théories.
    \item \textbf{Méthode :} La science n'est pas un amas de faits bruts ; elle est \emph{méthodique}. Elle suit des règles précises : observation, hypothèse, expérimentation, modélisation, publication, relecture par les pairs (\emph{peer review}). C'est ce qui la distingue du simple « essai-erreur » artisanal.
    \item \textbf{Vérifiabilité (et falsifiabilité) :} Toute proposition scientifique doit pouvoir être confrontée à l'expérience. Pour Karl Popper, une théorie est scientifique si elle est \textbf{falsifiable}, c'est-à-dire s'il existe une expérience possible qui pourrait la contredire. « Demain il pleuvra ou il ne pleuvra pas » n'est pas scientifique (non falsifiable). « Demain il pleuvra 15 mm à Yaoundé » est scientifique (vérifiable/falsifiable).
    \item \textbf{Cumulativité (et révisibilité) :} Le savoir scientifique s'accumule et se corrige. « La science progresse par rectifications successives » (Gaston Bachelard). Newton n'a pas été « détruit » par Einstein ; sa mécanique a été \emph{incluse} comme cas limite (vitesse faible) dans la relativité générale. Au Cameroun, les protocoles de traitement du paludisme ont évolué : chloroquine (années 60-80) $\to$ résistance $\to$ sulfadoxine-pyriméthamine $\to$ résistance $\to$ thérapies combinées à base d'artémisinine (ACT) actuelles. Chaque étape corrige la précédente.
\end{enumerate}

\begin{exemplebox}[La cumulativité en action : l'évolution du traitement du paludisme au Cameroun]
\begin{itemize}
    \item \textbf{Années 1950-1980 :} Chloroquine (résidu de la guerre). Efficace, peu coûteuse. \textbf{Problème :} Apparition de résistances de \emph{P. falciparum} (sélection naturelle).
    \item \textbf{Années 1990-2000 :} Sulfadoxine-Pyriméthamine (Fansidar). \textbf{Problème :} Résistance rapide, effets secondaires.
    \item \textbf{Depuis 2004 (OMS/Cameroun) :} Thérapies Combinées à base d'Artémisinine (ACT : Artemether-Luméfantrine, Artesunate-Amodiaquine). \textbf{Avancée :} Combinaison de deux molécules aux mécanismes différents pour retarder la résistance. Diagnostic biologique obligatoire avant traitement (TDR / goutte épaisse).
\end{itemize}
\emph{Ce n'est pas un "retour en arrière", c'est une correction progressive guidée par l'épidémiologie, la génétique du parasite et la pharmacologie.}
\end{exemplebox}

\courssub{La science comme connaissance par les causes (Aristote)}

Pour Aristote, dans la \emph{Métaphysique} et les \emph{Seconds Analytiques}, « \textbf{savoir, c'est connaître par les causes} » (\emph{scire est per causas cognoscere}). L'opinion se contente du \emph{quia} (le fait que : « il a fièvre »), la science atteint le \emph{propter quid} (la raison pourquoi : « il a fièvre \emph{parce que} le plasmodium détruit ses globules rouges et libère des toxines pyrogènes »).

Aristote distingue quatre causes, que la science moderne réinterprète :
\begin{itemize}
    \item \textbf{Cause matérielle :} De quoi est fait l'objet ? (Le parasite, l'ADN, le moustique).
    \item \textbf{Cause formelle :} Quelle est sa structure, sa loi ? (Le cycle de vie du plasmodium, le code génétique).
    \item \textbf{Cause efficiente :} Quel est l'agent du changement ? (La piqûre de l'anophèle, la réplication cellulaire).
    \item \textbf{Cause finale :} À quelle fin ? (Aristote : la nature agit pour une fin. Science moderne : \emph{fonction} ou \emph{avantage sélectif} — le parasite se reproduit).
\end{itemize}
La science moderne, depuis Galilée, privilégie les causes \emph{efficientes} et \emph{formelles} (lois mathématiques, mécanismes physico-chimiques) et met entre parenthèses la cause finale (finalité) pour la nature inanimée, tout en la réintégrant en biologie évolutive (fonction adaptative).

\courssub{Distinction science / technique : le vrai et l'efficace}

C'est une distinction classique (Aristote, Descartes, Heidegger) mais leur relation est dialectique.

\begin{center}
\begin{tabularx}{\linewidth}{|l|X|X|}\hline
\rowcolor{popblueL} \textbf{Critère} & \textbf{Science (\emph{Epistémè})} & \textbf{Technique (\emph{Tekhnè})} \\ \hline
\textbf{But} & Connaître le \textbf{Vrai} (comprendre, expliquer, prévoir) & Produire l'\textbf{Utile} / l'\textbf{Efficace} (agir, transformer, fabriquer) \\ \hline
\textbf{Critère de réussite} & Vérité, cohérence, preuve, reproductibilité & Efficacité, rendement, coût, satisfaction du besoin \\ \hline
\textbf{Statut de l'erreur} & L'erreur est un obstacle à la connaissance & L'erreur (échec technique) est une information pour améliorer le prototype \\ \hline
\textbf{Exemple} & Découvrir la structure de la théobromine dans le cacao ; comprendre la génétique du cacaoyer & Sélectionner des variétés hybrides résistantes aux maladies (cacao "Mercédés", "Big Five") ; optimiser la fermentation \\ \hline
\end{tabularx}
\end{center}

\begin{attention}[Erreur fréquente : Science et technique ne s'identifient pas]
\begin{itemize}
    \item \textbf{Technique sans science :} Les savoir-faire traditionnels (forge, poterie, pharmacopée traditionnelle, architecture en banco) ont existé des millénaires sans théorie scientifique explicite. Ils relèvent de l'\emph{expérience accumulée} et de la \emph{transmission orale}.
    \item \textbf{Science sans technique (immédiate) :} La physique des particules (boson de Higgs), l'astronomie (trous noirs), les mathématiques pures (théorème de Fermat) n'ont pas d'application technique directe ou immédiate. Elles visent la compréhension.
    \item \textbf{Interaction féconde :} La science fournit des théories à la technique (thermodynamique $\to$ machine à vapeur ; génétique $\to$ OGM / sélection variétale) ; la technique fournit des instruments à la science (microscope $\to$ biologie cellulaire ; accélérateur $\to$ physique des particules ; séquençage $\to$ génomique).
\end{itemize}
\end{attention}

\begin{savaistu}[Le cacao camerounais : science et technique main dans la main]
Le Cameroun est le 5\textsuperscript{e} producteur mondial de cacao. L'Institut de Recherche Agricole pour le Développement (IRAD) mène des travaux \textbf{scientifiques} (génomique, résistance aux maladies comme le chancre phyllostictique ou le miride) pour créer des variétés améliorées. Ces variétés sont ensuite diffusées par la \textbf{technique} (pépinières, greffage, vulgarisation agricole) auprès des planteurs du Sud, du Centre et du Littoral. Résultat : le rendement passe de 300-400 kg/ha (variétés traditionnelles) à plus de 1,5-2 tonnes/ha (variétés améliorées + bonnes pratiques). La science éclaire la technique ; la technique valide la science sur le terrain.
\end{savaistu}

\courssub{La science n'est pas un savoir absolu : révisabilité et perfectibilité}

Un dogme est une vérité intangible, imposée par une autorité, non discutable. La science est l'exact opposé : \textbf{tout savoir scientifique est provisoire, révisable, perfectible}. Cela ne signifie pas qu'il est « faux » ou « relatif » (relativisme), mais qu'il est \textbf{le meilleur savoir disponible à un instant T}, corrigeable par de nouvelles observations, de meilleurs instruments ou des théories plus englobantes.

\begin{exoresolu}[Analyse épistémologique : Pourquoi le traitement du paludisme a-t-il changé ?]
\textbf{Énoncé :} Expliquer, en mobilisant les notions de \emph{falsifiabilité}, \emph{cumulativité} et \emph{obstacle épistémologique}, pourquoi le protocole national de prise en charge du paludisme au Cameroun a remplacé la chloroquine par les ACT (Artemisinin-based Combination Therapies).
\tcblower
\textbf{Solution détaillée :}
\begin{enumerate}
    \item \textbf{Obstacle épistémologique (Bachelard) :} L'usage massif et unique de la chloroquine a constitué un « obstacle » : on croyait le problème résolu une fois pour toutes (illusion de la solution définitive). La réalité biologique (évolution darwinienne du parasite) a contredit cette certitude.
    \item \textbf{Falsifiabilité (Popper) :} L'hypothèse « La chloroquine guérit tout paludisme » était falsifiable. L'observation d'échecs thérapeutiques répétés (résistance gène \emph{pfcrt}) a \emph{falsifié} cette hypothèse. Elle a été rejetée comme vérité universelle.
    \item \textbf{Cumulativité :} On n'a pas jeté tout le savoir antérieur. On a gardé : le diagnostic biologique, la connaissance du cycle du parasite, la pharmacocinétique. On a \emph{ajouté} : la combinaison de molécules (artémisinine + partenaire) pour réduire la pression de sélection (probabilité double mutation quasi nulle). Le nouveau protocole \emph{englobe} l'ancien (la chloroquine reste utile pour \emph{P. vivax} ou \emph{P. ovale} dans certaines zones).
    \item \textbf{Conclusion :} Le changement de protocole n'est pas un aveu d'échec de la science, mais la \emph{preuve de son fonctionnement normal} : autocorrection par l'expérience, progrès par rectification.
\end{enumerate}
\end{exoresolu}

\begin{exemplebox}[Le "malaria" : de l'opinion à la science]
Le mot \emph{malaria} vient de l'italien \emph{mala aria} (« mauvais air »). Au Moyen Âge et jusqu'au XIX\textsuperscript{e} siècle, on attribuait les fièvres des marais à des émanations toxiques du sol (théorie miasmatique). C'était une \textbf{opinion scientifique de l'époque} (basée sur l'observation : fièvre + marais), mais fausse sur la cause.
En 1880, à l'hôpital militaire de Constantine (Algérie), le médecin militaire français \textbf{Charles Louis Alphonse Laveran} observe, au microscope, des corpuscules mobiles dans le sang d'un malade : le \emph{Plasmodium} (alors \emph{Oscillaria malariae}). Il reçoit le Prix Nobel en 1907.
En 1897, le britannique \textbf{Ronald Ross} (en Inde) démontre le cycle dans le moustique \emph{Anopheles}.
\emph{C'est le passage de la corrélation (marais/fièvre) à la causalité identifiée (parasite/vecteur) qui marque la naissance de la science du paludisme.}
\end{exemplebox}

\begin{savaistu}[Histoire des sciences : Laveran à Constantine]
Charles Laveran (1845-1922) était médecin militaire français. C'est dans le laboratoire de l'hôpital militaire de Constantine (Algérie), avec un microscope modeste, qu'il a découvert le parasite du paludisme dans le sang d'un caporal atteint de fièvre. Il a dû lutter contre le scepticisme de l'Académie de médecine (Pasteur le soutint). Cette découverte a transformé une maladie "mystérieuse" en une pathologie parasitaire identifiable, traçable, et donc combattable scientifiquement. Au Cameroun, le Centre Pasteur du Cameroun (CPC) et l'IRAD perpétuent cette tradition de recherche sur le terrain.
\end{savaistu}

\begin{popfigure}
\begin{tikzpicture}[
    node distance=1.2cm and 1.5cm,
    lab/.style={draw=popblue, thick, rounded corners, fill=popblueL, minimum height=1.5cm, text width=4.5cm, align=center, font=\small},
    trad/.style={draw=poporange, thick, rounded corners, fill=poporangeL, minimum height=1.5cm, text width=4.5cm, align=center, font=\small},
    arrow/.style={->, >=Stealth, thick, popdark},
    title/.style={font=\bfseries\large, text=popdark},
]
% Titre
\node[title] at (0,4.5) {Deux types de savoirs face au paludisme : Complémentarité ou Concurrence ?};

% Colonne Laboratoire IRAD
\node[lab, align=center] (lab1) at (-5,2){\textbf{Laboratoire IRAD / CPC / CHU (Nkolbisson, Yaoundé)}\\[4pt]
\textbf{Savoir Scientifique}\\
- Microscopie / PCR / Séquençage\\
- Identification \emph{Plasmodium} spp.\\
- Test de résistance (marqueurs moléculaires)\\
- Essais cliniques randomisés (double aveugle)\\
- Posologie standardisée (mg/kg)\\
- Pharmacovigilance (effets indésirables)\\
\textbf{Critère : Preuve, Reproductibilité, Universalité}};

% Colonne Marché Traditionnel
\node[trad, align=center] (trad1) at (5,2){\textbf{Marché Trad. / Tradipraticien (Mfoundi, Mokolo)}\\[4pt]
\textbf{Savoir Traditionnel / Expérientiel}\\
- Pharmacopée ancestrale (écorces, feuilles, racines)\\
- Diagnostic clinique / symbolique (pouls, regard, divination)\\
- Transmission orale (maître $\to$ disciple)\\
- Posologie empirique (verre, poignée, "jusqu'à guérison")\\
- Approche holistique (corps, esprit, ancêtres, social)\\
\textbf{Critère : Efficacité perçue, Cohérence culturelle, Accessibilité}};

% Flèches interaction
\draw[arrow, bend left=20] (lab1.east) to node[above, font=\small, text=poppurple] {\textbf{Valorisation :} Criblage phytochimique $\to$ Nouvelles molécules} (trad1.west);
\draw[arrow, bend right=20] (trad1.west) to node[below, font=\small, text=popgreen] {\textbf{Validation :} Essais cliniques $\to$ Médicament amélioré traditionnel (MAT)} (lab1.east);

% Bas : Patient
\node[draw=popgreen, thick, rounded corners, fill=popgreenL, minimum height=1.2cm, text width=10cm, align=center, font=\small] (patient) at (0,-1.5) {\textbf{LE PATIENT CAMEROUNAIS AU CŒUR DU SYSTÈME DE SANTÉ}\\[4pt]
Droit au meilleur soin possible (Déclaration d'Helsinki, Constitution OMS).\\
Nécessité d'une \textbf{médecine intégrée} : Diagnostic biologique certain (Science) + Prise en charge accessible, culturelle, observance (Tradition validée).\\
\textbf{Exemple :} TDR Palu positif $\to$ ACT (Science) \textbf{ET} Si rupture de stock / zone enclavée $\to$ MAT validé (ex. \emph{Artop} à base d'\emph{Artemisia annua} cultivée localement) sous surveillance.};

\draw[arrow] (lab1.south) -- (patient.north west);
\draw[arrow] (trad1.south) -- (patient.north east);
\end{tikzpicture}
\legende{Schéma comparatif : Laboratoire scientifique (IRAD/CPC) vs Savoir traditionnel (Marché/Tradipraticien). La flèche verte montre la convergence nécessaire vers le patient. Les flèches latérales illustrent la dialectique science/technique : la science valide la tradition (MAT), la tradition inspire la science (phytochimie).}
\end{popfigure}

\courssub{Synthèse et transition}

\emph{À retenir :} La science se définit comme un \textbf{savoir rationnel, objectif, universel, méthodique, vérifiable et cumulatif}. Elle se distingue radicalement de l'opinion (\emph{doxa}) par l'exigence de la preuve et de la reproductibilité. Elle se distingue de la technique par son but (le Vrai vs l'Utile), tout en entretenant avec elle une relation dialectique féconde (ex. cacao, paludisme). Enfin, la science n'est pas un dogme : elle est \textbf{révisable}, progressant par rectification d'obstacles épistémologiques (Bachelard) et falsification des théories (Popper).

La compréhension de cette nature de la science nous permet maintenant d'aborder la \textbf{section 2} : « Les différents types de sciences » (mathématiques, sciences de la matière, sciences humaines), car toutes ne partagent pas exactement la même méthode ni le même objet, malgré l'unité de l'esprit scientifique.

\coursec{Les différents types de sciences}

Dans la section précédente, nous avons défini la science par ses caractères essentiels : elle est un savoir rationnel, méthodique, objectif et vérifiable, qui se distingue de l'opinion et de la croyance. Mais dire « la science » au singulier est trompeur : un mathématicien qui démontre un théorème, un chimiste qui analyse un engrais au laboratoire de l'IRAD à Nkolbisson, et un sociologue qui enquête sur les tontines à Douala ne font pas exactement le même métier. Tous cherchent la vérité, mais ils ne regardent pas le même objet et n'utilisent pas la même méthode. Il faut donc maintenant \cle{classer} les sciences. Le problème qui guide cette section est le suivant : selon quels critères peut-on distinguer les sciences entre elles, et toutes les disciplines qui se prétendent scientifiques le sont-elles au même titre ?

\courssub{Le critère de classification : l'objet et la méthode}

Pour ranger les sciences, il ne suffit pas de regarder leur nom ou leur prestige. Deux critères décisifs permettent de les distinguer rigoureusement.

\begin{aretenir}[Les deux critères de classification des sciences]
Une science se définit par :
\begin{enumerate}
\item son \cle{objet} : ce qu'elle étudie (les nombres, la matière, le vivant, l'homme en société) ;
\item sa \cle{méthode} : la manière dont elle établit ses vérités (démonstration logique ou observation et expérimentation).
\end{enumerate}
De ces deux critères découle un troisième élément : le \cle{type de preuve} admis dans la discipline (preuve déductive ou preuve expérimentale).
\end{aretenir}

L'intuition est simple : on ne juge pas un poisson à sa capacité de grimper aux arbres. De même, on ne peut pas demander aux mathématiques de faire des expériences, ni à la physique de se passer du réel. Chaque type de science a son territoire et son outil propre. C'est pourquoi la classification classique distingue trois grandes familles : les \cle{sciences formelles}, les \cle{sciences de la matière} (ou sciences de la nature) et les \cle{sciences humaines}.

\begin{methode}[Comment classer une discipline en trois questions]
Face à une discipline donnée (par exemple la psychologie ou l'astronomie), procéder par étapes :
\begin{enumerate}
\item \textbf{Identifier l'objet} : de quoi parle cette science ? De formes abstraites, de la matière, ou de l'homme ?
\item \textbf{Identifier la méthode} : comment établit-elle ses résultats ? Par démonstration pure ou par observation et expérience ?
\item \textbf{Identifier le type de preuve} : qu'est-ce qui fait foi dans cette discipline ? La rigueur logique d'un raisonnement ou la reproduction d'une expérience ?
\end{enumerate}
La réponse à ces trois questions suffit à ranger la discipline dans la bonne famille et à éviter les confusions dans une copie de dissertation.
\end{methode}

\courssub{Les mathématiques : des sciences formelles et déductives}

Commençons par une expérience de pensée. Un élève calcule : « si j'emprunte \num{100000} FCFA à une tontine avec un intérêt de 10 \% par mois, je devrai \num{110000} FCFA le mois suivant ». Ce calcul est-il vrai parce qu'on l'a vérifié dans une tontine réelle à Bafoussam ? Non : il est vrai parce que les règles de l'arithmétique l'imposent. Même si la tontine fait faillite, même si personne ne rembourse, le calcul reste exact. Voilà l'intuition fondatrice : la vérité mathématique ne dépend pas de ce qui arrive dans le monde.

\begin{definition}[Science formelle]
Une \cle{science formelle} (mathématiques, logique) est une science qui étudie des formes de raisonnement et des relations entre symboles, indépendamment de tout contenu empirique. Ses vérités sont établies par \cle{démonstration}, c'est-à-dire par un enchaînement déductif de propositions à partir d'axiomes et de définitions.
\end{definition}

Les mathématiques sont donc des sciences \emph{déductives} : elles partent de principes posés au départ (les axiomes, par exemple « par un point extérieur à une droite, il passe une seule parallèle ») et en tirent des conséquences nécessaires. Leur grande force est la \cle{certitude} : un théorème démontré est vrai pour toujours et partout. Le théorème de Pythagore était vrai avant qu'on le découvre et restera vrai dans mille ans.

Mais cette certitude a un prix, souligné par Albert Einstein : « Pour autant que les propositions des mathématiques se rapportent à la réalité, elles ne sont pas certaines ; et pour autant qu'elles sont certaines, elles ne se rapportent pas à la réalité. » Autrement dit, les mathématiques ne parlent pas directement du réel : le nombre 3 n'existe nulle part dans la nature, on ne rencontre que trois chèvres, trois sacs de maïs, trois élèves. Les mathématiques décrivent des \emph{formes} que le réel peut incarner ou non.

\begin{exemplebox}[Les mathématiques au service de la vie quotidienne camerounaise]
Le calcul des intérêts dans une \textbf{tontine}, la répartition des parts dans un \textbf{njangui}, le calcul de la surface d'un champ de cacao, la facturation de l'ENEO ou de la CAMWATER : toutes ces opérations appliquent des vérités mathématiques certaines à des situations concrètes. Le calcul est infaillible ; c'est son application au réel (taux réellement pratiqué, mesures exactes du champ) qui peut comporter des erreurs.
\end{exemplebox}

\courssub{Les sciences de la matière : des sciences expérimentales}

Changeons maintenant de terrain. Un agronome veut savoir si un nouvel engrais augmente le rendement du maïs dans la plaine des Mbo. Peut-il le démontrer par un calcul ? Non : il doit semer, traiter une parcelle avec l'engrais, laisser une parcelle témoin sans traitement, récolter, peser et comparer. La vérité, ici, sort de la terre et non du seul raisonnement.

\begin{definition}[Science expérimentale]
Une \cle{science expérimentale} est une science qui fonde ses lois sur l'\emph{observation contrôlée} et l'\emph{expérimentation reproductible}. Elle formule des hypothèses sur le réel, les soumet à l'épreuve de l'expérience et n'accepte comme vrai que ce qui résiste à la vérification répétée.
\end{definition}

Les sciences de la matière — \cle{physique}, \cle{chimie}, \cle{biologie} — étudient la nature : la matière inanimée (physique et chimie) et la matière vivante (biologie). Leur méthode, que Claude Bernard a décrite dans son \emph{Introduction à l'étude de la médecine expérimentale} (1865), suit un parcours en quatre temps : observation d'un fait, formulation d'une hypothèse, vérification expérimentale, établissement d'une loi. Contrairement aux mathématiques, leurs vérités sont \emph{révisables} : une loi physique reste vraie jusqu'à ce qu'une expérience plus précise la mette en défaut. C'est ainsi que la mécanique de Newton a été complétée par celle d'Einstein.

\begin{exemplebox}[Les sciences de la matière au Cameroun]
\begin{itemize}
\item \textbf{Chimie des engrais} : les laboratoires de l'IRAD analysent les sols et les fertilisants pour conseiller les producteurs de cacao et de maïs ;
\item \textbf{Physique des télécommunications} : la propagation des ondes radio et la fibre optique expliquent le fonctionnement des réseaux MTN et Orange qui couvrent le pays ;
\item \textbf{Biologie médicale} : le diagnostic du paludisme au Centre Pasteur du Cameroun repose sur l'observation du parasite \emph{Plasmodium} au microscope.
\end{itemize}
\end{exemplebox}

Notons que ces sciences utilisent massivement les mathématiques comme \emph{instrument} : la physique exprime ses lois en équations. Mais la différence demeure : une équation physique peut être fausse si l'expérience la contredit, tandis qu'un théorème mathématique démontré ne peut pas être faux.

\courssub{Les sciences humaines : l'homme comme objet de science}

Reste le cas le plus délicat. Peut-on étudier l'homme comme on étudie une pierre ou une plante ? Les \cle{sciences humaines} — \cle{sociologie}, \cle{psychologie}, \cle{économie}, \cle{histoire}, anthropologie — prennent pour objet l'être humain : ses comportements, ses groupes, ses productions, son passé. Elles empruntent aux sciences de la nature le souci de la méthode : enquêtes, questionnaires, statistiques, comparaison de documents, entretiens.

Mais elles rencontrent une difficulté spécifique et radicale : \textbf{le savant fait partie de son objet}. Le physicien n'est pas une molécule, mais le sociologue camerounais qui étudie les conflits entre générations dans sa propre société est lui-même un homme, membre d'une génération, d'une ethnie, d'une classe sociale. Il a des opinions, des intérêts, des préjugés. Comment être sûr qu'il décrit la réalité et non ses propres désirs ? C'est le problème de l'\cle{objectivité} en sciences humaines.

De plus, l'objet humain se dérobe : un individu interrogé peut mentir, se trahir, modifier son comportement parce qu'il se sait observé ; et surtout, l'homme est libre et imprévisible, ce qui rend les lois humaines moins strictes que les lois physiques. On ne peut pas non plus expérimenter sur les hommes comme sur des éprouvettes : on ne déclenche pas une guerre ou une crise économique en laboratoire pour en étudier les causes.

\begin{attention}[Piège fréquent dans les copies]
N'écris jamais que les sciences humaines « ne sont pas des sciences ». C'est une thèse simpliste et mal notée. Écris plutôt qu'elles sont des sciences qui rencontrent des \textbf{difficultés spécifiques d'objectivité} : intrusion du chercheur dans son objet, impossibilité de l'expérimentation directe, complexité et variabilité des faits humains. Le débat porte sur le \emph{degré} de scientificité, non sur une exclusion pure et simple.
\end{attention}

\courssub{Problème : les sciences humaines sont-elles de vraies sciences ?}

Le débat est classique au Baccalauréat et oppose deux thèses.

\textbf{Thèse 1 : les sciences humaines ne sont pas de vraies sciences.} Pour les tenants d'une conception stricte de la science (héritée du positivisme le plus rigide), une vraie science doit être objective, expérimentale et capable de prévision exacte. Or les sciences humaines échouent sur ces trois points : le chercheur est engagé dans son objet, l'expérience est impossible, et nul ne peut prévoir les comportements humains comme on prévoit une éclipse. Max Weber reconnaît lui-même que le sociologue doit pratiquer une « neutralité axiologique » difficile, c'est-à-dire suspendre activement ses jugements de valeur — preuve que l'objectivité y est un effort permanent, jamais un acquis.

\textbf{Thèse 2 : les sciences humaines sont des sciences, mais d'un type particulier.} Émile Durkheim, dans \emph{Les Règles de la méthode sociologique} (1895), montre qu'on peut traiter les « faits sociaux comme des choses » : le taux de suicide, le taux de mariage, la délinquance présentent des régularités statistiques stables qu'on peut mesurer et expliquer objectivement. De même, l'économie élabore des modèles rigoureux, et l'histoire soumet ses documents à une critique méthodique sévère. La difficulté d'objectivité n'est pas une impossibilité : elle se surmonte par des procédures (échantillonnage, statistiques, croisement des sources, distanciation critique).

\begin{exemplebox}[Les mathématiques au service des sciences humaines : le cas de l'INS]
L'\textbf{Institut National de la Statistique (INS)} du Cameroun réalise des recensements de la population et calcule des indicateurs sociaux : taux de chômage, taux de scolarisation, indice des prix. Lors du troisième Recensement Général de la Population et de l'Habitat (RGPH), des dizaines de millions de questionnaires ont été traités par des méthodes statistiques rigoureuses (échantillonnage, moyennes, pourcentages). Ici, les mathématiques — science formelle et certaine — deviennent l'instrument qui confère aux sciences humaines une objectivité mesurable : on ne discute plus d'opinions, mais de chiffres vérifiables.
\end{exemplebox}

\textbf{Position synthétique.} Les sciences humaines sont des sciences en voie de consolidation : elles partagent avec les sciences de la nature l'exigence de méthode et de preuve, mais leur objet — l'homme libre et conscient — leur interdit la certitude expérimentale de la physique. Leur scientificité est réelle mais relative, et c'est précisément ce qui fait leur richesse : elles doivent sans cesse contrôler leurs propres préjugés.

\courssub{Complément pour le Bac : la classification d'Auguste Comte}

Auguste Comte, fondateur du positivisme, a proposé dans son \emph{Cours de philosophie positive} (1830-1842) une classification célèbre des sciences, ordonnée selon deux principes : de la plus \emph{simple} à la plus \emph{complexe}, et de la plus ancienne (la première à devenir positive) à la plus récente.

\begin{propriete}[La hiérarchie encyclopédique de Comte]
Comte ordonne les sciences fondamentales ainsi :
\[ \text{Mathématiques} \rightarrow \text{Astronomie} \rightarrow \text{Physique} \rightarrow \text{Chimie} \rightarrow \text{Biologie} \rightarrow \text{Sociologie} \]
Chaque science dépend de la précédente (la physique suppose les mathématiques, la biologie suppose la chimie) et étudie un objet plus complexe. La sociologie, science de l'homme en société, est la plus complexe et la dernière née : Comte l'appelle la « physique sociale ».
\end{propriete}

L'idée maîtresse est que la complexité croissante de l'objet explique la difficulté croissante de la science : il est plus facile d'établir des lois sur les nombres que sur les sociétés humaines. Cette classification éclaire directement notre problème : les sciences humaines sont au sommet de la hiérarchie parce que leur objet est le plus riche — et donc le plus rebelle à l'objectivité.

\begin{center}
\begin{popfigure}
\begin{tikzpicture}[
  node distance=0.55cm,
  boite/.style={draw=popblue, fill=popblueL, rounded corners=2pt, align=center, font=\small, text=popdark, minimum height=0.8cm},
  ex/.style={font=\scriptsize\itshape, text=popmuted, align=center},
  branche/.style={draw=popdark, thick}
]
\node[boite, fill=popgoldL, draw=popgold, font=\small\bfseries] (racine) {LA SCIENCE};
\node[boite, below left=1.1cm and 3.2cm of racine, align=center] (form){Sciences formelles\\ \textbf{Mathématiques, logique}\\ Méthode : démonstration};
\node[boite, below=1.1cm of racine, align=center] (mat){Sciences de la matière\\ \textbf{Physique, chimie, biologie}\\ Méthode : expérience};
\node[boite, below right=1.1cm and 3.2cm of racine, align=center] (hum){Sciences humaines\\ \textbf{Sociologie, psychologie,}\\ \textbf{économie, histoire}\\ Méthode : enquête, statistiques};
\draw[branche] (racine) -- (form);
\draw[branche] (racine) -- (mat);
\draw[branche] (racine) -- (hum);
\node[ex, below=0.15cm of form, align=center]{Calcul des intérêts\\ d'une tontine (njangui)};
\node[ex, below=0.15cm of mat, align=center]{Chimie des engrais (IRAD),\\ ondes des télécommunications};
\node[ex, below=0.15cm of hum, align=center]{Recensement de l'INS,\\ étude des migrations\\ vers Douala et Yaoundé};
\end{tikzpicture}
\legende{Arbre de classification des sciences selon l'objet et la méthode, avec des exemples camerounais pour chaque branche.}
\end{popfigure}
\end{center}

\begin{center}
\begin{popfigure}
\begin{tikzpicture}
\begin{axis}[
  width=0.85\linewidth, height=6.2cm,
  xlabel={Ordre d'apparition historique et de dépendance},
  ylabel={Complexité de l'objet},
  xmin=0, xmax=7, ymin=0, ymax=7,
  xtick={1,2,3,4,5,6},
  xticklabels={Math., Astron., Phys., Chim., Biol., Sociol.},
  x tick label style={font=\scriptsize, rotate=20, anchor=east},
  ytick=\empty,
  axis lines=left,
  every axis x label/.style={at={(ticklabel* cs:1)}, anchor=west, font=\scriptsize},
  every axis y label/.style={at={(ticklabel* cs:1)}, anchor=south, font=\scriptsize},
]
\addplot[popcurve, const plot, very thick] coordinates {(0,1) (1,2) (2,3) (3,4) (4,5) (5,6) (6.5,6)};
\addplot[only marks, mark=*, mark size=2.5pt, color=poporange] coordinates {(1,2) (2,3) (3,4) (4,5) (5,6)};
\end{axis}
\end{tikzpicture}
\legende{La hiérarchie de Comte en escalier : chaque marche correspond à une science dont l'objet est plus complexe que le précédent, des mathématiques (le plus simple) à la sociologie (le plus complexe).}
\end{popfigure}
\end{center}

\begin{exoresolu}[Classer une discipline : la médecine et l'économie]
Énoncé. Un sujet de dissertation demande : « Toutes les sciences se valent-elles ? » Pour préparer son introduction, un élève doit classer la médecine et l'économie selon leur objet, leur méthode et leur type de preuve. Rédiger ce classement en quelques lignes.
\tcblower
Solution détaillée. Appliquons la méthode en trois questions.
\begin{enumerate}
\item \textbf{La médecine} : son objet est le corps humain vivant, donc la matière vivante ; sa méthode repose sur l'observation clinique et l'expérimentation (essais thérapeutiques) ; son type de preuve est l'expérience reproductible. C'est une application de la \emph{biologie}, donc une science de la matière — à ceci près qu'elle vise aussi à soigner, ce qui lui donne une dimension pratique et éthique.
\item \textbf{L'économie} : son objet est l'activité humaine de production et d'échange (marchés de Mokolo, prix du cacao, politique budgétaire de l'État) ; sa méthode combine modèles mathématiques et enquêtes statistiques ; son type de preuve est mixte : démonstration formelle pour les modèles, vérification statistique pour les faits. C'est une \emph{science humaine} qui utilise les mathématiques comme instrument.
\end{enumerate}
Conclusion : toutes les sciences ne se valent pas au sens où elles n'ont ni le même objet, ni la même méthode, ni le même degré de certitude — mais chacune est légitime dans son domaine. Cette distinction prépare une problématique solide : la valeur d'une science dépend-elle de sa certitude ou de son utilité pour l'homme ?
\end{exoresolu}

\begin{exercice}[Pour s'entraîner]
1. En utilisant les deux critères de classification, range les disciplines suivantes : la logique, la géologie, la psychologie, l'astronomie, la démographie. Justifie chaque réponse en une phrase.
2. Sujet de réflexion : « Les mathématiques sont certaines mais vides ; les sciences humaines sont riches mais incertaines. » Discute cette affirmation en t'appuyant sur les notions de science formelle et de science expérimentale.
\end{exercice}

\begin{corrige}[Exercice n°1]
1. \textbf{La logique} : science formelle ; elle étudie les formes du raisonnement valide par démonstration, sans contenu empirique. \textbf{La géologie} : science de la matière ; elle étudie la structure de la Terre par observation et expérimentation (analyse d'échantillons). \textbf{La psychologie} : science humaine ; elle étudie les comportements et la vie mentale par l'observation et l'enquête, avec des difficultés d'objectivité. \textbf{L'astronomie} : science de la matière fondée sur l'observation (on ne peut pas expérimenter sur les astres, mais on mesure et on vérifie les prédictions). \textbf{La démographie} : science humaine qui utilise les statistiques, comme le montre le travail de l'INS au Cameroun.
2. L'affirmation est en partie juste : les mathématiques sont certaines parce qu'elles ne parlent pas du réel (sciences formelles), tandis que les sciences humaines parlent du réel le plus complexe — l'homme — au prix d'une objectivité difficile. Mais elle est excessive : « vides » est injuste, car les mathématiques s'appliquent puissamment au réel (tontines, télécommunications) ; et « incertaines » est relatif, car les méthodes statistiques donnent aux sciences humaines une fiabilité réelle, quoique moindre.
\end{corrige}

\begin{aretenir}[L'essentiel de la section]
\begin{itemize}
\item On classe les sciences selon leur \cle{objet} et leur \cle{méthode}, d'où trois familles : sciences formelles, sciences de la matière, sciences humaines.
\item Les \cle{mathématiques} sont déductives et certaines, mais ne parlent pas directement du réel.
\item Les \cle{sciences de la matière} (physique, chimie, biologie) sont expérimentales : leurs lois, tirées de l'observation et de l'expérience, sont révisables.
\item Les \cle{sciences humaines} étudient l'homme ; leur objectivité est difficile car le savant fait partie de son objet, mais les méthodes statistiques (INS, recensements) renforcent leur rigueur.
\item La classification de \cle{Comte} ordonne les sciences de la plus simple à la plus complexe : mathématiques, astronomie, physique, chimie, biologie, sociologie.
\end{itemize}
\end{aretenir}

Ayant distingué les types de sciences, une question nouvelle se pose : comment l'esprit humain est-il devenu capable de science ? Car penser scientifiquement ne va pas de soi : l'opinion, la tradition et l'imagination nous y résistent. C'est l'objet de la section suivante, consacrée à l'émergence de l'esprit scientifique et aux obstacles épistémologiques.

\coursec{L'émergence de l'esprit scientifique et les obstacles épistémologiques}

Nous avons vu dans la section précédente que la science n'est pas un bloc unique : les mathématiques, les sciences de la matière et les sciences humaines se distinguent par leurs objets et leurs méthodes. Mais une question demeure : comment l'esprit humain est-il parvenu à penser \emph{scientifiquement} ? L'homme n'est pas né savant. Pendant des millénaires, il a expliqué le monde par les mythes, les dieux, les esprits et l'autorité des anciens. La science est donc une \cle{conquête} historique, fragile et récente. Comprendre comment elle est née, et contre quoi elle a dû lutter, c'est comprendre ce que le philosophe Gaston Bachelard appelle la \cle{formation de l'esprit scientifique}.

\courssub{Les conditions historiques : la révolution scientifique des XVI\textsuperscript{e}-XVII\textsuperscript{e} siècles}

Commençons par une intuition simple. Aujourd'hui, tout écolier sait que la Terre tourne autour du Soleil. Pourtant, pendant près de deux mille ans, les plus grands savants ont affirmé le contraire : le Soleil tourne autour de la Terre. Et ils avaient de « bonnes raisons » : chaque matin, on \emph{voit} le Soleil se lever à l'est et se coucher à l'ouest ; personne ne \emph{sent} la Terre bouger. Comment a-t-on pu renverser une évidence aussi solide ?

C'est l'événement majeur de l'histoire des sciences : la \cle{révolution scientifique} des XVI\textsuperscript{e} et XVII\textsuperscript{e} siècles.

\begin{definition}[La révolution scientifique]
La \cle{révolution scientifique} désigne le bouleversement intellectuel qui, entre le milieu du XVI\textsuperscript{e} siècle et la fin du XVII\textsuperscript{e} siècle, substitue à la vision du monde héritée d'Aristote et de Ptolémée une conception nouvelle fondée sur l'observation, l'expérience et les mathématiques.
\end{definition}

Trois étapes essentielles marquent cette révolution :

\begin{itemize}
\item \textbf{Copernic (1543)}. Dans son ouvrage \emph{Des révolutions des sphères célestes}, l'astronome polonais Nicolas Copernic (1473-1543) propose le système \cle{héliocentrique} : c'est le Soleil, et non la Terre, qui occupe le centre du monde. La Terre devient une planète parmi d'autres. On passe ainsi du \cle{géocentrisme} (la Terre au centre) à l'\cle{héliocentrisme} (le Soleil au centre).
\item \textbf{Galilée (1610)}. Le savant italien Galilée (1564-1642) pointe vers le ciel une lunette astronomique qu'il a perfectionnée. Il observe les montagnes de la Lune, les satellites de Jupiter, les phases de Vénus : autant de faits que la cosmologie aristotélicienne ne peut pas expliquer. L'observation instrumentale \emph{confirme} Copernic.
\item \textbf{Newton (1687)}. Dans les \emph{Principia}, Isaac Newton (1642-1727) formule la loi de la gravitation universelle : une même loi mathématique explique la chute d'une pomme et le mouvement des planètes. Le ciel et la terre obéissent aux mêmes lois. La physique moderne est née.
\end{itemize}

Au XIX\textsuperscript{e} siècle, Louis Pasteur (1822-1895) prolonge ce mouvement dans le domaine du vivant : en démontrant le rôle des microbes dans les maladies, il fonde la microbiologie et montre que la méthode expérimentale s'applique aussi à la matière vivante.

\begin{popfigure}
\begin{center}
\begin{tikzpicture}[x=1cm,y=1cm]
% Frise chronologique
\draw[line width=1.2pt,popblue,-{Stealth[length=3mm]}] (0,0) -- (13.2,0);
\foreach \x/\annee in {1/1543, 4.2/1610, 7.4/1687, 11.5/XIX\textsuperscript{e} s.}{
  \draw[popblue,line width=1pt] (\x,-0.12) -- (\x,0.12);
  \node[below,popdark,font=\small\bfseries] at (\x,-0.15) {\annee};
}
% Jalons
\node[circle,fill=poporange,inner sep=2.2pt] (c) at (1,0) {};
\node[above,popdark,align=center,font=\small] at (1,0.35) {\textbf{Copernic}\\ héliocentrisme};
\node[circle,fill=popgreen,inner sep=2.2pt] at (4.2,0) {};
\node[above,popdark,align=center,font=\small] at (4.2,0.35) {\textbf{Galilée}\\ lunette, observation};
\node[circle,fill=poppurple,inner sep=2.2pt] at (7.4,0) {};
\node[above,popdark,align=center,font=\small] at (7.4,0.35) {\textbf{Newton}\\ gravitation universelle};
\node[circle,fill=poppink,inner sep=2.2pt] at (11.5,0) {};
\node[above,popdark,align=center,font=\small] at (11.5,0.35) {\textbf{Pasteur}\\ microbiologie};
\end{tikzpicture}
\end{center}
\legende{Frise chronologique de la révolution scientifique : de l'héliocentrisme de Copernic (1543) à la microbiologie de Pasteur (XIX\textsuperscript{e} siècle), en passant par les observations de Galilée (1610) et la synthèse de Newton (1687).}
\end{popfigure}

\begin{aretenir}[L'essentiel]
La science moderne naît d'une \cle{révolution} : le passage du géocentrisme à l'héliocentrisme symbolise le renversement de l'évidence sensible par le raisonnement et l'observation réglée. La science ne croit pas ce qu'elle voit immédiatement ; elle \emph{interroge} ce qu'elle voit.
\end{aretenir}

\courssub{Les conditions intellectuelles : le refus de l'argument d'autorité}

La révolution scientifique n'est pas seulement une affaire de dates et de découvertes. Elle repose sur un changement profond de \emph{mentalité} : le refus de l'\cle{argument d'autorité}.

\begin{definition}[Argument d'autorité]
L'\cle{argument d'autorité} consiste à tenir une proposition pour vraie non pas parce qu'elle a été vérifiée, mais parce qu'elle a été affirmée par une personne ou un texte prestigieux (un ancien, un maître, un livre sacré, la tradition).
\end{definition}

Pendant le Moyen Âge, la phrase « Aristote a dit » (\emph{ipse dixit}) suffisait à clore un débat. Or Aristote affirmait, par exemple, que les corps lourds tombent plus vite que les corps légers, proportionnellement à leur poids. Personne ne l'avait sérieusement vérifié. Galilée, lui, \emph{expérimente} : la tradition rapporte qu'il laisse tomber des objets de masses différentes du haut de la tour de Pise et constate qu'ils arrivent au sol en même temps (en négligeant la résistance de l'air). Contre l'autorité d'Aristote, Galilée oppose le verdict de l'\cle{expérience}.

C'est là le principe intellectuel fondateur de la science moderne : \textbf{la primauté de l'observation et de l'expérience sur l'autorité}. Une théorie, aussi respectable soit son auteur, doit être abandonnée si les faits la contredisent. Galilée résume cette attitude : la nature est un livre écrit en langage mathématique, et c'est ce livre qu'il faut lire directement, non les commentaires des anciens.

\begin{exemplebox}[Galilée contre Aristote]
Aristote affirme qu'un boulet de dix kilogrammes tombe dix fois plus vite qu'un boulet d'un kilogramme. Galilée objecte par un raisonnement : si l'on attache les deux boulets ensemble, le plus léger devrait ralentir le plus lourd ; mais l'ensemble, plus lourd, devrait tomber plus vite. Contradiction ! L'expérience tranche : les deux boulets tombent ensemble. L'autorité millénaire cède devant un fait.
\end{exemplebox}

\courssub{Les conditions sociales et techniques : instruments, imprimerie, institutions}

L'esprit scientifique n'aurait pu éclore sans des conditions matérielles et sociales favorables. Trois facteurs ont joué un rôle décisif.

\textbf{1. L'invention d'instruments.} La \cle{lunette astronomique} (perfectionnée par Galilée vers 1609) révèle un ciel inconnu ; le \cle{microscope} (XVII\textsuperscript{e} siècle, avec Leeuwenhoek) révèle un monde invisible : les « animalcules », ancêtres des microbes découverts par Pasteur. L'instrument prolonge les sens et corrige leurs limites : il permet d'\emph{observer ce que l'œil nu ne voit pas}. La science devient une observation \emph{armée}.

\textbf{2. L'imprimerie.} Inventée par Gutenberg vers 1450, l'imprimerie permet de diffuser rapidement et fidèlement les résultats scientifiques. L'ouvrage de Copernic circule dans toute l'Europe ; les découvertes deviennent vérifiables et discutables par tous. La science est un travail \emph{collectif} : elle suppose la communication publique des résultats.

\textbf{3. Les institutions scientifiques.} Naissent au XVII\textsuperscript{e} siècle les académies et sociétés savantes (Royal Society à Londres en 1660, Académie des sciences à Paris en 1666). Ces institutions organisent la confrontation des idées, la vérification des expériences et la publication des travaux. La vérité scientifique n'est plus l'affaire d'un homme seul, mais d'une \emph{communauté} qui contrôle et valide.

\begin{savaistu}[Gaston Bachelard, du bureau de poste à la Sorbonne]
Gaston Bachelard (1884-1962), le philosophe de la « rupture épistémologique », fut d'abord employé des postes à Remiremont, puis à Paris. Autodidacte passionné, il passe le baccalauréat sur le tard, devient professeur de physique et de chimie au collège de Bar-sur-Aube, puis obtient l'agrégation de philosophie et un doctorat. Il finit sa carrière comme professeur de philosophie des sciences à la Sorbonne. Son parcours illustre sa propre thèse : l'esprit scientifique se \emph{forme} par un effort permanent contre soi-même.
\end{savaistu}

\courssub{Bachelard et la rupture épistémologique}

Venons-en au cœur philosophique de la leçon. Pour Bachelard, dans \emph{La Formation de l'esprit scientifique} (1938), la science ne prolonge pas le savoir ordinaire : elle le \emph{contredit}. Entre l'opinion commune et la connaissance scientifique, il n'y a pas continuité, mais \cle{rupture}.

\begin{definition}[Rupture épistémologique]
La \cle{rupture épistémologique} est le moment où la science se sépare du sens commun, des images et de la philosophie spontanée pour fonder ses propres méthodes : observation réglée, expérimentation, mathématisation. Connaître scientifiquement, c'est d'abord \emph{rompre} avec ce que l'on croit savoir.
\end{definition}

Bachelard formule sa thèse en une phrase célèbre : « L'opinion a, en droit, toujours tort. » Pourquoi ? Parce que l'opinion \emph{pense mal} : elle ne pose pas de problèmes, elle se contente de répéter ce qui semble évident. Or « le fait scientifique est conquis, construit, constaté » : il n'est pas donné, il est le résultat d'un travail contre les évidences premières. Ainsi, l'évidence sensible dit que le Soleil tourne ; la science, après rupture, établit que c'est la Terre qui tourne.

\begin{attention}[Ce que Bachelard ne dit pas]
Bachelard ne dit pas que le sens commun est inutile ou méprisable. Dans la vie pratique, l'expérience ordinaire suffit souvent : savoir qu'il faut abriter la récolte avant la pluie ne demande pas un laboratoire. Le sens commun devient un \emph{obstacle} seulement quand il prétend \emph{expliquer} le monde à la place de la science. Le piège de la dissertation est de présenter Bachelard comme un contempteur du bon sens : il est le philosophe de la \emph{conquête} de l'objectivité, non de la négation de la vie quotidienne.
\end{attention}

\courssub{Les obstacles épistémologiques}

Si la science exige une rupture, c'est que notre esprit est encombré de ce que Bachelard appelle des \cle{obstacles épistémologiques}.

\begin{definition}[Obstacle épistémologique]
Un \cle{obstacle épistémologique} est une habitude de pensée, un préjugé ou une image qui s'oppose à la connaissance scientifique et ralentit sa formation. Il ne s'agit pas d'une simple ignorance (qu'on peut combler), mais d'un \emph{faux savoir} actif, qui résiste et se défend.
\end{definition}

Bachelard en décrit plusieurs. Retenons les principaux :

\begin{itemize}
\item \textbf{L'expérience première} : c'est l'obstacle des apparences sensibles. Ce que nous voyons d'abord nous trompe, et nous prenons cette première impression pour la vérité.
\item \textbf{L'obstacle verbal} : les mots vagues et les métaphores donnent l'illusion d'expliquer. Dire qu'une plante « aspire » l'eau ou qu'un corps « craint » le vide, c'est croire comprendre parce qu'on a nommé.
\item \textbf{L'obstacle animiste} : c'est la tendance à prêter une âme, des intentions, des désirs aux choses et aux phénomènes naturels (expliquer la maladie par un esprit, la foudre par une colère).
\item \textbf{L'obstacle substantialiste} : on croit qu'il existe des « substances » mystérieuses dotées de propriétés magiques (le « phlogistique » censé expliquer la combustion, avant la découverte de l'oxygène par Lavoisier).
\item \textbf{L'obstacle unitariste} : le désir de ramener tous les phénomènes à un principe unique, quitte à forcer les faits.
\item \textbf{L'obstacle généralisateur} : la tendance à étendre abusivement à tous les cas une loi vraie dans un domaine limité.
\end{itemize}

\begin{popfigure}
\begin{center}
\begin{tikzpicture}
\node[font=\small, align=center]{
\begin{tabular}{|p{3cm}|p{5.2cm}|p{5.2cm}|}
\hline
\rowcolor{popblueL} \textbf{Obstacle} & \textbf{Définition brève} & \textbf{Exemple concret} \\
\hline
Expérience première & Les apparences sensibles prises pour la réalité & « La terre est plate parce qu'on la voit plate » \\
\hline
Verbal & Les mots vagues qui donnent l'illusion d'expliquer & « La nature a horreur du vide » \\
\hline
Animiste & Prêter des intentions aux phénomènes naturels & « L'épidémie est l'œuvre d'un sorcier » \\
\hline
Substantialiste & Croire à des substances aux vertus mystérieuses & Le « phlogistique » censé expliquer le feu \\
\hline
Unitariste & Tout ramener à un principe unique & Tout expliquer par « les énergies » \\
\hline
Généralisateur & Étendre abusivement une loi à tous les cas & « Ce remède a guéri mon voisin, il guérit tout le monde » \\
\hline
\end{tabular}
};
\end{tikzpicture}
\end{center}
\legende{Les principaux obstacles épistémologiques selon Gaston Bachelard (\emph{La Formation de l'esprit scientifique}, 1938) : nom, définition et exemple.}
\end{popfigure}

\begin{exemplebox}[L'expérience première : la terre plate]
Croire que « la terre est plate parce qu'on la voit plate » est l'exemple type de l'obstacle de l'\cle{expérience première}. L'œil ne perçoit qu'une petite portion de la surface terrestre ; la courbure du globe échappe à la perception immédiate. Seul un raisonnement dépassant l'apparence (l'ombre de la Terre sur la Lune, les navires dont on voit d'abord le mât à l'horizon) permet de conclure que la Terre est ronde. La science commence donc par une \emph{méfiance} à l'égard du témoignage immédiat des sens.
\end{exemplebox}

\courssub{Application au contexte camerounais : sorcellerie ou microbes ?}

Ces analyses ne sont pas de l'histoire ancienne réservée à l'Europe. Elles éclairent directement notre réalité. Considérons une épidémie dans un village camerounais — par exemple une épidémie de choléra ou de méningite.

\textbf{L'explication traditionnelle} relève souvent de l'obstacle \cle{animiste} : on attribue la maladie à la sorcellerie, à un sort jeté par un ennemi, à la colère des ancêtres. Cette explication a une fonction sociale (elle désigne un coupable, elle donne un sens au malheur), mais elle ne permet ni de prévenir ni de soigner la maladie.

\textbf{L'explication scientifique}, elle, procède autrement : observation des cas, prélèvements, analyse au laboratoire, identification du microbe (le vibrion cholérique, le méningocoque), recherche de la source de contamination (eau souillée, promiscuité), puis mesures efficaces : traitement de l'eau, vaccination, antibiotiques. C'est exactement la démarche que Pasteur a inaugurée au XIX\textsuperscript{e} siècle.

Le contraste est frappant : là où l'explication animiste conduit à la chasse aux sorciers, l'explication microbienne conduit au chlore dans les puits et aux campagnes de vaccination — comme celles menées au Cameroun contre la méningite dans le Grand Nord ou contre la poliomyélite. La rupture épistémologique n'est pas un luxe de philosophe : elle sauve des vies.

\begin{attention}[Nuancer sans relativiser]
Il ne s'agit pas de mépriser les cultures traditionnelles : la pharmacopée traditionnelle contient des savoirs empiriques précieux (certains remèdes de plantes ont inspiré des médicaments). Mais un savoir devient scientifique seulement lorsqu'il accepte d'être \emph{vérifié expérimentalement} et \emph{corrigé} par les faits. L'obstacle commence là où une croyance se refuse à tout contrôle.
\end{attention}

C'est ici que l'\cle{école} joue son rôle décisif. En enseignant la démarche expérimentale — observer, émettre une hypothèse, vérifier, conclure —, elle forme l'esprit scientifique, c'est-à-dire un esprit capable de distinguer une croyance d'une connaissance, une rumeur d'un fait établi. L'élève qui apprend à demander « quelles sont les preuves ? » devient un citoyen moins vulnérable aux charlatans, aux fausses informations et aux rumeurs qui, on l'a vu lors des campagnes de vaccination, peuvent coûter cher à la santé publique.

\courssub{La psychanalyse de la connaissance objective : un travail permanent}

Dernière étape de la pensée de Bachelard : pourquoi parle-t-il de \cle{psychanalyse de la connaissance objective} ? Parce que les obstacles épistémologiques ne sont pas seulement des erreurs intellectuelles : ils ont des racines \emph{affectives}. Nous \emph{aimons} nos premières images, nos explications rassurantes, nos certitudes héritées. Y renoncer coûte, comme renoncer à une illusion en psychanalyse.

Former son esprit scientifiquement est donc un \emph{travail permanent} contre soi-même : il faut sans cesse surveiller ses préjugés, accepter d'avoir eu tort, préférer un fait déplaisant à une belle théorie. Bachelard écrit que l'esprit scientifique doit se former « contre la nature », contre nos penchants spontanés. Le savant n'est jamais définitivement guéri de ses obstacles : il les combat toute sa vie. C'est pourquoi la science progresse par \emph{erreurs rectifiées} : « la vérité scientifique est toujours une erreur rectifiée ».

\begin{methode}[Former son esprit scientifique : la démarche bachelardienne]
\begin{enumerate}
\item \textbf{Identifier} ses évidences premières : « qu'est-ce que je tiens pour vrai sans l'avoir vérifié ? »
\item \textbf{Suspecter} les apparences, les mots vagues et les explications par les intentions (animisme).
\item \textbf{Rompre} avec l'opinion : formuler une hypothèse claire et précise.
\item \textbf{Vérifier} par l'observation réglée et l'expérience.
\item \textbf{Rectifier} ses erreurs et recommencer : la connaissance objective est un progrès sans fin.
\end{enumerate}
\end{methode}

\begin{exoresolu}[Analyse d'une situation]
\textbf{Énoncé.} Dans un quartier de Douala, plusieurs cas de paludisme grave surviennent après de fortes pluies. Certains habitants affirment : « Cette maladie vient d'un mauvais sort jeté au quartier. » Un élève de Terminale répond : « Le paludisme est transmis par les moustiques qui prolifèrent dans les eaux stagnantes ; il faut assécher les flaques et dormir sous moustiquaire. » Identifiez l'obstacle épistémologique présent dans la première affirmation et montrez en quoi la seconde relève de l'esprit scientifique.

\tcblower
\textbf{Solution détaillée.} La première affirmation relève de l'\cle{obstacle animiste} : elle prête une intention (un « mauvais sort ») à un phénomène naturel, au lieu d'en chercher la cause matérielle. Cet obstacle est dangereux car il détourne des mesures efficaces : on cherchera un coupable au lieu d'éliminer les gîtes larvaires. La seconde affirmation illustre l'esprit scientifique : elle repose sur une \emph{observation} (les flaques d'eau après la pluie), sur une \emph{connaissance établie} (le rôle du moustique anophèle, découvert à la fin du XIX\textsuperscript{e} siècle) et elle aboutit à des \emph{conséquences pratiques vérifiables} (assèchement, moustiquaires imprégnées). Elle illustre la \cle{rupture épistémologique} : on a cessé d'expliquer par des intentions pour expliquer par des causes, ce qui seul permet d'agir efficacement.
\end{exoresolu}

\begin{aretenir}[À retenir pour la dissertation]
\begin{itemize}
\item La science naît d'une \cle{révolution} (Copernic, Galilée, Newton) qui substitue l'expérience et les mathématiques à l'autorité.
\item Pour Bachelard, connaître scientifiquement exige une \cle{rupture épistémologique} avec le sens commun.
\item L'esprit scientifique se forme \emph{contre} les \cle{obstacles épistémologiques} : expérience première, obstacle verbal, animiste, substantialiste, unitariste, généralisateur.
\item Cette formation est un travail permanent de rectification de soi : « la vérité est une erreur rectifiée ».
\end{itemize}
\end{aretenir}

Ainsi, l'esprit scientifique n'est pas un don, mais une conquête permanente contre les préjugés. Mais une fois formé, que vaut-il ? Peut-on tout attendre de la science, et ses applications sont-elles toujours bénéfiques ? C'est la question de la \emph{valeur} de la science, que nous aborderons dans les sections suivantes.

\coursec{La méthode expérimentale et ses applications}

Nous avons vu, dans la section précédente, comment l'esprit scientifique naît en rompant avec les obstacles épistémologiques : les opinions, les préjugés, les fausses évidences. Mais rompre avec l'opinion ne suffit pas. Encore faut-il savoir \emph{comment} interroger la réalité de manière rigoureuse. C'est précisément le rôle de la \cle{méthode expérimentale}. Dire qu'une affirmation est « scientifique », ce n'est pas dire qu'elle est prononcée par un savant en blouse blanche : c'est dire qu'elle a été obtenue et vérifiée selon une démarche précise, contrôlée, que n'importe quel autre chercheur peut reproduire. Cette section présente les étapes de cette démarche, le rôle central de l'hypothèse, l'œuvre fondatrice de Claude Bernard pour l'expérimentation sur le vivant, puis deux applications concrètes au Cameroun : l'essai clinique en médecine et les parcelles expérimentales en agronomie. Nous terminerons par les limites de l'expérimentation sur le vivant.

\courssub{Les étapes de la méthode expérimentale}

Commençons par une intuition simple. Un élève de Terminale technique remarque que ses plants de tomates, arrosés avec une eau contenant de l'engrais, poussent plus vite que ceux arrosés à l'eau simple. A-t-il le droit de conclure que l'engrais fait pousser les tomates ? Pas encore : peut-être les plants arrosés à l'engrais étaient-ils mieux exposés au soleil, ou dans une terre plus riche. Pour conclure rigoureusement, il faut \emph{contrôler} toutes les autres conditions et ne faire varier qu'un seul facteur. Toute la méthode expérimentale tient dans cette exigence.

\begin{definition}[La méthode expérimentale]
La \cle{méthode expérimentale} est la démarche rationnelle par laquelle le savant interroge la nature en soumettant une hypothèse explicative à l'épreuve contrôlée de l'expérience, afin de la valider ou de la réfuter. Elle se déploie en cinq étapes : \textbf{observation du problème}, \textbf{formulation d'une hypothèse}, \textbf{expérimentation} (avec variables contrôlées), \textbf{analyse des résultats}, \textbf{conclusion} (validation ou réfutation).
\end{definition}

Détaillons chaque étape.

\begin{enumerate}
\item \textbf{L'observation du problème.} Toute recherche part d'un fait étonnant, d'une question : pourquoi certains malades guérissent-ils et d'autres non ? Pourquoi telle variété de maïs résiste-t-elle à la sécheresse ? L'observation doit être précise, mesurée, débarrassée des préjugés (nous l'avons vu avec les obstacles épistémologiques).
\item \textbf{La formulation d'une hypothèse.} Le savant propose une explication provisoire : « c'est le paludisme qui cause cette fièvre », « c'est l'engrais qui accélère la croissance ». L'hypothèse n'est pas une vérité : c'est une conjecture à tester.
\item \textbf{L'expérimentation.} On construit un dispositif où l'on fait varier \emph{un seul} facteur (la \cle{variable indépendante}) tout en maintenant constantes toutes les autres conditions (les \cle{variables contrôlées}), et l'on mesure l'effet produit (la \cle{variable dépendante}). C'est le principe du « toutes choses égales par ailleurs ».
\item \textbf{L'analyse des résultats.} Les données recueillies (mesures, comptages, courbes) sont dépouillées, souvent à l'aide des mathématiques et des statistiques.
\item \textbf{La conclusion.} Deux cas se présentent : si les résultats confirment les prévisions de l'hypothèse, celle-ci est \emph{validée provisoirement} ; s'ils les contredisent, elle est \emph{réfutée}, et il faut en formuler une nouvelle. La démarche est donc une \textbf{boucle} : la science avance par essais et erreurs.
\end{enumerate}

\begin{popfigure}
\begin{center}
\begin{tikzpicture}[
node distance=0.55cm and 1.1cm,
etape/.style={rectangle, rounded corners=3pt, draw=popblue, fill=popblueL, text width=4.6cm, align=center, font=\small, inner sep=5pt},
decision/.style={rectangle, rounded corners=3pt, draw=poporange, fill=poporangeL, text width=3.6cm, align=center, font=\small, inner sep=5pt},
fleche/.style={-{Stealth[length=2.5mm]}, thick, popdark}]
\node[etape, align=center] (obs){\textbf{1. Observation}\\ repérage d'un problème};
\node[etape, below=of obs, align=center] (hyp){\textbf{2. Hypothèse}\\ explication provisoire testable};
\node[etape, below=of hyp, align=center] (exp){\textbf{3. Expérimentation}\\ variables contrôlées};
\node[etape, below=of exp, align=center] (res){\textbf{4. Résultats}\\ mesures et analyse};
\node[decision, below=of res, align=center] (concl){\textbf{5. Conclusion}\\ hypothèse confirmée ?};
\node[etape, right=2.2cm of concl, fill=popgreenL, draw=popgreen, align=center] (valid){Validation provisoire\\ (loi ou théorie)};
\draw[fleche] (obs) -- (hyp);
\draw[fleche] (hyp) -- (exp);
\draw[fleche] (exp) -- (res);
\draw[fleche] (res) -- (concl);
\draw[fleche] (concl) -- node[above, font=\small]{oui} (valid);
\draw[fleche, poppink] (concl.west) -- ++(-1.6,0) node[midway, above, font=\small]{non} |- (hyp.west);
\node[font=\small\itshape, poppink, rotate=90] at ($(concl.west)+(-1.85,2.6)$) {hypothèse réfutée : on recommence};
\end{tikzpicture}
\end{center}
\legende{Organigramme des cinq étapes de la méthode expérimentale. La boucle rose montre que la réfutation de l'hypothèse n'est pas un échec : elle oblige à formuler une hypothèse nouvelle, et fait ainsi progresser la connaissance.}
\end{popfigure}

\begin{methode}[Analyser un protocole expérimental]
Face à un protocole expérimental (sujet de dissertation, commentaire de texte ou exercice), procédez en trois étapes :
\begin{enumerate}
\item Repérez la \cle{variable indépendante} : c'est le facteur que l'expérimentateur modifie volontairement (ex. : administrer ou non le médicament).
\item Repérez la \cle{variable dépendante} : c'est ce que l'on mesure et qui dépend de la manipulation (ex. : le taux de guérison).
\item Repérez les \cle{variables contrôlées} : ce sont tous les facteurs maintenus identiques entre les groupes (âge des patients, dose, durée, conditions de vie), afin que la différence observée ne puisse être attribuée qu'à la variable indépendante.
\end{enumerate}
Un protocole où plusieurs facteurs varient en même temps est \emph{invalide} : on ne peut plus savoir lequel a produit l'effet.
\end{methode}

\courssub{Le rôle de l'hypothèse : testable et réfutable}

L'hypothèse est le moteur de la recherche. Sans elle, l'expérimentateur ne saurait ni quoi observer, ni quoi mesurer : comme le dit Claude Bernard, « l'hypothèse est le principe de toute recherche et de toute invention ». Mais toutes les affirmations ne sont pas des hypothèses scientifiques.

\begin{definition}[Hypothèse scientifique]
Une \cle{hypothèse scientifique} est une proposition explicative \textbf{provisoire} qui doit pouvoir être \textbf{mise à l'épreuve de l'expérience} : elle est \emph{testable} (on peut concevoir une expérience pour la vérifier) et \emph{réfutable} (on peut concevoir un résultat qui la contredirait).
\end{definition}

Le philosophe des sciences Karl \textbf{Popper} (\emph{La Logique de la découverte scientifique}, 1934) a tiré de cette exigence un critère célèbre de \cle{réfutabilité} (ou falsifiabilité) : une théorie n'est scientifique que si elle accepte le risque d'être contredite par les faits. Dire « tous les cygnes sont blancs » est scientifique, car il suffit d'un cygne noir pour réfuter l'affirmation. En revanche, dire « les esprits influencent nos destins, mais de manière indétectable » n'est pas une hypothèse scientifique : aucune expérience ne pourrait jamais la contredire. Une théorie qui explique tout et son contraire n'explique rien.

\begin{attention}[Corrélation n'est pas causalité]
Erreur classique du commentaire de données : conclure d'une \textbf{corrélation} (deux phénomènes varient ensemble) à une \textbf{causalité} (l'un produit l'autre). Exemple : on observe que les villages où l'on vend le plus de moustiquaires sont aussi ceux où l'on recense le plus de cas de paludisme déclarés. Faut-il conclure que les moustiquaires donnent le paludisme ? Non : c'est parce que le paludisme y est fréquent que l'on y vend des moustiquaires. Seule une expérimentation contrôlée (variables contrôlées, groupe témoin) permet d'établir une cause. Retenez la formule : \emph{la corrélation n'est pas la causalité}.
\end{attention}

\courssub{Claude Bernard et l'expérimentation sur la matière vivante}

Peut-on expérimenter sur le vivant comme on expérimente sur la matière inerte ? Au XIX\textsuperscript{e} siècle, beaucoup soutenaient que non : pour les « vitalistes », les êtres vivants seraient gouvernés par un principe mystérieux, la « force vitale », qui échapperait à toute loi. La médecine restait alors un art fondé sur l'habitude et l'autorité des anciens, non une science.

Claude \textbf{Bernard} (1813-1878), dans son \emph{Introduction à l'étude de la médecine expérimentale} (1865), renverse cette position : \textbf{le vivant obéit à des lois déterminées}, aussi rigoureuses que celles de la physique et de la chimie. Les phénomènes vitaux — digestion, circulation, sécrétion — ont des causes précises que l'expérience peut découvrir. La méthode est la même pour la matière brute et pour la matière vivante : observer, émettre une hypothèse, expérimenter en variant les conditions, conclure. Bernard distingue d'ailleurs le médecin \emph{observateur}, qui constate les faits, du médecin \emph{expérimentateur}, qui les \emph{provoque} volontairement pour en isoler les causes : « l'expérimentateur est celui qui, à l'aide de l'expérience, fait parler la nature ».

\begin{savaistu}[Le foie et le sucre : une découverte fondatrice]
En étudiant des chiens nourris de différentes façons, Claude Bernard remarque que le sang contient du sucre même lorsque l'animal n'en mange pas. Hypothèse : un organe \emph{fabrique} le sucre. En perfusant le foie, il montre que cet organe produit et libère du glucose dans le sang : c'est la découverte de la \textbf{fonction glycogénique du foie} (1857). Cette expérience rigoureuse fonde la médecine expérimentale moderne et prépare la compréhension de maladies comme le diabète. Elle illustre parfaitement la thèse de Bernard : le vivant obéit à des mécanismes déterminés que l'expérience peut mettre en lumière.
\end{savaistu}

\courssub{Application médicale : l'essai clinique}

L'application la plus directe de la méthode expérimentale au vivant est l'\cle{essai clinique}, qui permet de savoir si un traitement est réellement efficace. Son principe applique exactement la logique des variables contrôlées :

\begin{itemize}
\item On constitue \textbf{deux groupes} de patients comparables (même maladie, profils semblables) : le \textbf{groupe traité}, qui reçoit le médicament, et le \textbf{groupe témoin} (ou groupe de contrôle), qui ne le reçoit pas — ou reçoit un placebo ou le traitement habituel.
\item La \textbf{variable indépendante} est l'administration du traitement ; la \textbf{variable dépendante} est l'état de santé mesuré (taux de guérison, disparition du parasite) ; les \textbf{variables contrôlées} sont l'âge, la gravité initiale, la durée du suivi, etc.
\item Si le groupe traité guérit nettement plus que le groupe témoin, l'écart est attribué au traitement — à condition d'être confirmé par l'analyse statistique.
\end{itemize}

C'est ainsi que sont testés, y compris au Cameroun, les vaccins et les traitements antipaludiques : avant d'être recommandés, les nouveaux médicaments contre le paludisme et les candidats vaccins (comme le vaccin RTS,S/AS01 \emph{Mosquirix}, introduit depuis 2024 dans le programme élargi de vaccination du Cameroun) ont été évalués par des essais cliniques comparant des groupes vaccinés à des groupes témoins.

\begin{exemplebox}[Un essai chiffré]
Un essai compare 100 patients atteints de paludisme simple, traités par un nouvel antipaludique, à 100 patients témoins recevant le traitement habituel de référence insuffisant. Résultats : \textbf{85 \% de guérison} dans le groupe traité contre \textbf{20 \%} dans le groupe témoin. L'écart de 65 points \emph{suggère fortement} l'efficacité du traitement. Mais le chercheur prudent ne conclut pas immédiatement : il vérifie que les deux groupes étaient bien comparables (variables contrôlées) et que l'écart est \textbf{statistiquement significatif}, c'est-à-dire trop grand pour être dû au seul hasard. Ce n'est qu'ensuite que l'hypothèse « ce médicament est efficace » est validée provisoirement.
\end{exemplebox}

\begin{popfigure}
\begin{center}
\begin{tikzpicture}
\begin{axis}[
width=0.85\linewidth, height=6.2cm,
xlabel={Temps de suivi (jours)},
ylabel={Taux de guérison (\%)},
xmin=0, xmax=14, ymin=0, ymax=100,
xtick={0,2,4,6,8,10,12,14},
ytick={0,20,40,60,80,100},
grid=both, grid style={popline!40},
legend style={at={(0.03,0.97)}, anchor=north west, font=\small},
]
\addplot[popcurve, popblue, mark=*, mark size=1.5pt] coordinates {(0,0) (2,18) (4,42) (6,63) (8,76) (10,83) (12,85) (14,85)};
\addlegendentry{Groupe traité (nouvel antipaludique)}
\addplot[popcurve, poppink, mark=square*, mark size=1.5pt, dashed] coordinates {(0,0) (2,4) (4,8) (6,12) (8,16) (10,18) (12,20) (14,20)};
\addlegendentry{Groupe témoin (sans nouveau traitement)}
\end{axis}
\end{tikzpicture}
\end{center}
\legende{Courbe fictive d'un essai clinique : évolution du taux de guérison sur 14 jours. L'écart croissant entre les deux courbes (85 \% contre 20 \% au jour 14) suggère l'efficacité du traitement, à confirmer par l'analyse statistique. Données inventées à visée pédagogique.}
\end{popfigure}

\begin{exoresolu}[Analyser un protocole d'essai clinique]
Énoncé. Des chercheurs veulent tester l'efficacité d'un sirop à base de plantes contre la toux. Ils le donnent à 50 élèves d'un lycée de Douala qui toussent, et comparent leur guérison à celle de 50 élèves d'un autre lycée, non traités. Après une semaine, 70 \% des élèves traités ne toussent plus, contre 55 \% des autres. Les chercheurs concluent : « le sirop guérit la toux ». Cette conclusion est-elle rigoureuse ?
\tcblower
Solution détaillée. Appliquons la méthode d'analyse d'un protocole.
\begin{enumerate}
\item \emph{Variable indépendante} : l'administration du sirop. \emph{Variable dépendante} : la disparition de la toux après une semaine.
\item \emph{Variables contrôlées} : c'est ici que le protocole est défaillant. Les deux groupes appartiennent à deux lycées différents : ils peuvent différer par l'environnement (poussière, humidité), l'alimentation, l'accès aux soins, le type de toux. Plusieurs facteurs varient en même temps que le sirop : on ne peut donc pas attribuer l'écart au seul sirop.
\item \emph{Interprétation} : l'écart (70 \% contre 55 \%) est modeste ; il peut s'expliquer par ces différences de contexte, ou par la guérison spontanée (la plupart des toux guérissent seules en une semaine), ou par un effet placebo (croire qu'on est soigné améliore parfois l'état ressenti).
\end{enumerate}
\emph{Conclusion} : la conclusion des chercheurs n'est pas rigoureuse. Un bon protocole exigerait des groupes constitués par tirage au sort au sein du même milieu, un groupe témoin recevant un placebo, et une analyse statistique de l'écart. Cet exercice illustre la vigilance épistémologique : une expérience mal contrôlée produit l'illusion d'une preuve.
\end{exoresolu}

\courssub{Application agricole : les parcelles expérimentales}

La même logique s'applique aux plantes. Au Cameroun, l'\textbf{IRAD} (Institut de Recherche Agricole pour le Développement) conduit des expérimentations sur des \cle{parcelles expérimentales} pour améliorer les cultures vivrières et de rente : maïs, manioc, cacao. Le protocole est le suivant : on divise un champ homogène en plusieurs parcelles identiques (même sol, même exposition, même date de semis, même densité — autant de \textbf{variables contrôlées}) ; on fait varier un seul facteur, par exemple la variété semée ou la dose d'engrais (\textbf{variable indépendante}) ; on mesure ensuite le rendement à l'hectare ou la résistance aux maladies (\textbf{variable dépendante}).

C'est grâce à cette méthode que sont sélectionnées des variétés de maïs à haut rendement adaptées aux sols de l'Adamaoua et de l'Ouest, des boutures de manioc résistant à la mosaïque, ou des cacaoyers tolérants à la pourriture brune des cabosses. L'agriculteur camerounais qui adopte une « variété améliorée » bénéficie, souvent sans le savoir, d'années d'expérimentation rigoureuse. La science n'est donc pas une affaire de laboratoires lointains : elle est à l'œuvre dans les champs de Ngaoundéré, de Bafoussam ou d'Ambam.

\begin{exercice}[À vous de concevoir le protocole]
Un agronome de l'IRAD veut vérifier si une nouvelle variété de manioc résiste mieux à la mosaïque que la variété locale. Décrivez le protocole expérimental complet : variable indépendante, variable dépendante, variables contrôlées, nombre de parcelles, critère de conclusion.
\end{exercice}

\begin{corrige}[Exercice : parcelles de manioc]
\emph{Variable indépendante} : la variété plantée (nouvelle ou locale). \emph{Variable dépendante} : la proportion de plants atteints de mosaïque après une durée fixée (et éventuellement le rendement en tubercules). \emph{Variables contrôlées} : même sol, même champ ou champs voisins, même date de plantation, même densité, mêmes apports d'eau et d'entretien, même durée d'observation. \emph{Dispositif} : au moins plusieurs parcelles par variété (répétitions) pour écarter l'effet du hasard, réparties aléatoirement dans le champ. \emph{Critère de conclusion} : si la nouvelle variété présente une proportion de plants malades nettement inférieure, de façon statistiquement significative et sur plusieurs répétitions, l'hypothèse de sa meilleure résistance est validée provisoirement ; sinon, elle est réfutée et l'on teste une autre variété (boucle de la méthode expérimentale).
\end{corrige}

\courssub{Les limites de l'expérimentation sur le vivant}

La puissance de la méthode expérimentale ne doit pas faire oublier ses limites, précisément lorsqu'elle s'applique à la matière vivante.

\textbf{Première limite : la complexité du vivant.} Un organisme n'est pas une machine simple : les variables y sont innombrables et interagissent (patrimoine génétique, alimentation, stress, environnement). Isoler \emph{une seule} cause, comme l'exige la méthode, est souvent difficile. De plus, ce qui est vrai pour l'animal de laboratoire ne l'est pas toujours pour l'homme : un médicament efficace sur le rat peut être inactif ou toxique chez le patient. Enfin, l'expérimentation \emph{simplifie} le réel pour le comprendre : il faut ensuite vérifier que les conclusions valent dans les conditions réelles, plus complexes.

\textbf{Deuxième limite : les problèmes éthiques.} Expérimenter sur le vivant, c'est intervenir sur des êtres sensibles, parfois sur des personnes humaines. L'expérimentation animale pose la question de la souffrance infligée ; l'expérimentation humaine pose celle du \cle{consentement éclairé} : nul ne peut être soumis à un essai sans être informé des risques et sans accepter librement. L'histoire comporte des abus tragiques (expériences menées sur des populations sans leur consentement), qui ont conduit à des règles internationales strictes et à des comités d'éthique. La science expérimentale doit donc être encadrée : tout ce qui est techniquement possible n'est pas moralement permis. Cette question sera approfondie dans la dernière section, consacrée à la bioéthique.

\begin{aretenir}[L'essentiel de la section]
\begin{itemize}
\item La méthode expérimentale comporte cinq étapes en boucle : observation, hypothèse, expérimentation à variables contrôlées, résultats, conclusion (validation ou réfutation).
\item Une hypothèse scientifique est une explication provisoire, testable et réfutable (critère de Popper) ; une corrélation ne prouve pas une causalité.
\item Claude Bernard (\emph{Introduction à l'étude de la médecine expérimentale}, 1865) a montré que le vivant obéit à des lois déterminées et fondé la médecine expérimentale.
\item Applications : l'essai clinique (groupe traité / groupe témoin) pour les vaccins (ex. RTS,S contre le paludisme) et antipaludiques ; les parcelles expérimentales de l'IRAD pour le maïs, le manioc et le cacao.
\item L'expérimentation sur le vivant a des limites : complexité du vivant et exigences éthiques (consentement, respect de la personne et de l'animal).
\end{itemize}
\end{aretenir}

\coursec{La valeur de la science : utilité, limites et problèmes}

Dans la section précédente, nous avons montré comment la méthode expérimentale — observation, hypothèse, expérimentation, vérification — permet à la science d'établir des connaissances rigoureuses et de les appliquer, y compris sur la matière vivante. Mais une fois que la science a prouvé son efficacité, une nouvelle question surgit : la science est-elle un bien en soi ? Tout ce qu'elle rend possible est-il souhaitable ? Autrement dit, quelle est la \cle{valeur} de la science ? C'est à cette question, qui prépare directement la dissertation, que cette section est consacrée.

\courssub{Qu'appelle-t-on la « valeur » de la science ?}

\begin{definition}[La valeur de la science]
La \cle{valeur de la science} désigne à la fois sa \textbf{vérité} (l'exactitude de ses connaissances, garantie par la méthode) et son \textbf{utilité} (ce qu'elle apporte concrètement à l'humanité : santé, nourriture, communication, maîtrise de la nature). Interroger la valeur de la science, c'est donc se demander si elle dit vrai et si elle fait du bien.
\end{definition}

L'intuition est simple : lorsqu'un enfant fiévreux à Yaoundé guérit grâce à un médicament dosé avec précision, nous tenons la science pour précieuse. Mais lorsque le même mot « science » est associé à la bombe atomique ou à la pollution, nous hésitons. La valeur de la science est donc un \emph{problème} philosophique et non une évidence : il faut peser ses bienfaits et ses méfaits, et déterminer qui porte la responsabilité de ses usages.

\courssub{La valeur positive de la science : expliquer, prévoir, agir}

\courssub{Un triple pouvoir : expliquer, prévoir, pouvoir}

La science remplit d'abord trois fonctions que résume une formule célèbre d'Auguste \cle{Comte} (1798--1857), fondateur du positivisme :

\begin{aretenir}[La formule d'Auguste Comte]
« \textbf{Savoir pour prévoir, prévoir pour pouvoir.} » La science donne la connaissance des lois de la nature (\emph{savoir}) ; cette connaissance permet d'anticiper les phénomènes (\emph{prévoir}) ; cette prévision permet enfin d'agir efficacement sur le monde (\emph{pouvoir}).
\end{aretenir}

Prenons un exemple camerounais. Le paysan de l'Adamaoua qui connaît le cycle des pluies (\emph{savoir}) peut annoncer la saison favorable aux semis (\emph{prévoir}) et choisir le bon moment pour planter son maïs (\emph{pouvoir}). À une échelle plus grande, la météorologie scientifique prolonge exactement cette démarche avec des satellites et des modèles mathématiques.

\courssub{Les bienfaits concrets de la science}

Les applications de la science ont transformé la condition humaine. Trois domaines illustrent cette valeur positive :

\begin{itemize}
\item \textbf{La médecine.} La vaccination, l'hygiène et les antibiotiques ont fait reculer les maladies infectieuses. Au Cameroun, les campagnes de vaccination (BCG, poliomyélite, rougeole, et plus récemment le paludisme) menées dans le cadre du Programme élargi de vaccination (PEV) depuis 1982 ont contribué à réduire sensiblement la mortalité des enfants de moins de cinq ans au cours des dernières décennies. Là où la tradition ne pouvait qu'invoquer le sort, la science identifie l'agent de la maladie et le neutralise.
\item \textbf{L'agriculture.} La sélection des variétés, les engrais, l'irrigation et la lutte contre les parasites ont augmenté les rendements : le savoir agronomique permet de nourrir davantage d'hommes sur les mêmes surfaces.
\item \textbf{La communication.} Le téléphone mobile et l'Internet relient aujourd'hui un étudiant de Maroua aux bibliothèques du monde entier ; le transfert d'argent par téléphone (Mobile Money) a transformé le commerce jusque dans les villages.
\end{itemize}

\begin{exemplebox}[La vaccination au Cameroun]
Grâce aux programmes de vaccination systématique menés depuis les années 1980 (Programme élargi de vaccination), des maladies autrefois meurtrières comme la poliomyélite ont quasiment disparu du territoire camerounais (certification « sans polio sauvage » en 2020), et la mortalité des enfants de moins de cinq ans a fortement diminué. C'est une preuve directe de la valeur positive de la science : une connaissance exacte des microbes (Pasteur) a permis de prévoir l'infection et de la prévenir.
\end{exemplebox}

\courssub{Les problèmes de la science : quand les applications nuisent}

\courssub{Le revers du progrès}

La même science qui guérit peut aussi détruire. Ses applications techniques ont des effets pervers que l'on ne peut nier :

\begin{itemize}
\item \textbf{Les armes.} La physique nucléaire a rendu possible la bombe atomique ; la chimie a produit les gaz de combat ; l'informatique sert aujourd'hui aux drones armés.
\item \textbf{La pollution.} L'industrie chimique et pétrolière dégrade l'air, les sols et les fleuves ; les plastiques envahissent les côtes de Kribi et de Limbé.
\item \textbf{Les déchets technologiques.} À Douala comme dans de grandes villes africaines, des décharges de \cle{déchets électroniques} (vieux ordinateurs, téléphones, téléviseurs importés d'Europe ou d'Asie) libèrent du plomb, du mercure et du cadmium, qui empoisonnent les récupérateurs — souvent des enfants — et contaminent les nappes phréatiques.
\end{itemize}

\begin{popfigure}
\begin{center}
\begin{tikzpicture}[scale=0.9, every node/.style={font=\small}]
% sol
\fill[popcream] (0,0) rectangle (12,1.1);
\draw[popdark] (0,1.1) -- (12,1.1);
% tas de déchets
\fill[popmuted!60] (2,1.1) .. controls (3,3.2) and (5,3.6) .. (6.2,1.1) -- cycle;
\fill[popmuted!40] (5.5,1.1) .. controls (6.8,2.8) and (8.6,3.0) .. (9.8,1.1) -- cycle;
% objets
\draw[popdark, fill=popblueL] (3.2,1.6) rectangle (4.2,2.3);
\draw[popdark] (3.35,1.75) rectangle (4.05,2.15);
\node[font=\scriptsize] at (3.7,2.55) {écran};
\draw[popdark, fill=poporangeL] (6.3,1.5) rectangle (7.1,2.1);
\node[font=\scriptsize] at (6.7,2.35) {téléphone};
\draw[popdark, fill=poppurpleL] (8.0,1.6) rectangle (9.0,2.4);
\node[font=\scriptsize] at (8.5,2.65) {ordinateur};
% fumées toxiques
\draw[popgreen, thick, decorate, decoration={snake, amplitude=0.15cm, segment length=0.3cm}] (4.5,3.3) -- (4.5,4.6);
\draw[popgreen, thick, decorate, decoration={snake, amplitude=0.15cm, segment length=0.3cm}] (7.5,2.9) -- (7.5,4.2);
\node[popgreen, font=\scriptsize, align=center] at (6.0,4.9) {vapeurs toxiques\\(plomb, mercure, cadmium)};
% infiltration
\draw[popblue, thick, ->] (4.0,1.0) -- (4.0,0.2);
\draw[popblue, thick, ->] (7.0,1.0) -- (7.0,0.2);
\node[popblue, font=\scriptsize] at (9.6,0.55) {contamination des sols et de l'eau};
% enfant récupérateur
\draw[popdark, fill=poppinkL] (1.2,1.6) circle (0.22);
\draw[popdark] (1.2,1.38) -- (1.2,0.7);
\draw[popdark] (1.2,1.2) -- (0.8,0.95);
\draw[popdark] (1.2,1.2) -- (1.6,0.95);
\draw[popdark] (1.2,0.7) -- (0.95,0.15);
\draw[popdark] (1.2,0.7) -- (1.45,0.15);
\node[font=\scriptsize, align=center] at (1.3,-0.35) {récupérateur\\exposé};
\end{tikzpicture}
\end{center}
\legende{Schéma commenté d'une décharge de déchets électroniques (type décharge de Douala) : les appareils usagés (1) dégagent des vapeurs toxiques (2), contaminent les sols et l'eau (3) et exposent les récupérateurs, souvent très jeunes (4). La « face cachée » du progrès technique.}
\end{popfigure}

\courssub{La science dit le possible, pas le souhaitable}

Le problème central est le suivant : la science répond à la question « \emph{que pouvons-nous faire ?} », mais elle est incapable de répondre à la question « \emph{que devons-nous faire ?} ». Elle établit ce qui est \textbf{possible} ; elle ne dit pas ce qui est \textbf{souhaitable}. Fabriquer une arme chimique est possible ; est-ce souhaitable ? La science, en tant que connaissance, se tait sur ce point. C'est pourquoi le philosophe Bertrand Russell pouvait affirmer que la science augmente la puissance de l'homme sans augmenter nécessairement sa sagesse.

\begin{propriete}[La neutralité de la science quant aux fins]
La science est \cle{neutre} quant aux fins : le savoir en lui-même n'est ni bon ni mauvais ; c'est l'\textbf{usage} qui en est fait qui est bon ou mauvais. \emph{Remarque pour la dissertation : cette thèse doit être discutée et non admise d'emblée — certains objectent que la recherche elle-même est orientée par des intérêts (militaires, commerciaux) qui ne sont pas neutres.}
\end{propriete}

\begin{exemplebox}[Le même savoir, deux usages opposés]
La physique nucléaire illustre parfaitement la neutralité du savoir. La \emph{même} connaissance de la fission de l'atome d'uranium permet :
\begin{itemize}
\item de produire de l'\textbf{électricité} dans une centrale nucléaire et de soigner des cancers par radiothérapie ;
\item de fabriquer la \textbf{bombe atomique} qui a détruit Hiroshima et Nagasaki en 1945.
\end{itemize}
Le savoir est un ; les usages sont opposés. Ce n'est donc pas la connaissance qui est en cause, mais la décision de s'en servir.
\end{exemplebox}

Le tableau suivant synthétise ce double visage des applications scientifiques :

\begin{popfigure}
\begin{center}
\begin{tabular}{|l|p{5.2cm}|p{5.2cm}|}
\hline
\rowcolor{popblueL} \textbf{Domaine} & \textbf{Bienfaits} & \textbf{Méfaits possibles} \\
\hline
Médecine & Vaccins, antibiotiques, chirurgie : recul de la mortalité & Expérimentations sans consentement ; trafic d'organes ; surmédicalisation \\
\hline
Nucléaire & Électricité sans pétrole ; radiothérapie contre le cancer & Bombe atomique ; accidents de centrales ; déchets radioactifs \\
\hline
Numérique & Communication instantanée ; accès au savoir ; mobile money & Surveillance des citoyens ; désinformation ; déchets électroniques (Douala) \\
\hline
OGM & Rendements agricoles accrus ; résistance aux parasites & Risques sanitaires débattus ; dépendance aux semenciers ; perte de biodiversité \\
\hline
\end{tabular}
\end{center}
\legende{Tableau à double entrée : bienfaits et méfaits des grandes applications scientifiques. À réutiliser comme exemples dans une dissertation.}
\end{popfigure}

\courssub{La distinction entre le fait et la valeur}

Pour bien poser le problème, il faut maîtriser une distinction capitale au Baccalauréat technique : celle du \cle{fait} et de la \cle{valeur}.

\begin{definition}[Fait et valeur]
Un \textbf{fait} est ce qui est, ce qui peut être observé et vérifié (« l'eau bout à $100\,^\circ\mathrm{C}$ au niveau de la mer »). Une \textbf{valeur} est ce qui doit être, ce que l'on juge bon, juste ou beau (« il est mal de mentir »). La science établit des \emph{faits} ; l'\emph{éthique} (la morale) juge les conduites au nom de valeurs.
\end{definition}

Conséquence décisive : on ne peut pas déduire logiquement un devoir d'un simple fait. Que la science constate qu'il est \emph{possible} de cloner un être humain ne prouve ni qu'il faille le faire ni qu'il faille l'interdire. Le passage du « est » au « doit être » exige un jugement moral, qui relève de la réflexion éthique et du débat citoyen, non du laboratoire.

\begin{attention}[Erreur fréquente : accuser « la science »]
Dans les copies, on lit souvent : « la science a tué des millions de personnes à Hiroshima » ou « la science pollue la planète ». C'est une erreur d'analyse. Il faut distinguer trois niveaux :
\begin{enumerate}
\item la \textbf{connaissance} (le savoir théorique, neutre) ;
\item la \textbf{technique} (l'application du savoir, qui peut être bien ou mal orientée) ;
\item les \textbf{décisions politiques et économiques} (gouvernements, entreprises) qui commandent et financent ces applications.
\end{enumerate}
Ce ne sont pas les équations d'Einstein qui ont détruit Hiroshima, mais une décision politique et militaire utilisant une technique fondée sur ces équations. Accuser « la science » en bloc, c'est confondre le savoir avec l'usage qu'on en fait.
\end{attention}

\courssub{La responsabilité du savant et de la société}

Si la science est neutre mais que ses usages ne le sont pas, alors quelqu'un doit en répondre : c'est le problème de la \cle{responsabilité}.

\begin{savaistu}[Einstein et la bombe atomique]
En 1939, le physicien Albert Einstein, alerté par d'autres savants, signa une lettre au président américain Roosevelt pour l'avertir que l'Allemagne nazie pourrait fabriquer une bombe atomique — ce qui lança le programme américain. Après la destruction d'Hiroshima en 1945, Einstein, pacifiste convaincu, déclara regretter amèrement cette signature : « Si j'avais su que les Allemands ne réussiraient pas, je n'aurais rien fait. » Ce regret illustre dramatiquement le problème de la responsabilité du savant : celui qui produit le savoir peut-il rester indifférent à l'usage qui en sera fait ?
\end{savaistu}

\begin{savaistu}[Hans Jonas et le principe responsabilité]
Face aux menaces planétaires (nucléaire, écologique, biotechnologies), le philosophe Hans Jonas (1903--1993) propose un nouvel impératif éthique : « Agis de façon que les effets de ton action soient compatibles avec la permanence d'une vie authentiquement humaine sur terre. » Ce \emph{principe responsabilité} étend l'éthique aux générations futures et à la biosphère : la science et la technique doivent être évaluées à l'aune de leurs conséquences à long terme.
\end{savaistu}

Deux niveaux de responsabilité doivent être distingués :

\begin{itemize}
\item \textbf{La responsabilité du savant.} Le chercheur ne peut pas se réfugier derrière la pureté de sa recherche : il doit anticiper les conséquences possibles de ses travaux, alerter l'opinion (comme le firent Einstein puis les signataires du manifeste Russell-Einstein contre les armes nucléaires) et refuser de participer à certains programmes.
\item \textbf{La responsabilité de la société.} Les usages de la science ne doivent pas être décidés par les seuls savants ni par les seules entreprises : ils relèvent du \textbf{débat citoyen}. C'est pourquoi on a créé des \cle{comités d'éthique} (par exemple les comités nationaux d'éthique qui examinent les recherches médicales) où siègent des scientifiques, des philosophes, des juristes et des représentants de la société, pour dire ce qui est permis et ce qui est interdit.
\end{itemize}

\begin{methode}[Analyser un problème posé par une application scientifique (méthode de dissertation)]
Face à un sujet du type « La science est-elle un danger pour l'humanité ? », procéder en quatre étapes :
\begin{enumerate}
\item \textbf{Distinguer} savoir, technique et décision (ne pas accuser « la science » en bloc).
\item \textbf{Reconnaître} la valeur positive : la science explique, prévoit, permet d'agir (Comte ; exemples : vaccination, agriculture).
\item \textbf{Identifier} le vrai problème : la science dit le possible, pas le souhaitable ; distinguer fait et valeur.
\item \textbf{Conclure} sur la responsabilité : c'est à l'éthique, au droit et au débat citoyen de fixer les fins ; la science doit être éclairée par la sagesse, non rejetée.
\end{enumerate}
\end{methode}

\courssub{La science n'est pas toute la connaissance}

Dernière limite, et non la moindre : la science ne répond pas à toutes les questions humaines. Elle dit \emph{comment} le monde fonctionne, mais elle ne dit pas \emph{pourquoi} il existe ni quel \emph{sens} donner à la vie. D'autres formes de connaissance répondent à d'autres besoins :

\begin{itemize}
\item l'\textbf{art} exprime et fait éprouver le beau, la joie, la douleur — une chanson ou une sculpture apaisent là où aucune équation ne le peut ;
\item la \textbf{morale} répond à la question « que dois-je faire ? » en jugeant les conduites selon le bien et le mal ;
\item la \textbf{religion} et les sagesses traditionnelles répondent aux questions du sens de la vie, de la mort et du destin — questions auxquelles la science, par méthode, renonce à répondre.
\end{itemize}

Prétendre que seule la science dit le vrai, c'est tomber dans le \cle{scientisme}, attitude qui transforme la science en nouvelle religion et méprise les autres dimensions de l'existence humaine. La philosophie invite au contraire à donner à chaque forme de connaissance sa place : la science pour maîtriser la nature, l'éthique pour conduire notre liberté, l'art et la religion pour habiter le sens de notre vie.

\begin{exoresolu}[Sujet de dissertation : « La science est-elle bonne ou mauvaise ? »]
\textbf{Énoncé.} Un camarade affirme : « La science est mauvaise car elle a produit la bombe atomique et la pollution. » Rédiger une réponse argumentée en trois paragraphes (thèse, antithèse, synthèse).
\tcblower
\textbf{Solution détaillée.}
\emph{Thèse (la science semble nuisible).} Il est vrai que les applications de la science ont des effets destructeurs : la bombe atomique (Hiroshima, 1945), la pollution industrielle, les déchets électroniques qui empoisonnent les décharges de Douala. La science a démultiplié la puissance de destruction de l'homme.

\emph{Antithèse (la science est d'abord bienfaisante).} Mais la même science a réduit la mortalité infantile au Cameroun grâce à la vaccination, nourri des millions d'hommes grâce à l'agronomie et relié les villages au monde grâce au téléphone mobile. Selon la formule de Comte, « savoir pour prévoir, prévoir pour pouvoir » : la science est d'abord un instrument de maîtrise et de progrès.

\emph{Synthèse (dépasser l'alternative).} L'alternative « bonne ou mauvaise » est mal posée : la science est \emph{neutre} quant aux fins ; elle dit ce qui est possible, non ce qui est souhaitable. Le même savoir nucléaire produit l'électricité et la bombe. Il faut donc distinguer la connaissance, la technique et les décisions politico-économiques : ce sont les usages, décidés par les hommes, qui sont bons ou mauvais. La conclusion n'est pas de rejeter la science, mais d'exiger qu'elle soit guidée par l'éthique, les comités d'éthique et le débat citoyen, car la science sans conscience, disait déjà Rabelais, « n'est que ruine de l'âme ».
\end{exoresolu}

\begin{exercice}[Pour s'entraîner]
1. Définir la valeur de la science et distinguer fait et valeur en donnant un exemple de chaque.

2. Expliquer la formule de Comte « savoir pour prévoir, prévoir pour pouvoir » à partir d'un exemple camerounais de votre choix.

3. Sujet de dissertation : « Faut-il avoir peur de la science ? » Construire le plan détaillé (introduction problématisée, trois parties, conclusion).
\end{exercice}

\begin{corrige}[Exercice]
\textbf{1.} La valeur de la science désigne sa vérité (exactitude garantie par la méthode) et son utilité (ce qu'elle apporte à l'humanité). Un fait est ce qui est observable et vérifiable : « le virus de la poliomyélite se transmet par l'eau souillée ». Une valeur est un jugement sur ce qui doit être : « il est juste de vacciner tous les enfants ». La science établit les faits ; l'éthique juge selon les valeurs.

\textbf{2.} Exemple attendu en trois temps : connaissance d'une loi (savoir), anticipation d'un phénomène (prévoir), action efficace (pouvoir). Par exemple : l'agronome connaît le cycle du niébé et les parasites qui l'attaquent (savoir) ; il prévoit la période de traitement et de récolte (prévoir) ; il obtient une meilleure récolte et lutte contre la famine (pouvoir).

\textbf{3.} Plan attendu : \emph{Introduction} — la science effraie (bombe, pollution) mais elle guérit ; problème : la science est-elle en elle-même dangereuse, ou seulement ses usages ? \emph{I.} Les raisons de craindre : armes, pollution, déchets électroniques, manipulations du vivant. \emph{II.} Les raisons de ne pas craindre la science elle-même : neutralité du savoir, distinction fait/valeur, bienfaits (vaccination, agriculture, communication). \emph{III.} Synthèse : ce qu'il faut craindre, ce n'est pas la science mais l'usage incontrôlé ; d'où la responsabilité du savant, les comités d'éthique et le débat citoyen ; la science doit être éclairée par la sagesse. \emph{Conclusion} : réponse nuancée — non pas peur de la science, mais vigilance sur ses applications.
\end{corrige}

\begin{aretenir}[L'essentiel de la section]
\begin{itemize}
\item La valeur de la science = sa \textbf{vérité} + son \textbf{utilité}.
\item La science \textbf{explique}, \textbf{prévoit} et permet d'\textbf{agir} (Comte : « savoir pour prévoir, prévoir pour pouvoir ») : vaccination, agriculture, communication.
\item Mais ses applications peuvent \textbf{nuire} : armes, pollution, déchets électroniques. La science dit le \emph{possible}, pas le \emph{souhaitable}.
\item Distinction capitale : la science établit des \textbf{faits}, l'éthique juge selon des \textbf{valeurs}.
\item La science est \textbf{neutre} quant aux fins : c'est l'usage qui est bon ou mauvais (thèse à discuter).
\item La responsabilité incombe au \textbf{savant} et à la \textbf{société} (comités d'éthique, débat citoyen, principe responsabilité de Jonas).
\item La science n'est pas toute la connaissance : l'art, la morale et la religion répondent à d'autres questions (sens, valeurs).
\end{itemize}
\end{aretenir}

\coursec{Science, bioéthique et question écologique}

Nous avons vu dans la section précédente que la science possède une valeur indéniable : elle accroît notre puissance d'agir sur la nature et sur la vie. Mais cette puissance soulève une question décisive : tout ce que la science \emph{peut} faire est-il pour autant \emph{permis} ? Peut-on, parce que la technique le permet, prolonger indéfiniment la vie d'un malade, modifier le patrimoine génétique d'un embryon, ou exploiter sans limite les forêts ? C'est ici qu'interviennent la \cle{bioéthique} et la réflexion \cle{écologique} : elles montrent que la science, pour rester au service de l'homme, doit être accompagnée d'une réflexion morale sur ses propres applications.

\courssub{La bioéthique : définition et naissance}

\begin{definition}[La bioéthique]
La \cle{bioéthique} est l'étude des questions morales soulevées par les sciences de la vie et la médecine. Elle réfléchit aux conséquences éthiques des applications de la biologie, de la médecine, de la génétique et des biotechnologies sur l'être humain et sur le vivant.
\end{definition}

L'intuition de départ est simple. Imaginez qu'un médecin dispose d'une technique capable de modifier l'ADN d'un embryon pour supprimer une maladie héréditaire. La question scientifique est : « Est-ce techniquement possible ? » La question bioéthique est tout autre : « Est-ce moralement acceptable ? Qui décide ? Jusqu'où peut-on aller ? » La bioéthique naît ainsi du \emph{décalage} entre le progrès technique, très rapide, et la réflexion morale, qui doit le suivre.

Le mot « bioéthique » apparaît au XX\textsuperscript{e} siècle, lorsque les progrès de la médecine (transplantations d'organes, réanimation, fécondation \emph{in vitro}) rendent possibles des actes que l'humanité n'avait jamais eu à juger. En Afrique et au Cameroun, ces questions se posent avec une acuité particulière : comment concilier les exigences de l'éthique médicale moderne avec les traditions, les ressources limitées des hôpitaux et la place reconnue de la médecine traditionnelle ?

\begin{attention}[Ne pas confondre bioéthique et biologie]
La \cle{biologie} est une \emph{science des faits} : elle décrit et explique le vivant (ce qui \emph{est}). La \cle{bioéthique} est une \emph{réflexion sur les valeurs} : elle juge ce qu'il est bon ou mauvais de faire (ce qui \emph{doit être}). Dire « le clonage est possible » relève de la biologie ; dire « le clonage humain est inacceptable » relève de la bioéthique. Dans une dissertation, celui qui confond les deux commet l'erreur de déduire un devoir d'un simple fait.
\end{attention}

\courssub{Les grands problèmes bioéthiques}

\textbf{Le début et la fin de vie.}\par
La médecine moderne peut prolonger la vie au-delà de ce que la nature permettait autrefois. Dès lors, deux problèmes se posent. L'\cle{euthanasie} consiste à provoquer la mort d'un malade incurable pour abréger ses souffrances : est-ce un acte de compassion ou un meurtre ? L'\cle{acharnement thérapeutique} désigne au contraire la poursuite de traitements lourds qui prolongent la vie sans espoir de guérison : faut-il tout tenter, ou savoir s'arrêter ? Ces questions opposent le respect de la vie à celui de la dignité et de la liberté du patient.

\textbf{La procréation médicalement assistée (PMA).}\par
Grâce à la fécondation \emph{in vitro}, des couples stériles peuvent avoir des enfants. Mais la PMA soulève des questions : que faire des embryons surnuméraires ? Faut-il autoriser les mères porteuses ? L'enfant devient-il un « produit » que l'on fabrique ?

\textbf{Le clonage et les manipulations génétiques.}\par
Le \cle{clonage} reproductif (créer un être humain génétiquement identique à un autre) est presque universellement condamné, car il porte atteinte à l'unicité de la personne. Les \cle{manipulations génétiques} posent des problèmes plus nuancés : la \emph{thérapie génique} (corriger un gène malade pour soigner) semble légitime, mais la modification du génome transmissible aux descendants inquiète, car elle engage les générations futures sans leur consentement. Quant aux \cle{OGM} (organismes génétiquement modifiés), ils promettent des rendements agricoles accrus — enjeu crucial pour l'Afrique — mais suscitent des débats sur la santé, la biodiversité et la dépendance des paysans envers les firmes semencières.

\courssub{Les principes de la bioéthique}

Pour trancher ces dilemmes, la bioéthique contemporaine s'appuie sur quatre grands principes, formulés notamment par les philosophes américains Tom Beauchamp et James Childress.

\begin{aretenir}[Les quatre principes de la bioéthique]
\begin{enumerate}
\item L'\cle{autonomie} : respecter la volonté et les décisions du patient, maître de son corps et de sa vie.
\item La \cle{bienfaisance} : faire le bien du patient, agir dans son intérêt.
\item La \cle{non-malfaisance} : ne pas nuire (\emph{primum non nocere}, « d'abord ne pas nuire », principe hérité d'Hippocrate).
\item La \cle{justice} : répartir équitablement les soins et les ressources médicales entre tous.
\end{enumerate}
Ces principes se concrétisent dans le \cle{consentement éclairé} : aucun acte médical ni aucune expérimentation ne peut être pratiqué sur une personne sans qu'elle ait été informée clairement et ait accepté librement.
\end{aretenir}

\begin{popfigure}
\begin{center}
\begin{tikzpicture}[scale=1.05]
\node[circle, draw=popdark, fill=popblueL, minimum size=2.4cm, align=center, font=\bfseries] (P) at (0,0) {Le\\patient};
\node[ellipse, draw=popblue, fill=popblueL!60, align=center] (A) at (0,3) {\textbf{Autonomie}\\respecter sa volonté};
\node[ellipse, draw=popgreen, fill=popgreenL, align=center] (B) at (3.6,0.9) {\textbf{Bienfaisance}\\faire son bien};
\node[ellipse, draw=poporange, fill=poporangeL, align=center] (C) at (-3.6,0.9) {\textbf{Non-malfaisance}\\ne pas nuire};
\node[ellipse, draw=poppurple, fill=poppurpleL, align=center] (J) at (0,-3) {\textbf{Justice}\\soins équitables pour tous};
\draw[<->, thick, popblue] (P) -- (A);
\draw[<->, thick, popgreen] (P) -- (B);
\draw[<->, thick, poporange] (P) -- (C);
\draw[<->, thick, poppurple] (P) -- (J);
\end{tikzpicture}
\end{center}
\legende{Les quatre principes de la bioéthique organisés autour du patient. Toute décision médicale (soin, essai clinique, greffe) doit être évaluée à la lumière de ces quatre exigences, qui peuvent entrer en tension : l'autonomie du malade peut par exemple contredire ce que le médecin juge être son bien.}
\end{popfigure}

\courssub{Cas camerounais et africains}

La bioéthique n'est pas un luxe de pays riches : elle concerne directement le Cameroun et l'Afrique.

\begin{itemize}
\item \textbf{L'accès aux soins.} Le principe de justice se heurte à la réalité : les hôpitaux manquent de personnel et de matériel, et de nombreux malades renoncent aux soins faute d'argent. La répartition équitable des ressources sanitaires entre villes et campagnes est un problème bioéthique majeur.
\item \textbf{Les essais cliniques.} Des médicaments sont parfois testés dans des populations africaines vulnérables. Le consentement éclairé exige que les participants comprennent réellement les risques : tester un produit sur des personnes analphabètes sans information véritable est une violation de l'autonomie.
\item \textbf{Le trafic d'organes.} La pauvreté expose certaines populations à la tentation de vendre un organe. La bioéthique affirme que le corps humain n'est pas une marchandise : il y a là une atteinte à la dignité de la personne.
\item \textbf{La médecine traditionnelle.} L'Organisation mondiale de la santé (OMS) reconnaît la place de la médecine traditionnelle africaine, à condition qu'elle soit évaluée scientifiquement. Le défi est de faire dialoguer les savoirs traditionnels et la démarche expérimentale, sans mépris ni naïveté.
\end{itemize}

\begin{exemplebox}[Le bassin du Congo : la science mesure, la politique décide]
Le bassin du Congo, qui s'étend notamment sur le Sud et l'Est du Cameroun, abrite la deuxième forêt tropicale du monde après l'Amazonie. Les mesures satellitaires montrent qu'il perd chaque année des milliers d'hectares sous l'effet de l'exploitation forestière, de l'agriculture itinérante et de l'urbanisation. La science fournit ici des \emph{données} (taux de déforestation, émissions de carbone, espèces menacées), mais elle ne peut pas, à elle seule, décider de ce qu'il faut faire : c'est à la politique, éclairée par l'éthique, de fixer les règles de protection.
\end{exemplebox}

\courssub{La question écologique}

La science entretient avec l'environnement un rapport double. D'une part, elle \emph{révèle} les dégradations : c'est grâce aux mesures des climatologues que nous savons que la planète se réchauffe ; c'est grâce aux images satellites que nous constatons la déforestation du bassin du Congo et l'avancée de la sécheresse dans l'Extrême-Nord du Cameroun. D'autre part, elle \emph{propose des solutions} : énergies renouvelables (solaire, hydroélectricité), agriculture durable, traitement des déchets, reforestation.

\begin{popfigure}
\begin{center}
\begin{tikzpicture}[scale=0.95]
% Contour simplifié du Cameroun
\draw[thick, popdark, fill=popcream] (1.2,0) -- (3.4,0.2) -- (4.2,1.4) -- (4.6,3.2) -- (5.4,4.6) -- (5.0,6.6) -- (3.8,7.4) -- (3.4,8.8) -- (2.2,9.4) -- (1.6,8.0) -- (2.2,6.4) -- (1.4,4.8) -- (0.4,3.4) -- (0.2,1.6) -- cycle;
% Zones
\fill[popgreen!45] (0.6,0.5) -- (3.2,0.6) -- (3.8,1.8) -- (2.4,2.6) -- (0.7,2.4) -- cycle;
\fill[popgold!50] (2.4,6.8) -- (3.8,7.3) -- (3.4,8.7) -- (2.3,9.2) -- (1.8,8.0) -- (2.3,7.0) -- cycle;
\fill[popink!35] (0.5,1.7) circle (0.35);
% Villes
\fill[popdark] (1.6,2.9) circle (0.09); \node[right, font=\small] at (1.75,2.9) {Yaoundé};
\fill[popdark] (0.5,1.7) circle (0.09); \node[left, font=\small] at (0.35,1.45) {Douala};
\fill[popdark] (2.6,8.3) circle (0.09); \node[right, font=\small] at (2.75,8.3) {Maroua};
% Nord
\draw[->, thick, popdark] (5.6,0.4) -- (5.6,1.4); \node[right, font=\small] at (5.6,1.0) {N};
% Légendes internes
\node[font=\small\bfseries, popdark] at (2.0,1.4) {Forêt menacée};
\node[font=\small\bfseries, popdark] at (2.7,7.9) {Sécheresse};
\node[font=\small\bfseries, popdark] at (1.9,0.7) {(déforestation)};
\node[font=\small\bfseries, popdark] at (2.9,7.35) {(désertification)};
\end{tikzpicture}
\end{center}
\legende{Carte très simplifiée du Cameroun localisant trois enjeux écologiques majeurs : 1. au Sud, la forêt équatoriale du bassin du Congo menacée par la déforestation (zone verte) ; 2. à l'Extrême-Nord, la sécheresse et la désertification autour de Maroua (zone dorée) ; 3. à Douala, la pollution urbaine et industrielle (zone sombre). Croquis indicatif, sans précision cartographique.}
\end{popfigure}

Face à ces enjeux, un principe directeur s'impose dans les politiques modernes : le principe de précaution.

\begin{definition}[Le principe de précaution]
Le \cle{principe de précaution} impose d'agir pour prévenir un risque grave ou irréversible pour l'environnement ou la santé, même en l'absence de certitude scientifique complète. Autrement dit : l'incertitude de la science ne doit pas servir d'excuse pour retarder les mesures de protection.
\end{definition}

Ce principe renverse la charge de la preuve. Autrefois, on attendait qu'un danger soit prouvé pour interdire une activité ; désormais, face à un risque grave (réchauffement climatique, dissémination d'OGM, pollution des nappes phréatiques), c'est à celui qui veut agir de démontrer que son action est sans danger. Il ne s'agit pas de refuser la science, mais de l'accompagner d'une sagesse pratique : quand l'enjeu est irréversible, mieux vaut prévenir que guérir.

\begin{savaistu}[Le Cameroun et les énergies renouvelables]
Le Cameroun dispose d'un potentiel hydroélectrique considérable (barrages de Songloulou, Édéa, Lom Pangar) et d'un ensoleillement exceptionnel dans le Nord. La science et la technique offrent ainsi des solutions concrètes pour produire de l'énergie sans aggraver le réchauffement climatique : la question écologique n'oppose pas développement et nature, elle invite à un développement \emph{durable}.
\end{savaistu}

\courssub{Synthèse : vers une science responsable}

\begin{methode}[Traiter un sujet de dissertation sur la bioéthique]
\begin{enumerate}
\item \textbf{Thèse} : la science doit être libre, car toute limite entrave le progrès (« tout ce qui est possible doit être tenté »). Argument : interdire la recherche, c'est renoncer à guérir des maladies.
\item \textbf{Antithèse} : la science sans limites conduit aux pires dérives (expérimentations sur l'homme, eugénisme, destruction de la nature). La puissance technique ne fonde aucun droit.
\item \textbf{Synthèse} : il faut dépasser l'opposition par l'idée d'une \emph{science responsable} : la recherche reste libre, mais ses applications sont soumises à des principes éthiques (autonomie, non-malfaisance, justice, précaution) discutés démocratiquement. La limite n'empêche pas le progrès : elle l'humanise.
\end{enumerate}
\end{methode}

\begin{exoresolu}[La science peut-elle tout se permettre ?]
\textbf{Énoncé.} Un candidat affirme dans sa copie : « Puisque la science est fondée sur la raison, elle est toujours bonne et doit être libérée de toute contrainte morale. » Discutez cette affirmation en mobilisant les notions de la leçon.
\tcblower
\textbf{Solution détaillée.} L'affirmation repose sur une confusion entre la science comme \emph{connaissance} et la science comme \emph{puissance d'agir}. Comme connaissance, la science est effectivement une valeur : elle détruit les préjugés et nous libère de l'erreur (Bachelard). Mais comme puissance, ses applications peuvent nuire : les expérimentations médicales sans consentement, le trafic d'organes, la destruction du bassin du Congo montrent que le savoir peut être mis au service du pire. La bioéthique répond précisément à ce problème : elle n'interdit pas la recherche, mais soumet ses applications à des principes — autonomie, bienfaisance, non-malfaisance, justice — et au consentement éclairé du patient. De même, le principe de précaution impose de protéger l'environnement même sans certitude absolue sur les risques. Conclusion : la science ne peut pas tout se permettre, non parce qu'elle serait mauvaise, mais parce qu'elle est un \emph{pouvoir} ; or tout pouvoir appelle une responsabilité. La véritable grandeur de la science est d'accepter d'être une science responsable, au service de l'homme et de la nature.
\end{exoresolu}

\begin{aretenir}[L'essentiel de la section]
\begin{itemize}
\item La bioéthique réfléchit aux questions morales posées par la médecine et les biotechnologies (fin de vie, PMA, clonage, manipulations génétiques).
\item Elle repose sur quatre principes : autonomie, bienfaisance, non-malfaisance, justice, concrétisés par le consentement éclairé.
\item Au Cameroun et en Afrique, elle concerne l'accès aux soins, les essais cliniques, le trafic d'organes et le dialogue avec la médecine traditionnelle.
\item La science révèle les dégradations écologiques et propose des solutions ; le principe de précaution impose d'agir malgré l'incertitude face aux risques graves.
\item Enjeu final : construire une science responsable, au service de l'homme et de la nature.
\end{itemize}
\end{aretenir}

\begin{exercice}[Pour s'entraîner]
Sujet de dissertation : « Faut-il avoir peur de la science ? » Vous construirez un plan dialectique (thèse, antithèse, synthèse) en mobilisant la bioéthique, le principe de précaution et au moins un exemple camerounais ou africain.
\end{exercice}

\newpage

\coursec{Activité d'intégration}

\coursec{La science : de la méthode expérimentale aux enjeux bioéthiques}

\courssub{Objectifs de l'activité}

À l'issue de cette activité d'intégration, l'élève sera capable de :
\begin{itemize}
\item identifier, dans un protocole de recherche réaliste, les étapes de la \cle{méthode expérimentale} (observation, hypothèse, expérience, vérification, conclusion) ;
\item interpréter des données chiffrées d'un essai clinique et les représenter par un diagramme en bâtons comparatif ;
\item repérer, dans des réactions sociales concrètes, les \cle{obstacles épistémologiques} décrits par Bachelard (opinion, expérience première, animisme, obstacle verbal) ;
\item discuter les conditions éthiques d'une recherche sur l'être humain (consentement éclairé, groupe témoin, rapport bénéfices/risques) ;
\item rédiger une conclusion argumentée répondant à un problème philosophique : « la science peut-elle tout se permettre sur l'homme ? ».
\end{itemize}

Cette activité mobilise ainsi les deux savoir-faire officiels de l'unité : \emph{appliquer les notions de l'unité à des situations concrètes de la vie courante} et \emph{rédiger et justifier une démarche, une réponse ou une conclusion}.

\courssub{Situation}

\textbf{Un essai clinique à Buea : la science à l'épreuve de l'éthique.}

Le paludisme reste la première cause de consultation dans les hôpitaux camerounais. En 2025, une équipe de chercheurs de la Faculté des Sciences de la Santé de l'Université de Buea (région du Sud-Ouest) met au point un nouveau traitement antipaludique à base d'artémisinine associée à une plante locale, le \emph{MBO-25}. Avant toute mise sur le marché, un \cle{essai clinique} est organisé à l'hôpital régional de Buea, avec l'autorisation du Comité national d'éthique.

\textbf{Document 1 — Le protocole et les résultats.} Les chercheurs recrutent 400 patients volontaires atteints de paludisme simple, répartis au hasard en deux groupes de 200 : le \cle{groupe traité} reçoit le MBO-25, le \cle{groupe témoin} reçoit le traitement classique (quinine). Les deux groupes sont comparables (âge moyen 32 ans, même gravité initiale). Après 14 jours, on mesure le taux de guérison et la fréquence des effets secondaires :

\begin{center}
\begin{tabularx}{\linewidth}{|l|X|X|}
\hline
\rowcolor{popblueL} \textbf{Critère} & \textbf{Groupe traité (MBO-25)} & \textbf{Groupe témoin (quinine)} \\
\hline
Nombre de patients & 200 & 200 \\
\hline
Guérisons complètes au 14\textsuperscript{ème} jour & 176 & 148 \\
\hline
Effets secondaires légers (nausées) & 22 & 40 \\
\hline
Effets secondaires graves & 4 & 6 \\
\hline
\end{tabularx}
\end{center}

Hypothèse des chercheurs : « le MBO-25 guérit plus vite et provoque moins d'effets secondaires que la quinine ». Coût du traitement MBO-25 : \num{1500} FCFA, contre \num{3500} FCFA pour la quinine.

\textbf{Document 2 — Extrait de presse (\emph{Le Messager du Sud-Ouest}).} « Dans plusieurs villages autour de Buea, des familles refusent de participer à l'essai. "Le paludisme qui ne guérit pas, c'est la sorcellerie, pas les moustiques", déclare un notable. "Les Blancs et leurs comprimés veulent nous utiliser comme cobayes." Une rumeur affirme que le médicament rend stérile. Certains anciens affirment : "nos grands-parents soignaient la fièvre avec les écorces, cela suffisait ; pourquoi changer ?" »

\textbf{Document 3 — Le consentement éclairé.} « Toute recherche biomédicale sur l'être humain suppose le \cle{consentement libre et éclairé} du participant : celui-ci doit être informé, dans une langue qu'il comprend, du but de la recherche, de ses risques et de ses bénéfices, et il peut se retirer à tout moment sans perdre son droit aux soins. » (D'après la Déclaration d'Helsinki et le Comité national d'éthique du Cameroun.)

\courssub{Consignes}

Par groupes de quatre, puis individuellement pour la question 5 :

\begin{enumerate}
\item \textbf{Méthode expérimentale.} Relevez dans le document 1 les étapes de la méthode expérimentale : observation, hypothèse, expérience (avec variable et groupe témoin), résultats, conclusion. Pourquoi les chercheurs ont-ils constitué un groupe témoin ? Pourquoi la répartition au hasard ?
\item \textbf{Traitement des données.} Calculez le taux de guérison (en \%) de chaque groupe et le taux d'effets secondaires. Construisez un diagramme en bâtons comparant les taux de guérison des deux groupes. L'hypothèse des chercheurs est-elle confirmée ? Peut-on affirmer que le MBO-25 est « prouvé » efficace pour toute l'humanité ?
\item \textbf{Obstacles épistémologiques.} Dans le document 2, identifiez au moins trois obstacles épistémologiques (au sens de Bachelard) présents dans les réactions des villageois, et justifiez chaque identification.
\item \textbf{Débat éthique.} Le groupe témoin ne reçoit pas le nouveau traitement : est-ce moralement acceptable ? Le consentement éclairé est-il réellement possible dans un village où une partie de la population ne lit pas le français ? Faut-il payer les participants ? Organisez un débat : la moitié du groupe défend l'essai, l'autre moitié formule des réserves.
\item \textbf{Production individuelle.} Rédigez une conclusion argumentée d'environ 15 lignes répondant à la question : \emph{« la science peut-elle tout se permettre sur l'homme ? »} Vous mobiliserez au moins deux notions du cours et un exemple de la situation.
\end{enumerate}

\courssub{Ressources à mobiliser}

\begin{aretenir}[Notions utiles de la leçon]
\begin{itemize}
\item \textbf{Méthode expérimentale} (Claude Bernard) : observation $\rightarrow$ hypothèse $\rightarrow$ expérience contrôlée $\rightarrow$ vérification $\rightarrow$ conclusion provisoire. Toute conclusion scientifique est révisable.
\item \textbf{Obstacles épistémologiques} (Bachelard) : l'opinion, l'expérience première, l'animisme (expliquer par des esprits), l'obstacle verbal (les mots qui font croire qu'on sait), la tradition.
\item \textbf{Bioéthique} : la science dit ce qu'elle \emph{peut} faire, non ce qu'elle \emph{doit} faire ; le consentement éclairé, la dignité humaine et la proportion bénéfices/risques limitent la recherche.
\item \textbf{Calculs} : taux $= \dfrac{\text{effectif de la catégorie}}{\text{effectif total}} \times 100$.
\end{itemize}
\end{aretenir}

\courssub{Démarche et corrigé commenté}

\begin{exoresolu}[Corrigé commenté de l'activité]
\textbf{Énoncé.} Analyser le protocole de l'essai du MBO-25, interpréter les données, repérer les obstacles épistémologiques et conclure sur la question : « la science peut-elle tout se permettre sur l'homme ? »

\tcblower

\textbf{1. Les étapes de la méthode expérimentale.}
\begin{itemize}
\item \emph{Observation} : le paludisme résiste parfois aux traitements classiques ; la plante locale est utilisée traditionnellement contre la fièvre.
\item \emph{Hypothèse} : le MBO-25 guérit plus vite et provoque moins d'effets secondaires que la quinine. C'est une conjecture claire et vérifiable.
\item \emph{Expérience} : deux groupes de 200 patients, répartis au hasard pour éviter les biais (ni les chercheurs ni les préférences ne choisissent qui reçoit quoi). Le \cle{groupe témoin} sert de référence : sans lui, impossible de savoir si les guérisons sont dues au MBO-25 ou à l'évolution naturelle de la maladie.
\item \emph{Résultats et vérification} : comparaison chiffrée des deux groupes.
\item \emph{Conclusion provisoire} : à confirmer par d'autres essais sur des échantillons plus larges.
\end{itemize}

\textbf{2. Traitement des données.}
\begin{align*}
\text{Taux de guérison (MBO-25)} &= \frac{176}{200}\times 100 = 88\,\% \\
\text{Taux de guérison (quinine)} &= \frac{148}{200}\times 100 = 74\,\% \\
\text{Effets secondaires légers} &= \frac{22}{200}\times 100 = 11\,\% \quad \text{contre} \quad \frac{40}{200}\times 100 = 20\,\%
\end{align*}

\begin{center}
\begin{tikzpicture}
\begin{axis}[ybar, bar width=0.9cm, width=0.7\linewidth, height=5.5cm,
ylabel={Taux de guérison (\%)}, symbolic x coords={MBO-25, Quinine},
xtick=data, ymin=0, ymax=100, nodes near coords,
legend style={at={(0.5,-0.18)},anchor=north}]
\addplot[fill=popblue] coordinates {(MBO-25,88) (Quinine,74)};
\end{axis}
\end{tikzpicture}
\end{center}

L'écart de 14 points va dans le sens de l'hypothèse, qui est donc \emph{confirmée provisoirement} sur cet échantillon. Mais on ne peut pas dire que l'efficacité est « prouvée » pour toute l'humanité : la science expérimentale procède par induction, et ses vérités sont révisables (400 patients d'une seule région ne représentent pas tous les cas ; il faut répéter l'essai).

\textbf{3. Obstacles épistémologiques.}
\begin{itemize}
\item \emph{Obstacle animiste} : « c'est la sorcellerie, pas les moustiques » — on explique un phénomène naturel par des intentions d'esprits, ce qui dispense de chercher la cause réelle (le plasmodium transmis par l'anophèle).
\item \emph{Obstacle de l'expérience première et de la tradition} : « nos grands-parents soignaient avec les écorces, cela suffisait » — l'habitude et le prestige des anciens bloquent l'examen critique.
\item \emph{Obstacle de l'opinion} : la rumeur de stérilité, répétée sans preuve, fonctionne comme un préjugé qui se substitue à la vérification.
\end{itemize}
Attention cependant : la méfiance historique envers la recherche (« cobayes ») n'est pas un simple obstacle ; elle rappelle que l'éthique doit encadrer la science.

\textbf{4. Débat éthique.} Pour l'essai : le groupe témoin reçoit tout de même le traitement classique (personne n'est privé de soins) ; le bénéfice collectif est grand (traitement à \num{1500} FCFA contre \num{3500} FCFA). Réserves : le consentement éclairé exige une information dans une langue comprise (traduction en pidgin ou langues locales) ; payer trop les participants pourrait fausser la liberté du consentement des plus pauvres.

\textbf{5. Éléments de conclusion (question individuelle).} La science ne peut pas tout se permettre sur l'homme. Certes, la recherche exige l'expérimentation, et sans essais cliniques aucun progrès thérapeutique n'est possible : l'essai de Buea montre qu'une méthode rigoureuse peut améliorer la santé de tous à moindre coût. Mais l'homme n'est pas une matière comme les autres : il est une personne douée de dignité. C'est pourquoi la bioéthique impose des limites — consentement éclairé, proportion des risques, contrôle par un comité indépendant. La science répond à la question « que pouvons-nous faire ? », non à « que devons-nous faire ? » : c'est la conscience morale et le débat démocratique qui fixent les frontières du permis.
\end{exoresolu}

\courssub{Bilan de l'activité}

Cette activité a permis de mettre en évidence trois acquis essentiels. D'abord, la \cle{méthode expérimentale} n'est pas une recette abstraite : elle s'incarne dans des protocoles concrets où le groupe témoin, le hasard de la répartition et le calcul des taux garantissent l'objectivité. Ensuite, l'esprit scientifique se construit \emph{contre} des obstacles : sorcellerie, tradition, rumeur — autant d'obstacles épistémologiques que Bachelard invite à dépasser par la preuve. Enfin, la valeur de la science est réelle mais non absolue : l'expérimentation sur l'homme exige un cadre bioéthique (consentement éclairé, dignité, comité d'éthique). L'élève comprend ainsi que la rationalité scientifique et l'exigence morale ne s'opposent pas : la seconde protège les conditions humaines de la première. Cette articulation entre savoir, pouvoir et devoir constitue le cœur de la problématique de la leçon et prépare directement à la dissertation au Baccalauréat technique.

\newpage

\coursec{Méthodes et exercices résolus}

\coursec{Leçon 10 : La science}

\courssub{Qu'est-ce que la science ? Nature et caractères}

\begin{definition}[Science]
La \cle{science} (du latin \emph{scientia}, connaissance) est un ensemble de connaissances \textbf{rationnelles}, \textbf{objectives}, \textbf{méthodiques} et \textbf{vérifiables}, portant sur des objets déterminés, obtenues par l'observation, l'expérimentation et le raisonnement, et soumises au critère de la \cle{réfutabilité} (Popper) ou de la \cle{rupture épistémologique} (Bachelard).
\end{definition}

\begin{propriete}[Caractères de la scientificité]
Un savoir est scientifique s'il possède les caractères suivants :
\begin{enumerate}
\item \textbf{Rationalité} : il procède par concepts clairs, raisonnements logiques (déduction, induction, abduction) et arguments explicites.
\item \textbf{Objectivité} : ses résultats sont indépendants du sujet qui les produit ; ils sont intersubjectivement contrôlables (reproductibilité).
\item \textbf{Méthodicité} : il suit une démarche codifiée (méthode expérimentale, axiomatisation, enquête) et non le hasard ou l'intuition seule.
\item \textbf{Vérifiabilité / Réfutabilité} : ses énoncés peuvent être confrontés aux faits (expérience, observation, calcul) et être confirmés ou infirmés.
\item \textbf{Systématicité} : les connaissances s'organisent en théories cohérentes, non en opinions isolées.
\item \textbf{Cumulativité et historicité} : la science progresse en corrigeant ses erreurs (Bachelard : « le savoir scientifique est historique »).
\item \textbf{Universalité} : ses lois valent en tout lieu et en tout temps (dans leur domaine de validité).
\end{enumerate}
\end{propriete}

\begin{aretenir}[Distinctions fondamentales]
\begin{itemize}
\item \textbf{Science vs Opinion (\emph{doxa}) :} L'opinion est subjective, non vérifiée, liée au sentiment ou à la tradition. La science naît \emph{contre} l'opinion (rupture épistémologique).
\item \textbf{Science vs Technique :} La science est un \emph{savoir} (connaître pour comprendre) ; la technique est un \emph{savoir-faire} (connaître pour agir/transformer). La technique \emph{utilise} la science mais lui est antérieure historiquement.
\item \textbf{Science vs Philosophie / Religion :} La philosophie interroge le sens, les fondements et les limites de la science (épistémologie) ; la religion relève de la foi et du dogme, non de la démonstration.
\end{itemize}
\end{aretenir}

\begin{exemplebox}[Exemple : La médecine traditionnelle vs médecine scientifique]
Au Cameroun, un guérisseur traite le paludisme par une décoction d'écorces (savoir empirique, transmis oralement, efficace mais non standardisé). Un médecin prescrit de l'artémisinine (molécule isolée, dosage précis, efficacité prouvée par essais cliniques randomisés en double aveugle). La seconde relève de la science : méthode contrôlée, objectivité, reproductibilité, explication causale (mécanisme d'action sur le parasite).
\end{exemplebox}

\courssub{Les différents types de sciences}

\begin{definition}[Classification epistemologique des sciences]
On distingue traditionnellement trois grands types de sciences selon leur \textbf{objet} et leur \textbf{méthode} propre :
\end{definition}

\begin{center}
\begin{center}\tcbox[colback=popblue,colframe=popblue,coltext=white,arc=9pt,boxsep=4pt,left=6pt,right=6pt]{\bfseries\small Sciences}\end{center}
\vspace{-2pt}
\begin{tcbraster}[raster columns=2,raster equal height=rows,raster column skip=6pt,raster row skip=6pt,enhanced,blankest]
\begin{tcolorbox}[enhanced,colback=white,colframe=poppurple,boxrule=0.9pt,arc=3pt,title={Sciences formelles (Mathématiques, Logique)},fonttitle=\bfseries\scriptsize,coltitle=white,colbacktitle=poppurple,left=4pt,right=4pt,top=2pt,bottom=2pt,boxsep=1pt,toptitle=1.5pt,bottomtitle=1.5pt,halign=left]
\begin{itemize}[nosep,leftmargin=*,itemsep=1pt,font=\scriptsize]
\item Objet : Idéal, abstrait
\item Méthode : Déductive, axiomatique
\item Certitude : Absolue (démonstration)
\end{itemize}
\end{tcolorbox}
\begin{tcolorbox}[enhanced,colback=white,colframe=popgreen,boxrule=0.9pt,arc=3pt,title={Sciences de la matière (Physique, Chimie, Biologie)},fonttitle=\bfseries\scriptsize,coltitle=white,colbacktitle=popgreen,left=4pt,right=4pt,top=2pt,bottom=2pt,boxsep=1pt,toptitle=1.5pt,bottomtitle=1.5pt,halign=left]
\begin{itemize}[nosep,leftmargin=*,itemsep=1pt,font=\scriptsize]
\item Objet : Nature inanimée \& vivante
\item Méthode : Expérimentale / Hypothético-déductive
\item Certitude : Probabiliste, provisoire
\end{itemize}
\end{tcolorbox}
\begin{tcolorbox}[enhanced,colback=white,colframe=poporange,boxrule=0.9pt,arc=3pt,title={Sciences humaines (Sociologie, Psychologie, Économie, Histoire)},fonttitle=\bfseries\scriptsize,coltitle=white,colbacktitle=poporange,left=4pt,right=4pt,top=2pt,bottom=2pt,boxsep=1pt,toptitle=1.5pt,bottomtitle=1.5pt,halign=left]
\begin{itemize}[nosep,leftmargin=*,itemsep=1pt,font=\scriptsize]
\item Objet : Homme, société, culture
\item Méthode : Compréhension / Explication, Statistique, Enquête
\item Certitude : Probabiliste, contextuelle
\end{itemize}
\end{tcolorbox}
\end{tcbraster}

\end{center}
\legende{Figure 1 : Les trois grands types de sciences (carte mentale).}

\begin{propriete}[Spécificités de chaque type]
\begin{itemize}
\item \textbf{Sciences formelles (Mathématiques) :} Objet \emph{idéal} (nombres, figures, structures). Méthode \emph{axiomatique et déductive} : on part d'axiomes pour démontrer des théorèmes. Vérité \emph{analytique} (tautologique), nécessaire, a priori. Pas d'expérience sensible. Exemple : Théorème de Pythagore.
\item \textbf{Sciences de la matière (Physique, Chimie, Biologie) :} Objet \emph{réel, naturel}. Méthode \emph{expérimentale / hypothético-déductive} : observation $\to$ hypothèse $\to$ expérience contrôlée $\to$ loi/théorie. Vérité \emph{synthétique}, contingente, a posteriori, toujours révisable. Exemple : Loi de la gravitation, ADN.
\item \textbf{Sciences humaines (Sociologie, Économie, Psychologie, Histoire) :} Objet \emph{humain} (sujet \emph{et} objet). Méthode \emph{mixte} : compréhension (Weber : \emph{Verstehen}, sens subjectif) + explication causale (statistiques, enquêtes, comparaison). Difficulté : réflexivité (l'observateur influence l'observé), impossibilité souvent de l'expérimentation contrôlée. Vérité \emph{contextuelle}, probabiliste. Exemple : Enquête sur l'exode rural à l'INS (Institut National de la Statistique).
\end{itemize}
\end{propriete}

\begin{attention}[Piège fréquent]
Ne pas dire « les sciences humaines ne sont pas des sciences parce qu'on n'y fait pas d'expériences ». Elles sont des sciences \emph{à part entière} car elles appliquent l'exigence de méthode, d'objectivité et de vérification à un objet spécifique qui impose des méthodes adaptées (statistiques, enquête, analyse de discours, modélisation). Le critère de scientificité n'est pas l'identité de méthode mais la \textbf{rigueur méthodique} adaptée à l'objet.
\end{attention}

\courssub{L'émergence de l'esprit scientifique et les obstacles épistémologiques}

\begin{savaistu}[Naissance de l'esprit scientifique]
\begin{itemize}
\item \textbf{Antiquité grecque (VI\textsuperscript{e}--IV\textsuperscript{e} av. J.-C.) :} Passage du mythe au \emph{logos} (Thalès, Aristote). Recherche de causes naturelles (arché) et non surnaturelles. Mathématiques (Euclide) et astronomie (Ptolémée).
\item \textbf{Révolution scientifique (XVI\textsuperscript{e}--XVII\textsuperscript{e} s.) :} Galilée (mathématisation de la nature, expérience), Bacon (induction, « savoir c'est pouvoir »), Descartes (doute méthodique, mathesis universalis), Newton (synthèse mathématique/expérimentale). Naissance de la \cle{physique mathématique}.
\item \textbf{XIX\textsuperscript{e}--XX\textsuperscript{e} s. :} Professionnalisation, laboratoires, revues, spécialisation. Épistémologie comme discipline (Bachelard, Canguilhem, Popper, Kuhn).
\end{itemize}
\end{savaistu}

\begin{definition}[Obstacles épistémologiques (Gaston Bachelard, \emph{La Formation de l'esprit scientifique}, 1938)]
L'esprit scientifique ne va pas de soi : il se construit \emph{contre} des obstacles internes à la raison courante. La \cle{rupture épistémologique} est le franchissement de ces obstacles.
\end{definition}

\begin{center}
\begin{tabularx}{\linewidth}{|l|X|X|}\hline
\rowcolor{popblueL} \textbf{Obstacle} & \textbf{Description} & \textbf{Exemple concret} \\ \hline
\textbf{Expérience première} & Confiance naïve dans le donné sensible immédiat (voir, toucher). & « Le soleil tourne autour de la terre, je le vois. » \\ \hline
\textbf{Obstacle substantialiste} & Croire que la matière a des « essences » fixes, des qualités substantielles. & « Le feu contient le principe de chaleur (phlogistique). » \\ \hline
\textbf{Obstacle animiste / finaliste} & Attribuer une intention, une âme, une finalité à la nature. & « La nature fait bien les choses / la plante \emph{veut} grandir. » \\ \hline
\textbf{Obstacle réaliste / visuel} & Croire que connaître, c'est reproduire une image mentale (représentation). & Modèle atomique « billes » ; refuser la dualité onde/particule. \\ \hline
\textbf{Obstacle verbal} & Croire que nommer, c'est expliquer ; les mots figent le réel. & Dire « c'est la \emph{pesanteur} » sans donner la loi mathématique. \\ \hline
\textbf{Obstacle de l'autorité} & Accepter un savoir parce qu'un « maître » ou la tradition le dit. & « Aristote l'a dit / Le chef du village l'a dit / La rumeur court. » \\ \hline
\textbf{Obstacle affectif / passionnel} & La peur, le désir, l'espoir déforment le jugement. & Refuser un vaccin par peur de la stérilité (rumeur). \\ \hline
\textbf{Obstacle quantitatif / qualitatif} & Refuser la mathématisation (qualitatif) ou réduire le réel au chiffre (quantitatif). & Refuser les statistiques en sociologie / « tout est nombre ». \\ \hline
\end{tabularx}
\end{center}
\legende{Tableau 1 : Les principaux obstacles épistémologiques selon Bachelard.}

\begin{experience}[TP : Identifier les obstacles dans un discours courant]
\textbf{Matériel :} Extraits de forums, réseaux sociaux, conversations (ex. : « Les OGM sont des poisons », « Mon grand-père fumait et a vécu 90 ans, donc le tabac n'est pas dangereux », « Ce médicament est naturel donc sans danger »).
\textbf{Protocole :} En groupes, relever l'affirmation, identifier l'obstacle(s) bachelardien(s) à l'œuvre, formuler la question scientifique qui permettrait de lever l'obstacle.
\textbf{Interprétation :} Montre que l'esprit scientifique est une \emph{conquête} quotidienne, pas un don inné.
\end{experience}

\courssub{La méthode expérimentale et ses applications}

\begin{definition}[Méthode expérimentale (Claude Bernard, \emph{Introduction à l'étude de la médecine expérimentale}, 1865)]
Démarche rigoureuse qui permet de \textbf{provoquer} les phénomènes pour les \textbf{interroger} sous contrôle, afin d'établir des relations déterministes (causales). Elle distingue l'\cle{observateur} (voit ce qui se passe seul) de l'\cle{expérimentateur} (fait varier une cause pour mesurer l'effet).
\end{definition}

\begin{methode}[Étapes de la méthode expérimentale (type hypothético-déductif)]
\begin{enumerate}
\item \textbf{Observation / Problématisation :} Constater un fait, poser une question précise.
\item \textbf{Hypothèse :} Proposition d'explication testable, formulée en termes causaux (si cause $\to$ alors effet).
\item \textbf{Expérience contrôlée :} Mise en œuvre d'un \cle{protocole} où l'on isole \textbf{une seule variable indépendante} (cause) pour mesurer son effet sur la \textbf{variable dépendante} (effet), en neutralisant les variables parasites (groupe \cle{témoin}, randomisation, double aveugle si possible).
\item \textbf{Analyse des résultats :} Mesures, traitement statistique, comparaison aux prédictions de l'hypothèse.
\item \textbf{Conclusion :} Confirmation (corroboration), infirmation (réfutation) ou révision de l'hypothèse. Formulation d'une \cle{loi} ou intégration dans une \cle{théorie}.
\item \textbf{Communication / Validation par les pairs :} Publication, reproductibilité.
\end{enumerate}
\end{methode}

\begin{exemplebox}[Application à la matière vivante : Essai variétal à l'IRAD Nkolbisson]
\textbf{Problème :} Une nouvelle variété de maïs (VAR-NEW) résiste-t-elle mieux à la sécheresse que la variété locale (VAR-LOC) ?
\textbf{Hypothèse :} VAR-NEW donne un rendement supérieur sous stress hydrique.
\textbf{Protocole :} Deux parcelles adjacentes, même sol, même densité, même engrais. Parcel A : VAR-LOC (témoin). Parcel B : VAR-NEW. Arrêt irrigation à floraison. Mesure rendement (t/ha) à maturité. Répétition 4 blocs (statistique).
\textbf{Résultat :} Moyenne A = 2,1 t/ha ; Moyenne B = 3,8 t/ha (écart significatif $p<0,01$).
\textbf{Conclusion :} Hypothèse corroborée. VAR-NEW recommandée pour zone soudano-sahélienne.
\end{exemplebox}

\begin{attention}[Limites de l'expérimentation sur le vivant]
\begin{itemize}
\item \textbf{Complexité / Systémicité :} Organisme = système ouvert, non linéaire, multi-facteurs. Isoler une variable est souvent artificiel.
\item \textbf{Éthique :} On ne peut pas tout expérimenter sur l'homme (consentement, dignité, risque). $\to$ Éthique de la recherche, comités d'éthique (CNERS au Cameroun).
\item \textbf{Modélisation animale :} Extrapolation souris $\to$ homme incertaine.
\item \textbf{Statistique vs Certitude :} Résultats probabilistes (intervalles de confiance), pas de certitude absolue.
\end{itemize}
\end{attention}

\courssub{La valeur de la science : utilité, limites et problèmes}

\begin{aretenir}[La double valeur de la science]
\begin{itemize}
\item \textbf{Valeur cognitive (vérité) :} Connaître pour comprendre. Libération de l'ignorance, des peurs irrationnelles, des mythes. « La vérité nous rendra libres » (Spinoza, adapté).
\item \textbf{Valeur pratique (puissance / utilité) :} « Savoir, c'est pouvoir » (Bacon). Maîtrise de la nature : santé (vaccins, antibiotiques), alimentation (agronomie, OGM), énergie, communication, transports. Au Cameroun : barrages (Lom Pangar), téléphonie mobile, lutte anti-paludique.
\end{itemize}
\end{aretenir}

\begin{propriete}[Limites internes (épistémologiques)]
\begin{itemize}
\item \textbf{Provisoireté :} Toute théorie est falsifiable (Popper) et révisable (Kuhn : changements de paradigme). Ex. : Physique newtonienne $\to$ Relativité.
\item \textbf{Indétermination / Complexité :} Chaos déterministe, effet papillon (météorologie), principe d'incertitude (quantique), émergence (biologie).
\item \textbf{Dépendance aux modèles :} La science ne donne pas le « réel en soi » (Kant : noumène) mais des modèles mathématiques/expérimentaux validés.
\item \textbf{Incomplétude (Gödel) :} En mathématiques, tout système formel cohérent contient des énoncés indécidables.
\end{itemize}
\end{propriete}

\begin{propriete}[Limites externes (axiologiques / éthiques) -- Distinction Fait / Valeur (Hume, Weber)]
\begin{itemize}
\item La science décrit \textbf{ce qui est} (faits, lois, mécanismes). Elle ne dit pas \textbf{ce qui doit être} (valeurs, fins, bien/mal, justice).
\item Exemple : La science dit \emph{comment} faire une bombe atomique ou modifier un embryon (CRISPR-Cas9). Elle ne dit pas \emph{si} on doit le faire.
\item \textbf{Neutralité axiologique (Weber) :} Le savant \emph{en tant que savant} doit viser l'objectivité ; le citoyen \emph{en tant que citoyen} choisit les fins. Mais le savant a une \textbf{responsabilité} sur les usages de ses découvertes (Jonas, principe responsabilité).
\end{itemize}
\end{propriete}

\begin{savaistu}[Problèmes contemporains au Cameroun]
\begin{itemize}
\item \textbf{Technocratie :} Décisions techniques (barrages, OGM, urbanisme) imposées sans concertation démocratique.
\item \textbf{Fracture numérique / épistémique :} Accès inégal aux savoirs scientifiques (écoles rurales vs urbaines, genre).
\item \textbf{Pseudosciences / Charlatanisme :} Produits « miracles » non testés, exploitation de la crédulité (médias sociaux).
\item \textbf{Risque industriel :} Pollution (usines, mines, déchets plastiques), sécurité alimentaire (pesticides).
\end{itemize}
\end{savaistu}

\courssub{Science, bioéthique et question écologique}

\begin{definition}[Bioéthique]
Réflexion éthique interdisciplinaire (biologie, médecine, droit, philosophie, sociologie) sur les \textbf{applications des sciences de la vie} à l'être humain, aux animaux et à l'environnement. Née dans les années 1970 (Van Rensselaer Potter, André Hellegers) face aux pouvoirs nouveaux : procréation médicalement assistée (PMA), génomique, fin de vie, expérimentation.
\end{definition}

\begin{propriete}[Principes fondateurs de la bioéthique (Rapport Belmont, 1979 ; Déclaration d'Helsinki)]
\begin{enumerate}
\item \textbf{Respect de l'autonomie / Consentement libre et éclairé :} La personne est une fin (Kant), non un moyen. Information complète, compréhension, absence de contrainte.
\item \textbf{Bienfaisance :} Maximiser les bénéfices, minimiser les risques.
\item \textbf{Non-malfaisance :} \emph{Primum non nocere} (d'abord ne pas nuire).
\item \textbf{Justice / Équité :} Accès équitable aux bénéfices de la recherche et aux soins ; protection des populations vulnérables.
\end{enumerate}
\end{propriete}

\begin{definition}[Question écologique \& Responsabilité (Hans Jonas, \emph{Le Principe responsabilité}, 1979)]
L'écologie scientifique montre l'interdépendance des écosystèmes et la finitude des ressources. La technique moderne donne une \textbf{puissance démesurée} (nucléaire, chimie, climat) dont les effets sont \textbf{globaux, irréversibles, différés} et menacent l'\textbf{humanité future}. L'éthique traditionnelle (prochaine, immédiate) est inadaptée.
\end{definition}

\begin{aretenir}[Nouveaux impératifs éthiques]
\begin{itemize}
\item \textbf{Principe de précaution (Rio 1992, Constitution camerounaise 1996 art. 45) :} « En cas de risque de dommage grave ou irréversible, l'absence de certitude scientifique absolue ne doit pas servir de prétexte pour différer l'adoption de mesures effectives. »
\item \textbf{Responsabilité envers les générations futures (Jonas) :} « Agis de façon que les effets de ton action soient compatibles avec la permanence d'une vie authentiquement humaine sur la Terre. » (Heuristique de la peur).
\item \textbf{Développement durable :} Répondre aux besoins du présent sans compromettre la capacité des générations futures (Rapport Brundtland 1987). Trois piliers : Environnement, Social, Économie.
\item \textbf{Justice climatique :} Les pays du Sud (dont Cameroun) subissent le plus les effets (sécheresses, inondations, montée des eaux) tout en ayant le moins émis.
\end{itemize}
\end{aretenir}

\begin{exemplebox}[Cas concrets camerounais]
\begin{itemize}
\item \textbf{Bioéthique :} Essais cliniques vaccins (COVID-19, paludisme RTS,S) : respect du consentement éclairé en langues locales, indemnisation, comité d'éthique CNERS. Débat sur l'avortement sécurisé (Code pénal vs santé publique).
\item \textbf{Écologie :} Déforestation bassin du Congo (bois d'œuvre, huile de palme, agriculture sur brûlis) $\to$ perte biodiversité, érosion, climat. Pollution plastique (Douala, Yaoundé) $\to$ loi 2012 sur emballages non biodégradables (application difficile). Gestion déchets électroniques (Agbogbloshie type). Projet pipeline Tchad-Cameroun : impacts sociaux/environnementaux.
\end{itemize}
\end{exemplebox}

\coursec{Méthodologie : Dissertation, commentaire et application des notions}

\courssub{Méthode 1 : Appliquer les notions de l'unité à une situation concrète de la vie courante}

\begin{methode}[Relier une situation concrète aux notions du cours]
Face à un sujet qui part d'une situation de la vie courante (une rumeur sur les réseaux sociaux, une campagne de vaccination, un essai agricole, une croyance traditionnelle), il faut procéder par étapes :
\begin{enumerate}
\item \textbf{Lire et reformuler} la situation en une phrase simple, sans la juger.
\item \textbf{Dégager le problème philosophique} : que met-elle en jeu ? La différence entre opinion et savoir ? La méthode expérimentale ? Les limites de la science ?
\item \textbf{Choisir la notion du cours adaptée} (esprit scientifique, obstacle épistémologique, méthode expérimentale, bioéthique, question écologique) et la \textbf{définir} avec précision.
\item \textbf{Appliquer la notion à la situation} : montrer point par point en quoi la situation illustre ou contredit la notion.
\item \textbf{Mobiliser un auteur} (Bachelard, Claude Bernard, Popper, Jonas, Kant) pour appuyer l'analyse.
\item \textbf{Conclure} en répondant explicitement à la question posée, sans généraliser abusivement.
\end{enumerate}
\textbf{Réflexe :} toujours passer de la définition abstraite à l'exemple concret, puis de l'exemple concret à la conclusion. Jamais l'un sans l'autre.
\end{methode}

\begin{exoresolu}[Rumeur et esprit scientifique]
\textbf{Énoncé.} Dans un quartier de Douala, une rumeur affirme qu'un vaccin rend stérile. Beaucoup de parents refusent alors de faire vacciner leurs enfants, sans avoir vérifié l'information. En vous appuyant sur la notion d'obstacle épistémologique, analysez cette situation et montrez ce qu'exigerait une attitude scientifique.
\tcblower
\textbf{Solution détaillée.}

\textbf{1. Reformulation de la situation.} Des parents fondent une décision grave (refuser la vaccination) sur une simple rumeur, c'est-à-dire sur une croyance non vérifiée.

\textbf{2. Problème philosophique.} La situation oppose deux attitudes face au vrai : l'opinion (\emph{doxa}), qui accepte une affirmation parce qu'elle circule et fait peur, et le savoir scientifique, qui n'accepte une affirmation qu'après vérification méthodique.

\textbf{3. Notion mobilisée.} Selon Gaston Bachelard (\emph{La Formation de l'esprit scientifique}), la connaissance scientifique se construit \emph{contre} les obstacles épistémologiques : le premier obstacle est \cle{l'expérience première} (ce qu'on croit voir ou entendre directement), auquel s'ajoutent l'obstacle verbal (les mots qui impressionnent), l'obstacle de l'autorité (« on a dit que... ») et l'obstacle affectif (la peur, ici la peur de la stérilité, qui fait adhérer sans preuve).

\textbf{4. Application.} La rumeur illustre exactement ces obstacles : (a) obstacle de l'autorité — on croit parce que « tout le monde le dit » ; (b) obstacle affectif — l'angoisse remplace l'examen des faits ; (c) obstacle verbal — le mot « stérile » effraie et dispense de vérifier. Aucune observation contrôlée, aucune statistique, aucune expérience ne fonde la croyance.

\textbf{5. Attitude scientifique exigée.} Une attitude scientifique consisterait à : formuler une hypothèse claire, consulter les données des essais cliniques (échantillons, groupes témoins, taux d'effets secondaires), vérifier la source, accepter de réviser sa croyance si les faits la contredisent. Comme le dit Bachelard, « il faut connaître contre une opinion antérieure ».

\textbf{6. Conclusion.} Cette situation montre que l'esprit scientifique n'est pas un luxe de laboratoire : il protège la santé publique. Refuser la rumeur et exiger la preuve, c'est appliquer concrètement la rupture épistémologique entre opinion et savoir.
\end{exoresolu}

\courssub{Méthode 2 : Rédiger une introduction de dissertation sur la science}

\begin{methode}[Construire une introduction en 5 étapes]
\begin{enumerate}
\item \textbf{Amorce (accroche)} : partir d'un fait général et indiscutable lié au sujet (place de la science dans nos sociétés, progrès techniques au Cameroun).
\item \textbf{Analyse des termes du sujet} : définir brièvement chaque mot-clé (science, valeur, vérité, progrès...) en distinguant le sens courant du sens philosophique.
\item \textbf{Dégagement du problème} : montrer la tension ou le paradoxe contenu dans le sujet (ex. : la science est utile, mais ses applications peuvent nuire).
\item \textbf{Problématique} : formuler \textbf{une seule question} claire, qui reprend le problème sous forme interrogative.
\item \textbf{Annonce du plan} : présenter en une phrase les deux ou trois moments de la réflexion (thèse, antithèse, synthèse ou plan thématique).
\end{enumerate}
\textbf{Réflexe :} l'introduction tient en un paragraphe de 8 à 12 lignes ; elle ne contient aucun exemple détaillé ni aucune citation longue.
\end{methode}

\begin{exoresolu}[Introduction rédigée : « La science est-elle la seule voie d'accès à la vérité ? »]
\textbf{Énoncé.} Rédigez l'introduction complète d'une dissertation sur le sujet : « La science est-elle la seule voie d'accès à la vérité ? »
\tcblower
\textbf{Solution détaillée.}

\textbf{1. Amorce.} Depuis le XVII\textsuperscript{e} siècle, la science a transformé le monde : médecine, transports, communications. Au Cameroun comme ailleurs, on attend d'elle des réponses à presque tous les problèmes.

\textbf{2. Analyse des termes.} La \cle{science} désigne un ensemble de connaissances méthodiques, objectives et vérifiables, obtenues par l'observation, l'expérimentation et le raisonnement. La \cle{vérité} est l'accord de la pensée avec la réalité, ou d'un discours avec ce qu'il affirme. « Seule voie » signifie : à l'exclusion de toute autre (religion, art, philosophie, sens commun).

\textbf{3. Problème.} Or, si la science prouve ses résultats, elle ne répond pas à toutes les questions : elle dit \emph{comment} fonctionne le monde, mais non \emph{pourquoi} il existe, ni ce qui est bien ou mal. Dès lors, prétendre qu'elle est la seule voie vers la vérité revient à réduire la vérité à la seule vérité vérifiable expérimentalement.

\textbf{4. Problématique.} La science épuise-t-elle le champ du vrai, ou d'autres formes de connaissance peuvent-elles légitimement prétendre à la vérité ?

\textbf{5. Annonce du plan.} Nous montrerons d'abord que la science est bien une voie privilégiée et rigoureuse d'accès à la vérité ; nous verrons ensuite qu'elle possède des limites qui laissent place à d'autres formes de vérité ; enfin, nous demanderons comment articuler ces différentes voies sans les confondre.
\end{exoresolu}

\courssub{Méthode 3 : Construire un paragraphe argumenté (thèse d'un auteur)}

\begin{methode}[Le paragraphe type : idée, définition, auteur, exemple, transition]
Un paragraphe de dissertation philosophique suit un schéma rigoureux :
\begin{enumerate}
\item \textbf{Phrase d'idée} : énoncer clairement la thèse défendue dans le paragraphe.
\item \textbf{Définition} : définir la notion centrale mobilisée.
\item \textbf{Appui sur un auteur} : présenter brièvement la position d'un philosophe ou épistémologue (Bachelard, Claude Bernard, Popper, Jonas, Kant) avec une courte citation exacte si possible.
\item \textbf{Exemple concret} : illustrer par un cas précis, de préférence tiré de la vie courante ou du contexte camerounais.
\item \textbf{Analyse de l'exemple} : montrer explicitement en quoi l'exemple prouve l'idée (ne jamais laisser l'exemple « parler seul »).
\item \textbf{Conclusion partielle et transition} : tirer le bilan du paragraphe et annoncer l'étape suivante.
\end{enumerate}
\textbf{À retenir :} un paragraphe = une seule idée ; pas de liste d'exemples sans analyse.
\end{methode}

\begin{exoresolu}[Paragraphe rédigé sur la méthode expérimentale]
\textbf{Énoncé.} Rédigez un paragraphe argumenté montrant que la méthode expérimentale distingue la science de la simple opinion.
\tcblower
\textbf{Solution détaillée.}

\textbf{1. Phrase d'idée.} Ce qui sépare radicalement la science de l'opinion, c'est que la première soumet ses affirmations à l'épreuve méthodique de l'expérience, tandis que la seconde se contente de croire.

\textbf{2. Définition.} La \cle{méthode expérimentale} est une démarche qui part de l'observation d'un phénomène, formule une hypothèse, la soumet à une expérience contrôlée, puis conclut en confirmant ou en rejetant l'hypothèse.

\textbf{3. Appui sur un auteur.} Claude Bernard, dans l'\emph{Introduction à l'étude de la médecine expérimentale}, montre que l'expérimentateur ne subit pas les faits : il les provoque pour interroger la nature. L'expérience, écrit-il, est « la seule source de la vérité » en médecine : elle seule permet de dépasser les systèmes et les préjugés.

\textbf{4. Exemple concret.} Un agronome de l'IRAD à Nkolbisson veut savoir si une nouvelle variété de maïs résiste mieux à la sécheresse. Il ne se fie pas à l'apparence des plants : il cultive deux parcelles identiques, l'une avec l'ancienne variété (groupe témoin), l'autre avec la nouvelle, dans les mêmes conditions, puis compare les rendements mesurés.

\textbf{5. Analyse.} Ici, chaque étape de la méthode est visible : observation (la sécheresse détruit les récoltes), hypothèse (la nouvelle variété résiste mieux), expérience contrôlée (deux parcelles, une seule variable modifiée), conclusion fondée sur des mesures. L'opinion, elle, se serait contentée de dire « ce maïs a l'air plus solide ».

\textbf{6. Conclusion partielle et transition.} La méthode expérimentale garantit donc l'objectivité du savoir en le rendant vérifiable par tous. Mais cette rigueur a-t-elle des limites lorsqu'elle s'applique à la matière vivante et à l'homme ? C'est ce qu'il faut maintenant examiner.
\end{exoresolu}

\courssub{Méthode 4 : Problématiser et construire un plan dialectique}

\begin{methode}[Passer du sujet au plan en trois moments]
\begin{enumerate}
\item \textbf{Repérer l'alternative} contenue dans le sujet (ex. : la science est-elle bénéfique \emph{ou} dangereuse ?).
\item \textbf{Thèse (I)} : défendre d'abord la réponse la plus spontanée ou la plus forte, avec arguments, auteur et exemple.
\item \textbf{Antithèse (II)} : montrer les limites de cette réponse ou les raisons de la réponse opposée, avec arguments, auteur et exemple.
\item \textbf{Synthèse (III)} : dépasser l'opposition en distinguant les plans (ex. : la science comme connaissance est neutre ; ce sont ses applications qui posent problème) ou en posant des conditions (ex. : la science est bénéfique si elle est encadrée par l'éthique).
\item \textbf{Vérifier} que chaque partie répond à la problématique et que la synthèse n'est pas un simple compromis mou mais un véritable dépassement.
\end{enumerate}
\textbf{Réflexe :} annoncer le plan dans l'introduction et le suivre réellement ; une partie = au moins deux paragraphes argumentés.
\end{methode}

\begin{exoresolu}[Plan détaillé : « La science est-elle toujours bénéfique à l'homme ? »]
\textbf{Énoncé.} Établissez la problématique et le plan détaillé d'une dissertation sur le sujet : « La science est-elle toujours bénéfique à l'homme ? »
\tcblower
\textbf{Solution détaillée.}

\textbf{Problématique.} La science a libéré l'homme de bien des maux, mais ses applications (armes, pollution, manipulations du vivant) suscitent l'inquiétude : la science est-elle en elle-même un bien, ou sa valeur dépend-elle de l'usage que l'homme en fait ?

\textbf{I. La science est bénéfique : elle améliore la condition humaine.}
\begin{itemize}
\item Argument 1 : la science accroît la puissance de l'homme sur la nature (Bacon : « savoir, c'est pouvoir »). Exemple : l'électricité, les téléphones mobiles qui relient les villages camerounais.
\item Argument 2 : la médecine expérimentale sauve des vies (Claude Bernard). Exemple : les vaccins, la lutte contre le paludisme.
\end{itemize}

\textbf{II. Mais les applications de la science peuvent être nuisibles.}
\begin{itemize}
\item Argument 1 : la technique issue de la science peut détruire : armes de destruction massive, accidents industriels (Bhopal, Tchernobyl).
\item Argument 2 : la science appliquée au vivant pose des problèmes éthiques (bioéthique) : manipulation génétique, expérimentation sur l'embryon. Exemple : le clonage reproductif.
\item Argument 3 : la science ignore les fins : elle dit comment faire, non ce qu'il faut faire (elle est neutre quant aux valeurs, Weber).
\end{itemize}

\textbf{III. Synthèse : la science est un instrument dont la valeur dépend de l'usage éthique.}
\begin{itemize}
\item Distinction : la connaissance scientifique est en elle-même neutre ; le bien et le mal naissent de ses applications techniques et des choix de société.
\item Condition : il faut soumettre la science à une réflexion éthique (bioéthique, responsabilité du savant, principe de précaution) et à la question écologique, pour que le progrès technique reste un progrès humain (Jonas).
\end{itemize}

\textbf{Vérification.} Chaque partie répond à la problématique ; la synthèse dépasse l'opposition en distinguant science (connaissance) et technique (application), ce qui constitue un vrai dépassement et non un compromis.
\end{exoresolu}

\courssub{Méthode 5 : Justifier une réponse en commentaire de texte ou en question de synthèse}

\begin{methode}[Rédiger et justifier une réponse argumentée]
Pour toute question demandant de « justifier votre réponse » :
\begin{enumerate}
\item \textbf{Répondre d'abord clairement} : oui, non, ou « sous certaines conditions » — jamais de réponse évasive.
\item \textbf{Définir les notions} employées dans la question avant d'argumenter.
\item \textbf{Donner au moins deux arguments} distincts, chacun étayé par un auteur ou un exemple.
\item \textbf{Envisager l'objection} contraire et y répondre : c'est la marque d'une justification rigoureuse.
\item \textbf{Conclure} en reformulant la réponse justifiée, en une ou deux phrases.
\end{enumerate}
\textbf{Réflexe :} la justification ne repose jamais sur « je pense que » mais sur des raisons que n'importe quel lecteur peut examiner.
\end{methode}

\begin{exoresolu}[Question de synthèse : les sciences humaines sont-elles de vraies sciences ?]
\textbf{Énoncé.} Les sciences humaines (sociologie, psychologie, économie) sont-elles de véritables sciences au même titre que les mathématiques et les sciences de la matière ? Justifiez votre réponse.
\tcblower
\textbf{Solution détaillée.}

\textbf{1. Réponse claire.} Les sciences humaines sont de véritables sciences, mais des sciences dont l'objet particulier impose des méthodes adaptées et des limites propres.

\textbf{2. Définitions.} Une \cle{science} est un savoir méthodique, objectif et vérifiable. Les \cle{sciences humaines} ont pour objet l'homme : ses comportements, ses sociétés, ses productions.

\textbf{3. Premier argument.} Les sciences humaines remplissent les critères de la scientificité : elles observent des faits (enquêtes, statistiques, archives), formulent des hypothèses et les soumettent à vérification. Exemple : un sociologue qui étudie l'exode rural au Cameroun recueille des données chiffrées (INS, recensements), établit des corrélations (revenu, distance, services), publie des résultats contrôlables par d'autres chercheurs.

\textbf{4. Deuxième argument.} Elles se distinguent de l'opinion : le sociologue ne dit pas « les jeunes quittent les villages parce qu'ils sont paresseux » (jugement de valeur), il mesure des causes objectives (emploi, services, revenus, attractivité urbaine). Il y a donc bien rupture avec la \emph{doxa}, comme dans toute science (Bachelard).

\textbf{5. Objection et réponse.} Objection : l'expérimentation y est impossible ou immorale (on ne peut pas isoler un homme comme une molécule), et l'observateur fait partie de l'objet étudié, ce qui menace l'objectivité (réflexivité). Réponse : cette difficulté est réelle, mais elle est compensée par des méthodes spécifiques (comparaison inter-régionale, statistiques sur grands échantillons, distanciation méthodique, triangulation des sources). La certitude y est souvent probabiliste, mais une probabilité rigoureusement établie reste un savoir scientifique, à l'image de la météorologie ou de la sismologie.

\textbf{6. Conclusion.} Les sciences humaines sont donc de véritables sciences, non parce qu'elles copient la physique, mais parce qu'elles appliquent à l'homme l'exigence commune de toute science : la méthode, l'objectivité et la vérification.
\end{exoresolu}

\courssub{Méthode 6 : Traiter un sujet sur la bioéthique et la question écologique}

\begin{methode}[Articuler science, éthique et environnement]
Les sujets de bioéthique et d'écologie exigent de ne pas confondre les plans :
\begin{enumerate}
\item \textbf{Distinguer le « pouvoir » du « devoir »} : la science dit ce que l'homme \emph{peut} faire ; l'éthique demande ce qu'il \emph{doit} faire. Formule à retenir : « tout ce qui est techniquement possible n'est pas moralement permis ».
\item \textbf{Définir la bioéthique} : réflexion éthique sur les applications des sciences de la vie (médecine, génétique, procréation).
\item \textbf{Définir la question écologique} : problème des effets destructeurs des activités humaines et techniques sur la nature, et de la responsabilité envers les générations futures (Jonas).
\item \textbf{Illustrer par des cas concrets} : expérimentation médicale, OGM, déforestation, pollution industrielle, changement climatique.
\item \textbf{Proposer des principes de solution} : consentement éclairé, dignité humaine, principe de précaution, responsabilité, justice climatique.
\item \textbf{Conclure} sur la nécessité d'un encadrement éthique et juridique de la science, sans rejeter la science elle-même.
\end{enumerate}
\end{methode}

\begin{exoresolu}[Sujet d'application : faut-il tout ce que la science permet de faire ?]
\textbf{Énoncé.} Des chercheurs peuvent désormais modifier le patrimoine génétique d'embryons humains (CRISPR-Cas9). Faut-il autoriser tout ce que la science rend techniquement possible ? Rédigez une réponse argumentée en une vingtaine de lignes.
\tcblower
\textbf{Solution détaillée.}

\textbf{1. Réponse claire.} Non : tout ce qui est techniquement possible n'est pas pour autant moralement permis ; la science doit être encadrée par l'éthique.

\textbf{2. Distinction des plans.} La science répond à la question « que pouvons-nous faire ? » ; l'éthique pose la question « que devons-nous faire ? ». Confondre les deux, c'est faire de la puissance technique la norme de l'action, ce qui revient à abolir toute réflexion morale.

\textbf{3. Premier argument (bioéthique).} Modifier le génome d'un embryon touche à la dignité humaine : l'enfant à naître ne peut pas consentir, et ses descendants hériteront de modifications irréversibles (génome germinal). La \cle{bioéthique} impose donc des principes : respect de la personne, consentement éclairé, interdiction de traiter l'être humain comme un simple matériau. Kant rappelle qu'il faut toujours traiter l'humanité « comme une fin, jamais simplement comme un moyen ».

\textbf{4. Deuxième argument (précaution).} Les conséquences à long terme de telles manipulations sont imprévisibles (effets hors cible, mosaïcisme). Le \cle{principe de précaution} commande de s'abstenir lorsqu'un risque grave et irréversible est possible, même s'il n'est pas certain.

\textbf{5. Objection et réponse.} Objection : interdire, c'est freiner le progrès médical qui pourrait éliminer des maladies héréditaires graves (drépanocytose, myopathie). Réponse : il ne s'agit pas de rejeter la science, mais de distinguer les usages : soigner (thérapie génique somatique) peut être permis sous contrôle strict, tandis qu'améliorer ou « fabriquer » l'humain selon des préférences (eugénisme, bébé sur mesure) doit être interdit.

\textbf{6. Conclusion.} La science est une puissance admirable, mais une puissance sans norme devient un danger. C'est pourquoi la société, par la réflexion éthique et la loi (loi de bioéthique, conventions internationales), doit fixer des limites : la liberté de la recherche s'arrête où commence l'atteinte à la dignité humaine et à l'intégrité de l'espèce.
\end{exoresolu}

\begin{attention}[Les 5 erreurs qui coûtent des points au Bac]
\begin{enumerate}
\item \textbf{Confondre science et technique.} La science est une connaissance (savoir), la technique est une application (pouvoir-faire). Écrire « la science a inventé la bombe atomique » sans distinguer connaissance et usage rend l'argumentation confuse : c'est l'\emph{application} de la science qui est en cause.
\item \textbf{Ne pas définir les termes du sujet.} Commencer à argumenter sur « la science » ou « la vérité » sans les avoir définis dans l'introduction est sanctionné : toute dissertation philosophique exige l'analyse préalable des notions.
\item \textbf{Empiler les exemples sans les analyser.} Citer vaccin, téléphone, OGM, fusée sans montrer \emph{en quoi} chaque exemple prouve l'idée du paragraphe : un exemple non commenté ne vaut rien. Un seul exemple bien analysé vaut mieux que cinq exemples jetés en vrac.
\item \textbf{Attribuer à un auteur une thèse fausse ou vague.} Écrire « Bachelard dit que la science est dangereuse » est faux : Bachelard étudie les obstacles épistémologiques et la rupture entre opinion et savoir. Mieux vaut une référence sobre et exacte qu'une citation inventée.
\item \textbf{Oublier la problématique et la conclusion.} Une copie sans question clairement formulée dans l'introduction, ou qui s'arrête sans conclusion répondant explicitement au sujet, perd des points de structure. La conclusion doit reprendre la problématique et y répondre, non ouvrir sur un nouveau sujet.
\end{enumerate}
\end{attention}

\newpage

\coursec{Exercices}

\courssub{Je vérifie mes connaissances}

\begin{exercice}[QCM : Nature et types de sciences]
Parmi les affirmations suivantes, cochez la (ou les) bonne(s) réponse(s) et justifiez brièvement votre choix.
\begin{enumerate}
    \item La science se définit uniquement par son objet d'étude (la nature).
    \item Les mathématiques sont une science formelle car elles portent sur des objets idéaux et utilisent la déduction axiomatique.
    \item La sociologie est une science expérimentale au même titre que la physique.
    \item L'esprit scientifique naît spontanément du bon sens et de l'expérience courante.
    \item La méthode expérimentale comporte nécessairement les étapes : observation, hypothèse, expérience, conclusion.
\end{enumerate}
\end{exercice}

\begin{exercice}[Vrai / Faux justifié : Obstacles épistémologiques et méthode]
Indiquez si les propositions sont vraies ou fausses. Justifiez chaque réponse en une phrase en mobilisant les notions du cours (Bachelard, obstacle épistémologique, rupture, etc.).
\begin{enumerate}
    \item L'obstacle épistémologique est une simple erreur de raisonnement que l'on corrige facilement.
    \item Selon Bachelard, l'expérience première (le sens commun) est un obstacle majeur à la connaissance scientifique.
    \item La rupture épistémologique marque le passage de l'opinion au savoir scientifique.
    \item Une théorie scientifique est définitive et immuable une fois validée par l'expérience.
    \item La bioéthique a pour seul but de freiner le progrès scientifique.
\end{enumerate}
\end{exercice}

\begin{exercice}[Définitions et distinctions]
Définissez précisément les notions suivantes en deux ou trois lignes chacune, en donnant un exemple concret pour chacune :
\begin{enumerate}
    \item Science formelle (ex. : \cle{mathématiques}, \cle{logique}).
    \item Science de la nature / science expérimentale (ex. : \cle{physique}, \cle{biologie}).
    \item Science humaine / science de l'esprit (ex. : \cle{sociologie}, \cle{psychologie}).
    \item Obstacle épistémologique (selon Bachelard).
    \item Bioéthique.
\end{enumerate}
\end{exercice}

\begin{exercice}[Calcul direct et interprétation : Données écologiques au Cameroun]
Le tableau ci-dessous présente l'évolution de la surface forestière au Cameroun (en millions d'hectares) selon des données FAO/Global Forest Watch.
\begin{center}
\begin{tabularx}{\linewidth}{|c|X|X|X|X|X|}\hline
Année & 1990 & 2000 & 2010 & 2020 & 2023 \\ \hline
Surface forestière (Mha) & 22,5 & 21,2 & 20,1 & 18,9 & 18,5 \\ \hline
\end{tabularx}
\end{center}
\begin{enumerate}
    \item Calculez le taux de variation moyen annuel de la surface forestière sur la période 1990-2020 (en \% par an). Arrondissez à deux décimales.
    \item Calculez la surface moyenne perdue par an (en hectares) sur cette même période.
    \item En vous basant uniquement sur ces chiffres, la science (ici la mesure et la modélisation) permet-elle de \emph{décider} de la politique forestière à mener ? Justifiez en une phrase la distinction entre \cle{fait scientifique} et \cle{valeur/choix politique}.
\end{enumerate}
\end{exercice}

\courssub{Je m'entraîne}

\begin{exercice}[Application : Classer les disciplines]
Rangez les dix disciplines suivantes dans le tableau à trois colonnes (Sciences formelles, Sciences de la nature / expérimentales, Sciences humaines / sociales). Justifiez le classement de trois d'entre elles par une phrase chacune.
\begin{center}
Physique -- Mathématiques -- Histoire -- Chimie -- Sociologie -- Informatique théorique -- Biologie -- Économie -- Linguistique -- Logique
\end{center}
\end{exercice}

\begin{exercice}[Application : La méthode expérimentale en situation concrète]
Un agronome camerounais observe que dans une plantation de cacaoyers, certains plants résistent à la pourriture brune (causée par \emph{Phytophthora megakarya}) alors que d'autres meurent. Il veut comprendre pourquoi.
\begin{enumerate}
    \item Formalisez l'\cle{observation} initiale et la \cle{problème} scientifique posé.
    \item Proposez deux \cle{hypothèses} explicatives distinctes (l'une génétique, l'autre environnementale).
    \item Décrivez le \cle{protocole expérimental} (groupes témoins, variables contrôlées, variable indépendante, variable dépendante) pour tester l'hypothèse génétique.
    \item Quelle \cle{conclusion} provisoire peut-on tirer si les plants du groupe « variété résistante » survivent à 90\% contre 10\% pour le groupe « variété sensible » dans les mêmes conditions ?
\end{enumerate}
\end{exercice}

\begin{exercice}[Application : Obstacles épistémologiques chez Bachelard]
Pour chacun des énoncés ci-dessous, identifiez l'obstacle épistémologique principal (selon la typologie de Bachelard : obstacle de l'expérience première, obstacle du généralisable, obstacle réaliste, obstacle verbal, obstacle animiste, obstacle substantialiste, etc.) et expliquez en une phrase pourquoi c'est un obstacle pour la science.
\begin{enumerate}
    \item « L'eau bout à 100°C, c'est une vérité absolue. »
    \item « La chaleur est un fluide invisible (le calorique) qui passe du corps chaud au corps froid. »
    \item « Le soleil se lève à l'est et se couche à l'ouest, donc il tourne autour de la Terre. »
    \item « L'atome est indivisible, c'est dans son nom : a-tome. »
    \item « Cette plante guérit parce qu'elle a une "vertu" cachée donnée par la nature. »
\end{enumerate}
\end{exercice}

\begin{exercice}[Application : Science, bioéthique et écologie -- Cas concret]
Au Cameroun, un projet industriel propose l'épandage aérien d'un nouveau pesticide très efficace contre les ravageurs du cotonnier, mais des études toxicologiques indépendantes montrent un risque avéré pour les pollinisateurs (abeilles) et une persistance dans les nappes phréatiques.
\begin{enumerate}
    \item Identifiez le \cle{conflit de valeurs} entre l'intérêt économique court terme et l'intérêt écologique/sanitaire long terme.
    \item En quoi la \cle{science} (écotoxicologie, agronomie) est-elle \emph{nécessaire} pour éclairer le débat ? Donnez deux arguments.
    \item En quoi la science est-elle \emph{insuffisante} pour trancher seule la décision d'autoriser ou d'interdire ce pesticide ? Donnez deux arguments (rôle de l'éthique, du politique, du principe de précaution).
    \item Formulez une \cle{règle bioéthique} (inspirée du principe de responsabilité de Jonas ou du principe de précaution) qui pourrait guider la réglementation.
\end{enumerate}
\end{exercice}

\courssub{Je me prépare au Bac}

\begin{exercice}[Sujet de dissertation type Baccalauréat : « La science peut-elle tout expliquer ? »]
\textbf{Consigne :} Vous traiterez ce sujet en suivant une démarche dialectique rigoureuse (Thèse / Antithèse / Synthèse). Votre devoir comportera une introduction (accroche, définition des termes, problématique, annonce du plan), un développement en trois parties équilibrées, et une conclusion (bilan, ouverture).
\begin{enumerate}
    \item \textbf{Thèse :} Montrez que la science, par sa méthode (expérimentale, mathématique) et son extension continue, tend à expliquer l'ensemble du réel (matière, vie, société). Illustrer par des exemples (modélisation climatique, génomique, neurosciences).
    \item \textbf{Antithèse :} Montrez les limites intrinsèques de l'explication scientifique : distinction fait/valeur (Hume), indétermination quantique, complexité irréductible (systèmes chaotiques), subjectivité/conscience (qualia), questions de sens (métaphysique, éthique, esthétique) hors de portée de la méthode.
    \item \textbf{Synthèse :} Dépassez l'opposition. La science explique le « comment » (mécanismes, lois), mais ne peut pas trancher le « pourquoi » final (sens, valeurs, finalités). La rationalité scientifique est nécessaire mais non suffisante ; elle s'articule avec d'autres formes de rationalité (philosophique, éthique, politique, artistique) pour une compréhension plurielle du monde.
\end{enumerate}
\end{exercice}

\begin{exercice}[Explication de texte : Bachelard, \emph{La Formation de l'esprit scientifique} (1938)]
\textbf{Texte :} « La connaissance scientifique n'est pas une connaissance naturelle. Elle ne va pas de soi. Elle ne se donne pas par une intuition immédiate. Elle est une conquête. [...] Le premier obstacle que rencontre l'esprit scientifique, c'est l'expérience première. Elle est immédiate, substantielle, qualitative. La science, au contraire, est médiate, relationnelle, quantitative. [...] Il faut rompre avec les préjugés de l'expérience première, il faut faire une rupture épistémologique. »
\emph{(Gaston Bachelard, La Formation de l'esprit scientifique, Vrin, 1938, p. 17-18, adapté.)}

\textbf{Questions :}
\begin{enumerate}
    \item \textbf{Analyse :} Dégagez la thèse principale du texte et les arguments utilisés par Bachelard pour la soutenir.
    \item \textbf{Expliquez :} Les oppositions suivantes présentes dans le texte : « immédiate / médiate », « substantielle / relationnelle », « qualitative / quantitative ».
    \item \textbf{Concept :} Définissez la notion de « \cle{rupture épistémologique} » à partir du texte et de vos connaissances du cours.
    \item \textbf{Discussion :} L'expérience première (le sens commun) est-elle \emph{uniquement} un obstacle, ou peut-elle aussi être un point de départ heuristique pour la science ? Argumentez en vous appuyant sur un exemple concret (ex. : observation traditionnelle des plantes médicinales au Cameroun et recherche pharmacologique).
\end{enumerate}
\end{exercice}

\begin{exercice}[Question de synthèse : Science et décision écologique -- Déforestation au Cameroun]
\textbf{Document :} Extrait d'un rapport simulé (données inspirées du Global Forest Watch / MINEPDED).
\begin{center}
\begin{tabularx}{\linewidth}{|c|X|X|X|}\hline
Période & Perte annuelle moyenne (kha/an) & Cause principale identifiée (\% de la perte) & Émissions CO\textsubscript{2} associées (Mt CO\textsubscript{2}e/an) \\ \hline
1990-2000 & 85 & Agriculture itinérante (60\%), Exploitation forestière (30\%) & 45 \\ \hline
2000-2010 & 110 & Agriculture itinérante (50\%), Agro-industrie palmier/hévéa (35\%) & 62 \\ \hline
2010-2020 & 145 & Agro-industrie (55\%), Agriculture itinérante (30\%), Infrastructures (15\%) & 85 \\ \hline
\end{tabularx}
\end{center}

\textbf{Travail demandé :}
\begin{enumerate}
    \item \textbf{Analyse scientifique des données :} À partir du tableau, caractérisez l'évolution de la déforestation (quantitative) et de ses causes (qualitative) sur 30 ans. Quel rôle joue la science (télédétection, écologie, économie) dans la production de ce diagnostic ?
    \item \textbf{Limites de la science :} Expliquez pourquoi ces données scientifiques, aussi précises soient-elles, ne suffisent pas à \emph{déterminer} la politique forestière à adopter (ex. : moratoire total, zonage, intensification agricole, paiement pour services environnementaux). Mobilisez la distinction \cle{fait / valeur} et la notion de \cle{choix politique}.
    \item \textbf{Articulation science / éthique / politique :} Proposez un cadre de décision (type « science pour la décision ») intégrant : l'expertise scientifique (diagnostic, scénarios), l'évaluation éthique (justice intergénérationnelle, droits des populations locales/Baka/Bagyeli), et la délibération démocratique (participation des parties prenantes). Concluez sur la place de la science : « lampe » qui éclaire, pas « boussole » qui décide.
\end{enumerate}
\end{exercice}

\newpage

\coursec{Corrigés des exercices}

\coursec{Leçon 10 : La science — Exercices et corrigés types}

\begin{corrige}[Exercice 1 — QCM : Nature et types de sciences]
\begin{enumerate}
    \item \textbf{Faux.} La science ne se définit pas uniquement par son objet (la nature), mais d'abord par sa \cle{méthode} : démarche rigoureuse, vérifiable, fondée sur l'observation, l'expérimentation et la démonstration. Une même nature peut être étudiée non scientifiquement (mythe, opinion).
    \item \textbf{Vrai.} Les mathématiques sont une \cle{science formelle} : leurs objets (nombres, figures, structures) sont idéaux, et leurs vérités sont établies par \cle{déduction axiomatique} (démonstration à partir d'axiomes), indépendamment de l'expérience sensible.
    \item \textbf{Faux.} La sociologie est une \cle{science humaine} : elle utilise l'observation, les enquêtes et les statistiques, mais elle ne peut pas expérimenter sur l'homme comme la physique expérimente sur la matière (pas de laboratoire, pas de répétition contrôlée des conditions). Elle est empirique, non expérimentale au sens strict.
    \item \textbf{Faux.} Selon \cle{Bachelard}, l'esprit scientifique naît \emph{contre} le bon sens et l'expérience première, par une \cle{rupture épistémologique} avec les opinions spontanées (ex. : le soleil « semble » tourner autour de la Terre).
    \item \textbf{Vrai.} La \cle{méthode expérimentale} comporte nécessairement les étapes : observation, formulation du problème, hypothèse, expérience de vérification, conclusion (validation ou réfutation de l'hypothèse).
\end{enumerate}
\textbf{Points de vigilance :} ne pas confondre « science empirique » et « science expérimentale » ; citer Bachelard pour la question 4.
\end{corrige}

\begin{corrige}[Exercice 2 — Vrai / Faux justifié : Obstacles épistémologiques et méthode]
\begin{enumerate}
    \item \textbf{Faux.} L'\cle{obstacle épistémologique} (Bachelard) n'est pas une simple erreur : c'est un mode de pensée ancien, cohérent et séduisant (expérience première, images, mots) qui s'oppose durablement à la connaissance scientifique et qu'il faut \emph{rompre}, non simplement corriger.
    \item \textbf{Vrai.} Pour Bachelard, l'\cle{expérience première} (immédiate, qualitative, substantielle) est le premier et principal obstacle : elle donne l'illusion de connaître alors qu'elle livre des opinions non vérifiées.
    \item \textbf{Vrai.} La \cle{rupture épistémologique} désigne précisément le passage de l'\cle{opinion} (sens commun, préjugés) au \cle{savoir scientifique} (concepts construits, vérifiés, mathématisés).
    \item \textbf{Faux.} Une théorie scientifique n'est jamais définitive : elle reste \cle{révisable} et peut être réfutée ou dépassée par une théorie plus générale (ex. : la mécanique newtonienne dépassée par la relativité). C'est le critère de \cle{falsifiabilité} (Popper) qui fait la scientificité ; Bachelard parle de \emph{rectification} continue.
    \item \textbf{Faux.} La \cle{bioéthique} ne vise pas à freiner la science mais à l'\emph{orienter} : elle fixe des limites et des conditions (dignité humaine, consentement, responsabilité) pour que le progrès scientifique serve l'homme sans le détruire.
\end{enumerate}
\textbf{Points de vigilance :} la justification doit mobiliser explicitement une notion du cours (rupture, falsifiabilité, responsabilité) ; une réponse sans justification ne vaut aucun point.
\end{corrige}

\begin{corrige}[Exercice 3 — Définitions et distinctions]
\begin{enumerate}
    \item \textbf{Science formelle :} science dont les objets sont des \cle{êtres idéaux} (construits par l'esprit) et dont les vérités sont démontrées par \cle{déduction} à partir d'axiomes, sans recours à l'expérience. \emph{Exemple :} en mathématiques, le théorème de Pythagore est vrai par démonstration, non par mesure.
    \item \textbf{Science de la nature / science expérimentale :} science qui étudie les phénomènes matériels et vivants par la \cle{méthode expérimentale} (observation, hypothèse, expérience contrôlée, lois). \emph{Exemple :} la biologie teste l'efficacité d'un vaccin sur des groupes de patients.
    \item \textbf{Science humaine / science de l'esprit :} science qui étudie l'homme, ses conduites et ses productions sociales, par l'observation, l'enquête, les statistiques et l'interprétation, l'expérimentation y étant limitée. \emph{Exemple :} la sociologie étudie l'urbanisation de Douala par enquêtes et recensements.
    \item \textbf{Obstacle épistémologique (Bachelard) :} ensemble des habitudes de pensée (expérience première, images, mots, substantialisme) qui font obstacle à la formation de la connaissance scientifique et dont celle-ci doit se libérer par rupture. \emph{Exemple :} croire que la chaleur est un « fluide » (le calorique).
    \item \textbf{Bioéthique :} réflexion éthique sur les problèmes moraux posés par les applications de la science à la matière vivante et à la vie humaine (médecine, génétique, procréation). \emph{Exemple :} le débat sur la modification du génome humain ou sur les essais cliniques en Afrique.
\end{enumerate}
\textbf{Points de vigilance :} chaque définition doit comporter un critère distinctif (objet + méthode) et un exemple ; éviter les définitions circulaires (« la science formelle est une science qui est formelle »).
\end{corrige}

\begin{corrige}[Exercice 4 — Calcul direct et interprétation : Données écologiques au Cameroun]
\begin{enumerate}
    \item \textbf{Taux de variation moyen annuel (1990-2020).} Variation totale sur 30 ans :
    \[ \frac{\num{18,9} - \num{22,5}}{\num{22,5}} = \frac{-\num{3,6}}{\num{22,5}} = -\num{0,16} \quad \text{soit} \quad -\num{16}\,\% \text{ sur 30 ans.} \]
    Taux moyen annuel (approximation linéaire) : $-\num{16}\,\% / 30 \approx -\num{0,53}\,\%$ par an.
    
    \emph{Vérification par le coefficient multiplicateur annuel moyen (taux d'évolution moyen annuel géométrique) :} $(\num{18,9}/\num{22,5})^{1/30} = (\num{0,84})^{1/30} \approx \num{0,9942}$, soit environ $-\num{0,58}\,\%$ par an. Les deux ordres de grandeur sont cohérents ; on retient $\approx -\num{0,53}\,\%$ par an (méthode linéaire attendue en l'absence de précision).
    \item \textbf{Surface moyenne perdue par an :}
    \[ \frac{\num{22,5} - \num{18,9}}{30} = \frac{\num{3,6}}{30} = \num{0,12} \text{ Mha/an} = \SI{120000}{hectares} \text{ perdus par an.} \]
    \item \textbf{Fait scientifique et choix politique :} Non. La science établit des \cle{faits} (mesure de la perte, modélisation des tendances), mais elle ne dit pas ce qu'il \emph{faut} faire : décider d'un moratoire, d'un zonage ou d'une intensification relève d'un \cle{choix de valeurs} (développement économique, justice, préservation), donc du politique et de l'éthique. C'est la distinction entre le \emph{constat} (ce qui est) et la \emph{norme} (ce qui doit être), distinction fait/valeur (Hume).
\end{enumerate}
\textbf{Points de vigilance :} convertir correctement Mha en hectares ($1$ Mha $= 10^{6}$ ha) ; arrondir à deux décimales ; pour la question 3, citer explicitement la distinction fait/valeur.
\end{corrige}

\begin{corrige}[Exercice 5 — Application : Classer les disciplines]
\begin{center}
\begin{tabularx}{\linewidth}{|l|X|X|}\hline
\rowcolor{popblueL} \textbf{Sciences formelles} & \textbf{Sciences de la nature / expérimentales} & \textbf{Sciences humaines / sociales} \\ \hline
Mathématiques & Physique & Histoire \\ \hline
Logique & Chimie & Sociologie \\ \hline
Informatique théorique & Biologie & Économie \\ \hline
 & & Linguistique \\ \hline
\end{tabularx}
\end{center}
\textbf{Justifications (trois disciplines) :}
\begin{itemize}
    \item \textbf{Mathématiques} : science formelle, car elle porte sur des objets idéaux (nombres, structures) et procède par démonstration déductive, sans expérience sensible.
    \item \textbf{Biologie} : science expérimentale de la nature, car elle étudie le vivant par la méthode expérimentale (hypothèses testées en laboratoire ou sur le terrain).
    \item \textbf{Économie} : science humaine, car elle étudie les comportements et les échanges des hommes en société par l'observation, les statistiques et les modèles, sans expérimentation contrôlée possible à grande échelle.
\end{itemize}
\textbf{Points de vigilance :} l'informatique \emph{théorique} (algorithmique, calculabilité) est formelle ; la linguistique est une science humaine bien qu'elle utilise des méthodes formelles — c'est son \emph{objet} (le langage humain) qui décide du classement.
\end{corrige}

\begin{corrige}[Exercice 6 — Application : La méthode expérimentale en situation concrète]
\begin{enumerate}
    \item \textbf{Observation :} dans une même plantation, certains cacaoyers résistent à la pourriture brune, d'autres en meurent. \textbf{Problème scientifique :} quelle est la cause de cette différence de résistance au parasite \emph{Phytophthora megakarya} ?
    \item \textbf{Hypothèses :}
    \begin{itemize}
        \item \emph{Hypothèse génétique :} les plants résistants appartiennent à une variété possédant des gènes de résistance au parasite.
        \item \emph{Hypothèse environnementale :} les plants résistants bénéficient de conditions locales favorables (meilleur drainage du sol, ombrage, éloignement des foyers d'infection).
    \end{itemize}
    \item \textbf{Protocole expérimental (test de l'hypothèse génétique) :}
    \begin{itemize}
        \item \emph{Variable indépendante :} la variété du cacaoyer (variété supposée résistante / variété supposée sensible).
        \item \emph{Variable dépendante :} le taux de survie (ou le degré d'infection) après inoculation.
        \item \emph{Variables contrôlées :} même sol, même arrosage, même ombrage, même dose d'inoculation du parasite, même âge des plants.
        \item \emph{Dispositif :} deux groupes de plants (groupe A : variété résistante ; groupe B témoin : variété sensible), placés dans des conditions strictement identiques et inoculés avec le parasite ; comptage des survivants après une durée fixée.
    \end{itemize}
    \item \textbf{Conclusion provisoire :} si \num{90}\,\% des plants de la variété « résistante » survivent contre \num{10}\,\% pour la variété « sensible » dans des conditions identiques, l'hypothèse génétique est \emph{confortée} : la résistance dépend bien (au moins en partie) de facteurs génétiques. La conclusion reste \cle{provisoire} : elle devra être confirmée par la répétition de l'expérience et ne prouve pas que le facteur environnemental n'a aucun rôle.
\end{enumerate}
\textbf{Points de vigilance :} exiger un \cle{groupe témoin} et l'unicité de la variable testée (toutes les autres conditions constantes) ; la conclusion doit rester prudente (validation provisoire, jamais certitude absolue).
\end{corrige}

\begin{corrige}[Exercice 7 — Application : Obstacles épistémologiques chez Bachelard]
\begin{enumerate}
    \item \textbf{Obstacle du généralisable (et de l'expérience première) :} on érige en vérité absolue une constatation valable seulement dans des conditions précises ; or l'eau bout à \SI{100}{\celsius} seulement à pression atmosphérique normale (elle bout vers \SI{90}{\celsius} en altitude, comme sur les monts Bamboutos). La science précise les \cle{conditions de validité} des lois.
    \item \textbf{Obstacle substantialiste :} on imagine la chaleur comme une \emph{substance} (un fluide, le calorique) au lieu de la concevoir comme un \emph{phénomène} (agitation moléculaire, transfert d'énergie). La science substitue des relations mesurables à des substances imaginaires.
    \item \textbf{Obstacle de l'expérience première :} l'apparence immédiate (le soleil « se lève ») est prise pour la réalité ; la science (Copernic, Galilée) montre que c'est la Terre qui tourne. L'évidence sensible est trompeuse.
    \item \textbf{Obstacle verbal :} on croit expliquer la réalité par le \emph{mot} (« a-tome » = indivisible) ; or la physique a montré que l'atome est divisible (noyau, électrons, particules). Le langage ne tient pas lieu de démonstration.
    \item \textbf{Obstacle animiste (et substantialiste) :} attribuer à la plante une « vertu » cachée, quasi intentionnelle, c'est substituer une explication magique à l'analyse des \cle{principes actifs} chimiques que la pharmacologie peut isoler, doser et tester.
\end{enumerate}
\textbf{Points de vigilance :} nommer précisément l'obstacle selon la typologie de Bachelard et dire en quoi il bloque la démarche scientifique (il arrête la recherche au lieu de la provoquer).
\end{corrige}

\begin{corrige}[Exercice 8 — Application : Science, bioéthique et écologie -- Cas concret]
\begin{enumerate}
    \item \textbf{Conflit de valeurs :} d'un côté, l'\cle{intérêt économique} immédiat (rendements du coton, revenus des producteurs, devises pour la SODECOTON et l'État) ; de l'autre, l'\cle{intérêt écologique et sanitaire} à long terme (survie des pollinisateurs, qualité de l'eau potable, santé des populations). C'est un conflit entre le court terme et le long terme, entre le profit privé et le bien commun.
    \item \textbf{Nécessité de la science :}
    \begin{itemize}
        \item L'\cle{écotoxicologie} fournit les faits : doses létales pour les abeilles, durée de persistance dans les nappes phréatiques, concentrations mesurées. Sans elle, le débat se fonderait sur des rumeurs ou des intérêts.
        \item L'\cle{agronomie} permet d'évaluer les alternatives (lutte biologique, variétés résistantes) et de comparer objectivement les coûts et les bénéfices de chaque scénario.
    \end{itemize}
    \item \textbf{Insuffisance de la science :}
    \begin{itemize}
        \item La science dit ce qui \emph{est} et ce qui \emph{risque} d'arriver, mais pas ce qu'il \emph{faut} faire : accepter un risque pour les générations futures est un \cle{choix de valeurs} (distinction fait/valeur), qui relève de l'éthique et du politique.
        \item Face à l'incertitude scientifique elle-même (risque « avéré » mais ampleur discutée), seul le \cle{principe de précaution} et la délibération démocratique peuvent fixer le niveau de risque jugé acceptable par la société.
    \end{itemize}
    \item \textbf{Règle bioéthique possible :} « Nul produit ne peut être autorisé tant que son innocuité pour les écosystèmes et pour les générations futures n'est pas établie ; en cas de doute sérieux, la charge de la preuve incombe à l'exploitant, et l'intérêt des générations futures prime sur le profit présent » (inspirée du \cle{principe de responsabilité} de Hans Jonas et du \cle{principe de précaution}).
\end{enumerate}
\textbf{Points de vigilance :} bien distinguer ce que la science apporte (diagnostic, scénarios) de ce qu'elle ne peut pas décider (normes, valeurs) ; citer Jonas ou le principe de précaution dans la règle formulée.
\end{corrige}

\begin{corrige}[Exercice 9 — Sujet de dissertation type Baccalauréat : « La science peut-elle tout expliquer ? »]
\textbf{Corrigé détaillé (plan et contenus attendus) :}

\textbf{Introduction.} \emph{Accroche :} la science moderne a transformé la condition humaine (médecine, techniques, communication). \emph{Définition :} la science est un savoir méthodique, vérifiable, fondé sur l'expérience et la démonstration ; « expliquer » signifie rendre compte d'un phénomène par des causes et des lois. \emph{Problématique :} la méthode scientifique, victorieuse dans le domaine de la matière, de la vie et même des sociétés, peut-elle s'étendre à la totalité du réel, ou rencontre-t-elle des limites intrinsèques ? \emph{Annonce du plan :} puissance de l'explication scientifique (I), ses limites (II), dépassement : une rationalité nécessaire mais non suffisante (III).

\textbf{I. Thèse — La science tend à expliquer l'ensemble du réel.}
\begin{itemize}
    \item La \cle{méthode expérimentale} et la mathématisation ont permis de formuler des lois universelles (physique, chimie) et de maîtriser la matière.
    \item L'extension continue du champ scientifique : la biologie et la \cle{génomique} expliquent l'hérédité et les maladies ; la modélisation climatique prévoit l'évolution des températures ; les sciences humaines (économie, sociologie) dégagent des régularités dans les conduites collectives.
    \item L'histoire des sciences montre que des domaines jugés « inexplicables » (foudre, maladies infectieuses) sont devenus explicables : rien ne semble en principe hors de portée.
\end{itemize}

\textbf{II. Antithèse — Les limites intrinsèques de l'explication scientifique.}
\begin{itemize}
    \item \emph{Limite logique :} la distinction \cle{fait/valeur} (Hume) — la science constate ce qui est, elle ne peut pas fonder ce qui doit être (le bien, le juste).
    \item \emph{Limites internes :} indétermination en physique quantique, imprévisibilité des systèmes complexes et chaotiques (météorologie à long terme), révisabilité permanente des théories.
    \item \emph{Limites de domaine :} la subjectivité vécue (la douleur, la conscience), les questions de \cle{sens} (pourquoi vivre ? qu'est-ce qu'une vie bonne ? la beauté) échappent à la méthode expérimentale : ce sont des questions métaphysiques, éthiques et esthétiques.
\end{itemize}

\textbf{III. Synthèse — Une rationalité nécessaire mais non suffisante.}
\begin{itemize}
    \item La science explique le \emph{comment} (mécanismes, lois, conditions) mais ne tranche pas le \emph{pourquoi} final (sens, finalités, valeurs).
    \item Elle demeure \cle{indispensable} : nulle décision éclairée (bioéthique, écologie) ne peut ignorer ses diagnostics.
    \item Mais elle doit s'articuler avec d'autres formes de rationalité — philosophique, éthique, politique, artistique — pour une compréhension plurielle du monde et pour orienter ses propres applications (Jonas : responsabilité).
\end{itemize}

\textbf{Conclusion.} \emph{Bilan :} la science peut expliquer beaucoup, et toujours davantage, mais non tout : elle laisse hors d'elle les questions de valeur et de sens. \emph{Ouverture :} faut-il alors souhaiter que la science soit encadrée par l'éthique, ou l'éthique elle-même doit-elle devenir scientifique ?

\textbf{Barème indicatif (sur 20) :} introduction (accroche, définitions, problématique, plan) : 4 pts ; thèse argumentée et illustrée : 4 pts ; antithèse argumentée et illustrée : 4 pts ; synthèse (dépassement réel, non simple compromis) : 4 pts ; cohérence, langue, exemples pertinents (dont exemples camerounais valorisés) : 4 pts.

\textbf{Points de vigilance :} définir les termes du sujet dès l'introduction ; éviter la juxtaposition (la synthèse doit \emph{dépasser} l'opposition) ; citer au moins deux auteurs (Hume, Bachelard, Jonas) ; chaque partie doit comporter au moins un exemple concret.
\end{corrige}

\begin{corrige}[Exercice 10 — Explication de texte : Bachelard, \emph{La Formation de l'esprit scientifique} (1938)]
\begin{enumerate}
    \item \textbf{Thèse et arguments.} \emph{Thèse :} la connaissance scientifique n'est pas un prolongement naturel de l'expérience courante, mais une \cle{conquête} contre elle. \emph{Arguments :} (a) la science ne se donne pas par intuition immédiate, elle se construit ; (b) l'expérience première est le premier obstacle, car elle est immédiate, substantielle et qualitative, alors que la science est médiate, relationnelle et quantitative ; (c) il faut donc opérer une \cle{rupture épistémologique} avec les préjugés du sens commun.
    \item \textbf{Explication des oppositions :}
    \begin{itemize}
        \item \emph{Immédiate / médiate :} l'expérience première saisit les choses directement par les sens ; la science passe par des \emph{intermédiaires} (instruments, mesures, théories) qui corrigent et dépassent le donné sensible.
        \item \emph{Substantielle / relationnelle :} le sens commun croit connaître des substances (le « feu », la « chaleur ») ; la science connaît des \emph{relations} entre phénomènes (lois, rapports mesurables).
        \item \emph{Qualitative / quantitative :} l'expérience première décrit des qualités (chaud, lourd, amer) ; la science \emph{mesure} et mathématise (température en degrés, masse en kilogrammes).
    \end{itemize}
    \item \textbf{Définition du concept.} La \cle{rupture épistémologique} est l'acte par lequel l'esprit scientifique rompt avec les évidences, les images et les opinions de l'expérience première pour construire des concepts nouveaux, définis, mesurables et vérifiables. Elle n'est pas un simple progrès continu, mais une discontinuité, une réforme de la pensée (ex. : rompre avec l'idée que le soleil tourne autour de la Terre).
    \item \textbf{Discussion.} L'expérience première n'est pas \emph{uniquement} un obstacle : elle peut fournir un \cle{point de départ heuristique}, c'est-à-dire une indication précieuse pour la recherche. \emph{Exemple :} la médecine traditionnelle camerounaise utilise depuis longtemps certaines plantes (écorces, feuilles) contre le paludisme ou les infections ; cette observation empirique a orienté la recherche \cle{pharmacologique}, qui a pu isoler, doser et tester les principes actifs (comme l'artémisinine fut isolée à partir d'une plante de la pharmacopée chinoise). Mais l'expérience première ne devient science qu'à condition d'être \emph{dépassée} : vérification expérimentale, dosage, essais cliniques. Elle est donc un tremplin possible, jamais une preuve : c'est la rupture méthodique qui la transforme en savoir.
\end{enumerate}
\textbf{Barème indicatif (sur 20) :} dégagement de la thèse : 4 pts ; analyse des arguments : 4 pts ; explication des trois oppositions : 4 pts ; définition de la rupture épistémologique : 3 pts ; discussion argumentée avec exemple : 4 pts ; présentation et langue : 1 pt.

\textbf{Points de vigilance :} ne pas paraphraser le texte mais l'\emph{expliquer} (dire ce que chaque formule signifie) ; dans la discussion, éviter de contredire purement Bachelard : montrer que l'expérience première peut être un point de départ \emph{à condition d'être rompue et vérifiée} ; l'exemple camerounais (pharmacopée traditionnelle) est fortement valorisé.
\end{corrige}

\begin{corrige}[Exercice 11 — Question de synthèse : Science et décision écologique -- Déforestation au Cameroun]
\begin{enumerate}
    \item \textbf{Analyse scientifique des données.}
    \begin{itemize}
        \item \emph{Évolution quantitative :} la perte annuelle moyenne passe de \num{85} kha/an (1990-2000) à \num{110} kha/an (2000-2010) puis \num{145} kha/an (2010-2020), soit une hausse d'environ \num{71}\,\% en 30 ans ($145/85 \approx 1,71$) ; les émissions de CO\textsubscript{2} suivent la même tendance (\num{45} à \num{85} Mt CO\textsubscript{2}e/an, soit $+\num{89}\,\%$). La déforestation s'\emph{accélère}.
        \item \emph{Évolution qualitative :} la cause dominante se déplace : l'agriculture itinérante (\num{60}\,\% en 1990-2000) recule (\num{30}\,\% en 2010-2020) tandis que l'\cle{agro-industrie} (palmier à huile, hévéa) devient le premier facteur (\num{55}\,\%). La déforestation passe d'une logique paysanne de subsistance à une logique industrielle et marchande.
        \item \emph{Rôle de la science :} la \cle{télédétection} (images satellites) mesure les surfaces perdues ; l'\cle{écologie} identifie les mécanismes et les conséquences (biodiversité, carbone) ; l'\cle{économie} et la statistique attribuent les causes. Sans ces disciplines, aucun diagnostic objectif ne serait possible.
    \end{itemize}
    \item \textbf{Limites de la science.} Ces données, même exactes, ne dictent aucune politique. Choisir entre un moratoire total, un zonage, l'intensification agricole ou le paiement pour services environnementaux suppose d'\emph{arbitrer entre des valeurs} : développement économique et emplois \emph{contre} préservation du climat et de la biodiversité ; droits des populations actuelles \emph{contre} devoirs envers les générations futures. Or, selon la distinction \cle{fait/valeur}, aucun constat scientifique ne peut fonder à lui seul une norme : « ce qui est » ne dit pas « ce qui doit être ». La décision relève d'un \cle{choix politique}, c'est-à-dire d'une délibération collective sur le type de société souhaité.
    \item \textbf{Articulation science / éthique / politique.} Un cadre de décision de type « science pour la décision » comprend trois moments :
    \begin{itemize}
        \item \emph{Expertise scientifique :} établir le diagnostic (données du tableau), construire des \cle{scénarios} chiffrés (que se passe-t-il si l'on intensifie l'agriculture ? si l'on crée des corridors forestiers ?) et évaluer les risques de chaque option.
        \item \emph{Évaluation éthique :} interroger chaque scénario au regard de la \cle{justice intergénérationnelle} (Jonas : ne pas hypothéquer l'avenir) et des \cle{droits des populations locales} (Baka, Bagyeli), dont les modes de vie dépendent de la forêt.
        \item \emph{Délibération démocratique :} associer toutes les parties prenantes (État, agro-industriels, communautés locales et autochtones, scientifiques, ONG) à une décision transparente et révisable.
    \end{itemize}
    \emph{Conclusion :} la science est une \cle{lampe} qui éclaire les conséquences de nos choix, non une \cle{boussole} qui décide à notre place : elle rend la décision possible et responsable, mais la décision demeure un acte éthique et politique.
\end{enumerate}
\textbf{Barème indicatif (sur 20) :} analyse quantitative correcte (calculs de variation) : 4 pts ; analyse qualitative des causes : 3 pts ; rôle de la science dans le diagnostic : 3 pts ; distinction fait/valeur et limites de la science : 4 pts ; cadre de décision articulé (science/éthique/politique) : 4 pts ; conclusion et qualité de l'expression : 2 pts.

\textbf{Points de vigilance :} vérifier les calculs de pourcentages ($+\num{71}\,\%$ de perte, $+\num{89}\,\%$ d'émissions) ; ne pas écrire que la science « ne sert à rien » : elle est \emph{nécessaire mais non suffisante} ; citer explicitement la distinction fait/valeur et nommer les populations concernées (Baka, Bagyeli) dans le volet éthique.
\end{corrige}

\newpage

\coursec{Fiche bilan}

\begin{popfigure}
\begin{center}\tcbox[colback=popink,colframe=popink,coltext=white,arc=9pt,boxsep=4pt,left=6pt,right=6pt]{\bfseries\small LA SCIENCE}\end{center}
\vspace{-2pt}
\begin{tcbraster}[raster columns=3,raster equal height=rows,raster column skip=6pt,raster row skip=6pt,enhanced,blankest]
\begin{tcolorbox}[enhanced,colback=white,colframe=popblue,boxrule=0.9pt,arc=3pt,title={Nature et caract\`eres},fonttitle=\bfseries\scriptsize,coltitle=white,colbacktitle=popblue,left=4pt,right=4pt,top=2pt,bottom=2pt,boxsep=1pt,toptitle=1.5pt,bottomtitle=1.5pt,halign=left]
\begin{itemize}[nosep,leftmargin=*,itemsep=1pt,font=\scriptsize]
\item Savoir objectif, v\'erifiable
\item Conna\^itre par les causes
\end{itemize}
\end{tcolorbox}
\begin{tcolorbox}[enhanced,colback=white,colframe=poporange,boxrule=0.9pt,arc=3pt,title={Types de sciences},fonttitle=\bfseries\scriptsize,coltitle=white,colbacktitle=poporange,left=4pt,right=4pt,top=2pt,bottom=2pt,boxsep=1pt,toptitle=1.5pt,bottomtitle=1.5pt,halign=left]
\begin{itemize}[nosep,leftmargin=*,itemsep=1pt,font=\scriptsize]
\item Math\'ematiques
\item Sciences de la mati\`ere
\item Sciences humaines
\end{itemize}
\end{tcolorbox}
\begin{tcolorbox}[enhanced,colback=white,colframe=popgreen,boxrule=0.9pt,arc=3pt,title={Esprit scientifique},fonttitle=\bfseries\scriptsize,coltitle=white,colbacktitle=popgreen,left=4pt,right=4pt,top=2pt,bottom=2pt,boxsep=1pt,toptitle=1.5pt,bottomtitle=1.5pt,halign=left]
\begin{itemize}[nosep,leftmargin=*,itemsep=1pt,font=\scriptsize]
\item Conditions d'\'emergence
\item Obstacles \'epist\'emologiques (Bachelard)
\end{itemize}
\end{tcolorbox}
\begin{tcolorbox}[enhanced,colback=white,colframe=poppurple,boxrule=0.9pt,arc=3pt,title={M\'ethode exp\'erimentale},fonttitle=\bfseries\scriptsize,coltitle=white,colbacktitle=poppurple,left=4pt,right=4pt,top=2pt,bottom=2pt,boxsep=1pt,toptitle=1.5pt,bottomtitle=1.5pt,halign=left]
\begin{itemize}[nosep,leftmargin=*,itemsep=1pt,font=\scriptsize]
\item Observation, hypoth\`ese
\item Exp\'erience, loi
\end{itemize}
\end{tcolorbox}
\begin{tcolorbox}[enhanced,colback=white,colframe=poppink,boxrule=0.9pt,arc=3pt,title={Valeur de la science},fonttitle=\bfseries\scriptsize,coltitle=white,colbacktitle=poppink,left=4pt,right=4pt,top=2pt,bottom=2pt,boxsep=1pt,toptitle=1.5pt,bottomtitle=1.5pt,halign=left]
\begin{itemize}[nosep,leftmargin=*,itemsep=1pt,font=\scriptsize]
\item Utilit\'e et progr\`es
\item Limites et probl\`emes
\end{itemize}
\end{tcolorbox}
\begin{tcolorbox}[enhanced,colback=white,colframe=popgold,boxrule=0.9pt,arc=3pt,title={Bio\'ethique et \'ecologie},fonttitle=\bfseries\scriptsize,coltitle=white,colbacktitle=popgold,left=4pt,right=4pt,top=2pt,bottom=2pt,boxsep=1pt,toptitle=1.5pt,bottomtitle=1.5pt,halign=left]
\begin{itemize}[nosep,leftmargin=*,itemsep=1pt,font=\scriptsize]
\item Respect du vivant
\item Responsabilit\'e (Jonas)
\end{itemize}
\end{tcolorbox}
\end{tcbraster}

\legende{Carte mentale de la le\c{c}on : les six branches essentielles de \og La science \fg{} au programme de Terminale technique.}
\end{popfigure}

\begin{bilan}
\textbf{1. D\'efinitions cl\'es}
\begin{itemize}
\item \cle{Science} : ensemble de connaissances \emph{objectives}, m\'ethodiques et v\'erifiables, qui expliquent le r\'eel par ses \emph{causes} et ses \emph{lois}, et non par les opinions ou les croyances.
\item \cle{Objectivit\'e} : caract\`ere d'un savoir ind\'ependant des go\^uts, des passions et des pr\'ejug\'es du sujet ; il vaut pour tous.
\item \cle{Obstacle \'epist\'emologique} (Bachelard) : erreur ou habitude de pens\'ee (opinion, premi\`ere exp\'erience, images, anthropocentrisme) qui emp\^eche la pens\'ee scientifique de se former.
\item \cle{M\'ethode exp\'erimentale} (Claude Bernard) : d\'emarche rigoureuse allant de l'observation \`a la loi, en passant par l'hypoth\`ese et l'exp\'erience.
\item \cle{Bio\'ethique} : r\'eflexion sur les conditions d'application de la science \`a la mati\`ere vivante (homme, animal, environnement) afin de prot\'eger la vie et la dignit\'e humaine.
\end{itemize}

\textbf{2. Les types de sciences \`a conna\^itre par c\oe ur}
\begin{center}
\begin{tabularx}{\linewidth}{|l|X|X|}\hline
\rowcolor{popblueL} \textbf{Type} & \textbf{Objet} & \textbf{Caract\`ere} \\ \hline
Math\'ematiques & Nombres, figures, relations & Sciences \emph{formelles} : v\'erit\'es d\'emontr\'ees, certaines \\ \hline
Sciences de la mati\`ere & Physique, chimie, biologie & Sciences \emph{exp\'erimentales} : v\'erit\'es v\'erifi\'ees par l'exp\'erience \\ \hline
Sciences humaines & Homme et soci\'et\'e (histoire, sociologie\ldots) & Objectivit\'e plus difficile : le chercheur fait partie de son objet \\ \hline
\end{tabularx}
\end{center}

\textbf{3. Conditions d'\'emergence de l'esprit scientifique}
\begin{itemize}
\item Le \emph{doute m\'ethodique} : se m\'efier des \'evidences et des opinions re\c{c}ues.
\item Le \emph{refus de l'autorit\'e} : n'accepter que ce qui est prouv\'e (Galil\'ee contre Aristote).
\item La \emph{rigueur logique} et l'\emph{observation exacte} (instruments, mesures).
\item La \emph{rupture} avec le sens commun : \og il faut d\'etruire l'opinion \fg{} (Bachelard).
\end{itemize}

\textbf{4. Les \'etapes de la m\'ethode exp\'erimentale (Claude Bernard)}
\begin{enumerate}
\item \textbf{Observation} exacte du ph\'enom\`ene ;
\item \textbf{Hypoth\`ese} : explication provisoire \`a tester ;
\item \textbf{Exp\'erience} : v\'erification contr\^ol\'ee de l'hypoth\`ese ;
\item \textbf{Conclusion} : confirmation ou rejet de l'hypoth\`ese, \'etablissement d'une \textbf{loi}.
\end{enumerate}

\textbf{5. Valeur et limites de la science}
\begin{itemize}
\item \emph{Valeur th\'eorique} : elle fait conna\^itre la v\'erit\'e et lib\`ere l'esprit de l'erreur et de la superstition.
\item \emph{Valeur pratique} : elle rend la technique possible (m\'edecine, agriculture, vaccins au Cameroun).
\item \emph{Limites} : la science dit \emph{comment} agir, non \emph{pourquoi} il faut agir ; elle ne r\'epond pas aux questions de sens (Bergson : la science ne dit pas tout du r\'eel).
\item \emph{Probl\`emes} : usages dangereux (armes, manipulations g\'en\'etiques, pollution) $\Rightarrow$ n\'ecessit\'e de la \cle{bio\'ethique} et du \emph{principe de responsabilit\'e} (Hans Jonas) envers les g\'en\'erations futures et la nature.
\end{itemize}

\textbf{6. M\'ethode express pour la dissertation}
\begin{enumerate}
\item D\'efinir la science d\`es l'introduction (opposer science / opinion) ;
\item Probl\'ematiser : la science est-elle toujours un bien ? Suffit-elle \`a tout expliquer ?
\item Mobiliser au moins deux auteurs : Bachelard (obstacles), Claude Bernard (m\'ethode), Jonas (responsabilit\'e) ;
\item Illustrer par un exemple concret camerounais (vaccination, OGM, d\'eforestation) ;
\item Conclure en nuancant : la science est utile mais doit \^etre \'eclair\'ee par l'\'ethique.
\end{enumerate}
\end{bilan}

\begin{aretenir}[Checklist avant le Bac]
\begin{itemize}
\item[$\square$] Je sais d\'efinir la science et la distinguer de l'opinion et de la croyance.
\item[$\square$] Je connais les caract\`eres de la science : objectivit\'e, m\'ethode, v\'erifiabilit\'e, explication par les causes.
\item[$\square$] Je sais citer les trois types de sciences (math\'ematiques, sciences de la mati\`ere, sciences humaines) avec un exemple pour chacun.
\item[$\square$] Je connais les conditions d'\'emergence de l'esprit scientifique (doute, observation, rupture avec le sens commun).
\item[$\square$] Je sais d\'efinir l'obstacle \'epist\'emologique et citer Bachelard.
\item[$\square$] Je ma\^itrise les quatre \'etapes de la m\'ethode exp\'erimentale de Claude Bernard.
\item[$\square$] Je sais expliquer pourquoi l'application de la science au vivant pose un probl\`eme moral (bio\'ethique).
\item[$\square$] Je connais le principe de responsabilit\'e de Hans Jonas face \`a la question \'ecologique.
\item[$\square$] Je sais discuter la valeur de la science : utilit\'e, progr\`es, mais aussi limites et dangers.
\item[$\square$] Je dispose d'exemples camerounais concrets (vaccination, agriculture, environnement).
\item[$\square$] Je sais construire une probl\'ematique et r\'ediger une conclusion nuanc\'ee et justifi\'ee.
\end{itemize}
\end{aretenir}

\findecours

\end{document}
